% Leçon 2 — Grammaire : 第一……第二…. ; 不但……而且…. — Chinois 1ère A — Cours signé OnBuch+
% Programme officiel MINESEC. Compiler avec : tectonic ZHA02-grammaire.tex
\newcommand{\DOCMATIERE}{Chinois}
\newcommand{\DOCNIVEAU}{1ère A}
\newcommand{\DOCMODULE}{Module 1 : Vie familiale et intégration sociale}
\newcommand{\DOCLECON}{ZHA02}
\newcommand{\DOCTITRE}{Leçon 2 — Grammaire : 第一……第二…. ; 不但……而且….}
% OnBuch+ — préambule « cours signés » ENRICHI (compilé avec tectonic / XeLaTeX)
% Système visuel « Pop » d'OnBuch+ : fond crème (jamais blanc), bordures encre
% épaisses, ombres « dures » décalées, titres Archivo Black en capitales.
% Version MATHÉMATIQUES : graphes (pgfplots/tikz), boîtes pédagogiques
% (définition, propriété, méthode, attention, le savais-tu, exercice résolu,
% exercice, TP, bilan), page de garde et signature OnBuch+ ancrée en bas.
%
% Macros attendues dans main.tex AVANT \input{preamble} :
%   \DOCMATIERE  \DOCNIVEAU  \DOCMODULE  \DOCLECON (numéro)  \DOCTITRE
\documentclass[11pt]{article}

\usepackage{fontspec}
\usepackage[french]{babel} % français = langue principale ; le chinois via \zh{..} / textebox (xeCJK)
\usepackage{xeCJK}
\usepackage{pifont} % \ding{51} \ding{55} utilisés par les modèles
\usepackage[a4paper,margin=2cm,top=3.4cm,bottom=2.8cm,headheight=1.4cm,footskip=1.3cm]{geometry}
\usepackage{amsmath,amssymb,amsfonts}
\usepackage{mathtools} % \xmapsto, \coloneqq... souvent utilisés par les modèles
\usepackage{cancel} % simplifications barrées (bilans de forces)
% \abs{z} (module, valeur absolue) et \norm{u} (norme), utilisés par les modèles
\providecommand{\abs}{}\let\abs\relax\DeclarePairedDelimiter{\abs}{\lvert}{\rvert}
\providecommand{\norm}{}\let\norm\relax\DeclarePairedDelimiter{\norm}{\lVert}{\rVert}
\usepackage[table]{xcolor} % [table] : fournit aussi \rowcolors (alternance de lignes)
\usepackage{pagecolor}
\usepackage{tikz}
\usetikzlibrary{arrows.meta,positioning,calc,shapes.geometric,shapes.misc,shapes.arrows,decorations.pathmorphing,decorations.pathreplacing,decorations.markings,patterns,shadows,mindmap,trees,fit,backgrounds,matrix,intersections,angles,quotes}
\usepackage{pgfplots}
\usepackage[version=4]{mhchem} % équations nucléaires \ce{^{235}_{92}U}
\usepackage[european,siunitx]{circuitikz} % schémas électriques
\pgfplotsset{compat=1.18}
\usepgfplotslibrary{fillbetween,groupplots} % aires entre courbes (calcul intégral)
\usepackage{enumitem}
\usepackage{fancyhdr}
% [most] charge toutes les bibliothèques tcolorbox (listings, minted, raster,
% documentation...) et gonfle inutilement la mémoire TeX sur un document long
% (~300 boîtes) : on ne charge que ce qui est réellement utilisé.
\usepackage[skins,breakable]{tcolorbox}
\usepackage{graphicx}
\usepackage{colortbl}
\usepackage{booktabs}
\usepackage{tabularx}
\usepackage{diagbox} % case d'en-tête diagonale des tableaux à double entrée
% Colonnes extensibles centrée / gauche / droite (C, L, R), souvent utilisées par les modèles
\newcolumntype{C}{>{\centering\arraybackslash}X}
\newcolumntype{L}{>{\raggedright\arraybackslash}X}
\newcolumntype{R}{>{\raggedleft\arraybackslash}X}
\usepackage{array}
\usepackage{multicol}
\usepackage{multirow}
\usepackage{siunitx}
% parse-numbers=false : accepte aussi \SI{2,5 \times 10^{-3}}{...}
\sisetup{locale=FR,output-decimal-marker={,},group-minimum-digits=4,parse-numbers=false}
\DeclareSIUnit\ohm{\ensuremath{\symup{\Omega}}} % U+2126 absent de la police
\DeclareSIUnit\division{div} % réglages d'oscilloscope (ms/div, V/div)
\DeclareSIUnit\tr{tr} % tours (vitesse de rotation, ex. \SI{1500}{\tr\per\minute})
\DeclareSIUnit\molar{M} % concentration molaire (ex. \SI{3}{\molar}, \SI{1}{\micro\molar})
\DeclareSIUnit\an{an} % année (le modèle confond parfois avec \year, un compteur TeX) — ex. \SI{5}{\kilo\meter\per\an}
\DeclareSIUnit\FCFA{FCFA} % franc CFA (ex. \SI{500}{\FCFA\per\kilo\gram})
\DeclareSIUnit\degreeBrix{\textdegree Brix} % teneur en sucre d'un sirop/jus (réfractométrie)
\DeclareSIUnit\month{mois} % le modèle utilise parfois \month, un compteur TeX, comme unité de durée
\usepackage{qrcode}
% \degree : commande fréquemment utilisée par les modèles pour un angle ou une
% température en dehors de \SI{}{} (non fournie par les paquets chargés).
\providecommand{\degree}{^{\circ}}
% \qed (fin de démonstration), utilisé par les modèles sans amsthm
\providecommand{\qed}{\hfill$\square$}
% \barre : double barre des valeurs interdites dans un tableau de variations (tkz-tab)
\providecommand{\barre}{\ensuremath{\|}}
% environnement proof (amsthm non chargé) utilisé par les modèles
\ifdefined\proof\else\newenvironment{proof}[1][Démonstration]{\par\noindent\textit{#1.}\ }{\qed\par}\fi
\usepackage{lastpage}
\usepackage{hyperref}
\hypersetup{hidelinks}

% ============================================================
% Palette « Pop » OnBuch+ — mêmes hex que la classe P (plus_ui.dart)
% ============================================================
\definecolor{poporange}{HTML}{F59321}
\definecolor{popdark}{HTML}{E07A0C}
\definecolor{popcream}{HTML}{FFF6E8}
\definecolor{popink}{HTML}{1C1714}
\definecolor{poppurple}{HTML}{7A5AE0}
\definecolor{popgreen}{HTML}{2E9E6A}
\definecolor{poppink}{HTML}{E0457B}
\definecolor{popblue}{HTML}{2F7DE1}
\definecolor{popgold}{HTML}{C98A12}
\definecolor{popmuted}{HTML}{8A7F76}
\definecolor{popline}{HTML}{EADFD2}
\definecolor{popgray}{HTML}{8A7F76}
\colorlet{popgrey}{popgray}
% Alias tolérants (le modèle nomme parfois les couleurs différemment)
\colorlet{popwhite}{white}
\colorlet{popblack}{popink}
\colorlet{popred}{poppink}
\colorlet{popyellow}{popgold}
% Teintes claires (fonds de tableaux/encadrés) — le modèle nomme parfois
% les couleurs avec un suffixe L (ex. popblueL) sans les avoir définies.
\colorlet{poporangeL}{poporange!20}
\colorlet{popdarkL}{popdark!20}
\colorlet{poppurpleL}{poppurple!20}
\colorlet{popgreenL}{popgreen!20}
\colorlet{poppinkL}{poppink!20}
\colorlet{popblueL}{popblue!20}
\colorlet{popgoldL}{popgold!20}
\colorlet{popredL}{popred!20}
\colorlet{popgrayL}{popgray!20}
% Teintes claires pour fonds de tableaux / zones
\colorlet{popblueL}{popblue!10!white}
\colorlet{poporangeL}{poporange!14!white}
\colorlet{popgreenL}{popgreen!12!white}
\colorlet{poppurpleL}{poppurple!12!white}
\colorlet{poppinkL}{poppink!12!white}
\colorlet{popgoldL}{popgold!14!white}

\pagecolor{popcream}

% ============================================================
% Typographie
% ============================================================
\defaultfontfeatures{Ligatures=TeX}
\setmainfont{PlusJakartaSans-Medium.ttf}[Path=../../../fonts/,BoldFont=PlusJakartaSans-Bold.ttf]
% Symboles, API et emojis absents de la police principale : repli sur DejaVu Sans (caractères actifs)
\newfontface\symfont{DejaVuSans.ttf}[Path=../../../fonts/]
\makeatletter
\newcommand\symactive[1]{\catcode#1=13 \begingroup\lccode`\~=#1 \lowercase{\endgroup\def~}{{\symfont\char#1}}}
\newcommand\symblank[1]{\catcode#1=13 \begingroup\lccode`\~=#1 \lowercase{\endgroup\def~}{}}
\newcommand\symhyph[1]{\catcode#1=13 \begingroup\lccode`\~=#1 \lowercase{\endgroup\def~}{\mbox{-}}}
\newcount\symc
\newcommand\symrange[2]{\symc=#1 \@whilenum\symc<#2 \do{\expandafter\symactive\expandafter{\the\symc}\advance\symc by 1 }}
\newcommand\symblankrange[2]{\symc=#1 \@whilenum\symc<#2 \do{\expandafter\symblank\expandafter{\the\symc}\advance\symc by 1 }}
\makeatother
\symrange{"0250}{"02FF}\symactive{"2713}\symactive{"2714}\symactive{"2717}\symactive{"2718}\symactive{"2610}\symactive{"2611}\symactive{"2612}
\symrange{"2600}{"27BF}
\catcode"2705=13 \begingroup\lccode`\~="2705 \lowercase{\endgroup\def~}{{\symfont\char"2713}}\symblank{"26BD}\symblank{"26A1}
\symrange{"2460}{"257F}\symactive{"223C}\symactive{"2248}\symactive{"2260}
\symhyph{"2011}\symhyph{"2E1A}\symblankrange{"1F300}{"1FAFF}\symblank{"FE0F}
% Caractères chinois : WenQuanYi Zen Hei (sous-ensemble GB2312 + pinyin), gras/italique simulés
\setCJKmainfont{WenQuanYiZenHei-subset.ttf}[Path=../../../fonts/,AutoFakeBold=3,AutoFakeSlant=0.2]
\xeCJKsetup{CJKspace=false}
% Pinyin : voyelles à caron (ǎ ǐ ǒ ǔ ǖ ǘ ǚ ǜ) absentes de la police principale -> caractères actifs
\newfontface\pyfont{WenQuanYiZenHei-subset.ttf}[Path=../../../fonts/]
\newcommand\pyactive[1]{\catcode#1=13 \begingroup\lccode`\~=#1 \lowercase{\endgroup\def~}{{\pyfont\char#1}}}
\pyactive{"01CD}\pyactive{"01CE}\pyactive{"01CF}\pyactive{"01D0}\pyactive{"01D1}\pyactive{"01D2}\pyactive{"01D3}\pyactive{"01D4}\pyactive{"01D5}\pyactive{"01D6}\pyactive{"01D7}\pyactive{"01D8}\pyactive{"01D9}\pyactive{"01DA}\pyactive{"01DB}\pyactive{"01DC}
% Maths en sans-serif (Fira Math) avec chiffres et lettres droites de Plus
% Jakarta, pour que les formules se fondent dans le texte.
\usepackage[math-style=ISO,bold-style=ISO]{unicode-math}
\setmathfont{FiraMath-Regular.otf}[Path=../../../fonts/]
\setmathfont{PlusJakartaSans-Medium.ttf}[Path=../../../fonts/,range={up/num,up/{Latin,latin},\mathup}]
\usepackage{icomma} % virgule décimale sans espace en mode math
% Caractères mathématiques Unicode absents des polices de texte (Plus Jakarta,
% Archivo Black) : rendus par leur équivalent mathématique, en texte comme en
% formule — sinon ils s'affichent en carré vide (ex. « Arithmétique dans ℤ »).
\usepackage{newunicodechar}
\newunicodechar{ℝ}{\ensuremath{\mathbb{R}}}
\newunicodechar{ℤ}{\ensuremath{\mathbb{Z}}}
\newunicodechar{ℂ}{\ensuremath{\mathbb{C}}}
\newunicodechar{ℕ}{\ensuremath{\mathbb{N}}}
\newunicodechar{ℚ}{\ensuremath{\mathbb{Q}}}
\newunicodechar{𝔻}{\ensuremath{\mathbb{D}}}
\newunicodechar{α}{\ensuremath{\alpha}}
\newunicodechar{β}{\ensuremath{\beta}}
\newunicodechar{γ}{\ensuremath{\gamma}}
\newunicodechar{δ}{\ensuremath{\delta}}
\newunicodechar{ε}{\ensuremath{\varepsilon}}
\newunicodechar{θ}{\ensuremath{\theta}}
\newunicodechar{λ}{\ensuremath{\lambda}}
\newunicodechar{μ}{\ensuremath{\mu}}
\newunicodechar{σ}{\ensuremath{\sigma}}
\newunicodechar{τ}{\ensuremath{\tau}}
\newunicodechar{φ}{\ensuremath{\varphi}}
\newunicodechar{ψ}{\ensuremath{\psi}}
\newunicodechar{ω}{\ensuremath{\omega}}
\newunicodechar{Σ}{\ensuremath{\Sigma}}
% majuscules produites par \MakeUppercase dans les titres (α-aminés -> Α)
\newunicodechar{Α}{\ensuremath{\alpha}}
\newunicodechar{Β}{\ensuremath{\beta}}
\newunicodechar{Γ}{\ensuremath{\Gamma}}
\newunicodechar{Δ}{\ensuremath{\Delta}}
\newunicodechar{Ε}{\ensuremath{\varepsilon}}
\newunicodechar{Θ}{\ensuremath{\Theta}}
\newunicodechar{Λ}{\ensuremath{\Lambda}}
\newunicodechar{Μ}{\ensuremath{\mu}}
\newunicodechar{Τ}{\ensuremath{\tau}}
\newunicodechar{Φ}{\ensuremath{\Phi}}
\newunicodechar{Ψ}{\ensuremath{\Psi}}
\newunicodechar{Ω}{\ensuremath{\Omega}}
\newunicodechar{↦}{\ensuremath{\mapsto}}
\newunicodechar{∈}{\ensuremath{\in}}
\newunicodechar{∉}{\ensuremath{\notin}}
\newunicodechar{∧}{\ensuremath{\wedge}}
\newunicodechar{⇔}{\ensuremath{\Leftrightarrow}}
\newunicodechar{⟺}{\ensuremath{\Longleftrightarrow}}
\newunicodechar{⇒}{\ensuremath{\Rightarrow}}
\newunicodechar{⟹}{\ensuremath{\Longrightarrow}}
\newunicodechar{∘}{\ensuremath{\circ}}
\newunicodechar{∪}{\ensuremath{\cup}}
\newunicodechar{∩}{\ensuremath{\cap}}
\newunicodechar{≡}{\ensuremath{\equiv}}
\newunicodechar{⊂}{\ensuremath{\subset}}
\newunicodechar{⊥}{\ensuremath{\perp}}
\newunicodechar{∃}{\ensuremath{\exists}}
\newunicodechar{∀}{\ensuremath{\forall}}
\newunicodechar{∛}{\ensuremath{\sqrt[3]{\,}}}
\newunicodechar{∜}{\ensuremath{\sqrt[4]{\,}}}
\newunicodechar{⃗}{\ensuremath{{}^{\to}}}
\newunicodechar{ⁿ}{\ifmmode^{n}\else\textsuperscript{\ensuremath{n}}\fi}
\newunicodechar{ˣ}{\ifmmode^{x}\else\textsuperscript{\ensuremath{x}}\fi}
\newunicodechar{ᵇ}{\ifmmode^{b}\else\textsuperscript{\ensuremath{b}}\fi}
\newunicodechar{ʳ}{\ifmmode^{r}\else\textsuperscript{\ensuremath{r}}\fi}
\newunicodechar{ᵘ}{\ifmmode^{u}\else\textsuperscript{\ensuremath{u}}\fi}
\newunicodechar{ʸ}{\ifmmode^{y}\else\textsuperscript{\ensuremath{y}}\fi}
\newunicodechar{ᵃ}{\ifmmode^{a}\else\textsuperscript{\ensuremath{a}}\fi}
\newunicodechar{ᵉ}{\ifmmode^{e}\else\textsuperscript{\ensuremath{e}}\fi}
\newunicodechar{ᵏ}{\ifmmode^{k}\else\textsuperscript{\ensuremath{k}}\fi}
\newunicodechar{ᵖ}{\ifmmode^{p}\else\textsuperscript{\ensuremath{p}}\fi}
\newunicodechar{ᴺ}{\ifmmode^{N}\else\textsuperscript{\ensuremath{N}}\fi}
\newunicodechar{ᶜ}{\ifmmode^{c}\else\textsuperscript{\ensuremath{c}}\fi}
\newunicodechar{ʰ}{\ifmmode^{h}\else\textsuperscript{\ensuremath{h}}\fi}
\newunicodechar{ᵠ}{\ifmmode^{q}\else\textsuperscript{\ensuremath{q}}\fi}
\newunicodechar{ₙ}{\ifmmode_{n}\else\textsubscript{\ensuremath{n}}\fi}
\newunicodechar{ₐ}{\ifmmode_{a}\else\textsubscript{\ensuremath{a}}\fi}
\newunicodechar{ᵢ}{\ifmmode_{i}\else\textsubscript{\ensuremath{i}}\fi}
\newunicodechar{ᵣ}{\ifmmode_{r}\else\textsubscript{\ensuremath{r}}\fi}
\newunicodechar{ₖ}{\ifmmode_{k}\else\textsubscript{\ensuremath{k}}\fi}
\newunicodechar{ᵦ}{\ifmmode_{\beta}\else\textsubscript{\ensuremath{\beta}}\fi}
\newunicodechar{ₓ}{\ifmmode_{x}\else\textsubscript{\ensuremath{x}}\fi}
\newfontfamily\popdisplay{ArchivoBlack-Regular.ttf}[Path=../../../fonts/]
\newfontfamily\poplogo{SpaceGrotesk-Bold.ttf}[Path=../../../fonts/]
\newfontfamily\popsemi{PlusJakartaSans-SemiBold.ttf}[Path=../../../fonts/]
\newfontfamily\popxbold{PlusJakartaSans-ExtraBold.ttf}[Path=../../../fonts/]
\newfontfamily\popxboldls{PlusJakartaSans-ExtraBold.ttf}[Path=../../../fonts/,LetterSpace=3.0]

\setlist[enumerate]{leftmargin=1.7em,itemsep=5pt,topsep=4pt}
\setlist[itemize]{leftmargin=1.7em,itemsep=4pt,topsep=4pt,label={\color{poporange}\textbullet}}
\setlength{\parindent}{0pt}
\setlength{\parskip}{5pt}
\renewcommand{\arraystretch}{1.25}

% pgfplots : style Pop par défaut
\pgfplotsset{
  every axis/.append style={axis line style={popink,line width=1pt},tick style={popink},
    grid style={popline,line width=0.5pt},label style={font=\small},tick label style={font=\footnotesize}},
  popcurve/.style={poporange,line width=1.8pt,no markers,smooth},
}
% Même style disponible dans un \draw TikZ simple (hors pgfplots)
\tikzset{popcurve/.style={poporange,line width=1.8pt}}
\tikzset{popgrid/.style={popline,very thin}}
\colorlet{popcurve}{poporange} % le nom du style est parfois utilisé comme couleur

% ============================================================
% Pill / squiggle
% ============================================================
\newcommand{\poppill}[3]{%
  \tikz[baseline=-0.55ex]\node[draw=popink,line width=1.2pt,rounded corners=2.6pt,%
    fill=#2,inner xsep=6.5pt,inner ysep=3pt]{%
    \popxboldls\fontsize{7}{7}\selectfont\color{#3}\MakeUppercase{#1}};%
}
\newcommand{\popsquiggle}[1]{%
  \tikz[baseline=-0.4ex]{
    \draw[#1,line width=2.6pt,line cap=round]
      plot [smooth] coordinates {(0,0.09) (0.34,0.01) (0.68,0.21) (1.02,0.01) (1.36,0.10)};
  }%
}
% Mot-clé surligné (marqueur orange sous le texte)
\newcommand{\cle}[1]{{\popxbold\color{popdark}#1}}

% ============================================================
% En-tête / pied
% ============================================================
\pagestyle{fancy}
\fancyhf{}
\renewcommand{\headrulewidth}{0pt}
\renewcommand{\footrulewidth}{0pt}
\lhead{\raisebox{-0.15ex}{\poplogo\fontsize{13}{13}\selectfont\color{popink}OnBuch\color{poporange}+}}
\rhead{\poppill{\DOCMATIERE\ \textbullet\ \DOCNIVEAU\ \textbullet\ Leçon \DOCLECON}{white}{popink}}
\lfoot{\popxbold\fontsize{7.6}{7.6}\selectfont\color{popmuted}\MakeUppercase{Cours signé OnBuch+}\ \textbullet\ {\popsemi\DOCTITRE}}
\rfoot{\tikz[baseline=-0.6ex]\node[rounded corners=8.5pt,fill=popink,minimum height=17pt,minimum width=17pt,inner xsep=5pt,inner sep=0]{\popxbold\fontsize{8}{8}\selectfont\color{popcream}\thepage\,/\,\pageref*{LastPage}};}

% ============================================================
% Titres
% ============================================================
\newcommand{\chaptitre}[2]{%
  \begin{tcolorbox}[enhanced,colback=poporange,colframe=popink,boxrule=2.6pt,arc=11pt,
      drop shadow=popink,left=15pt,right=15pt,top=12pt,bottom=11pt,before skip=3pt,after skip=18pt]
    \poppill{#2}{popink}{popcream}\par\vspace{9pt}
    {\raggedright\hyphenpenalty=10000\exhyphenpenalty=10000\popdisplay\fontsize{22}{24}\selectfont\color{popcream}\MakeUppercase{#1}\par}
  \end{tcolorbox}
}
% Section numérotée : pastille orange avec le numéro + titre Archivo
\newcounter{popsec}
\newcommand{\coursec}[1]{%
  \stepcounter{popsec}\setcounter{popsub}{0}%
  \par\vspace{16pt}\noindent
  \begin{minipage}{\linewidth}
  \tikz[baseline=-0.7ex]\node[draw=popink,line width=1.6pt,fill=poporange,rounded corners=4pt,minimum size=22pt,inner sep=0]{\popdisplay\fontsize{12}{12}\selectfont\color{popcream}\thepopsec};%
  \hspace{8pt}{\popdisplay\fontsize{15}{17}\selectfont\color{popink}\MakeUppercase{#1}}\par
  \vspace{4pt}\noindent{\color{poporange}\rule{36pt}{3.6pt}}
  \end{minipage}\par\nopagebreak\vspace{8pt}
}
\newcounter{popsub}
\newcommand{\courssub}[1]{%
  \stepcounter{popsub}%
  \par\vspace{9pt}\noindent{\popxbold\fontsize{11.5}{13}\selectfont\color{popink}{\color{poporange}\thepopsec.\thepopsub}\hspace{6pt}#1}\par\nopagebreak\vspace{4pt}
}

% ============================================================
% Boîtes pédagogiques — même recette : fond blanc, bordure épaisse colorée,
% ombre dure encre, badge plein. Titre optionnel [#1].
% ============================================================
\tcbset{popbase/.style={breakable,enhanced,colback=white,boxrule=2.3pt,arc=9pt,
  shadow={3pt}{-3pt}{0pt}{popink},left=14pt,right=14pt,top=11pt,bottom=11pt,before skip=11pt,after skip=15pt,
  fontupper=\popsemi\fontsize{10.6}{14.5}\selectfont\color{popink},
  fontlower=\popsemi\fontsize{10.6}{14.5}\selectfont\color{popink}}}
% Coche dessinée (le glyphe ✓ n'existe pas dans les polices du document)
\newcommand{\popcheck}[1]{\tikz[baseline=-0.5ex]\draw[#1,line width=1.6pt,line cap=round,line join=round](-0.13,0.0)--(-0.04,-0.09)--(0.13,0.1);}
\providecommand{\coche}{\popcheck{popgreen}}
\providecommand{\resultat}[1]{\par\smallskip\noindent\fcolorbox{popgreen}{popgreenL}{\parbox{\dimexpr\linewidth-2\fboxsep-2\fboxrule\relax}{\textbf{Résultat.} #1}}\par\smallskip}
\newcommand{\popboxhead}[3]{\poppill{#1}{#2}{white}\ifx\relax#3\relax\else\hspace{6pt}{\popxbold\fontsize{10.6}{12}\selectfont\color{popink}#3}\fi\par\vspace{7pt}}

\newtcolorbox{aretenir}[1][]{popbase,colframe=popgreen,before upper={\popboxhead{À retenir}{popgreen}{#1}}}
% Texte en chinois (document de lecture, dialogue, extrait) : cadre
% d'étude. Un titre optionnel (source, auteur) s'affiche dans l'en-tête.
\newtcolorbox{textebox}[1][]{popbase,colframe=popblue,before upper={\popboxhead{文本 · Texte}{popblue}{#1}},after upper={}}
% Marques ✓ / ✗ (forme correcte / fautive) souvent utilisées par les modèles
\providecommand{\cross}{\textcolor{poppink}{\ensuremath{\times}}}
\providecommand{\xmark}{\cross}\providecommand{\crossmark}{\cross}\providecommand{\cmark}{\ensuremath{\checkmark}}
% Ligne de réponse / justification à compléter (inventée par les modèles)
\providecommand{\justification}[1]{\par\noindent\textit{Justification~:} \dotfill\par}
% \textipa (package tipa non chargé) : la police ne couvre pas l'API -> simple passage
\providecommand{\textipa}[1]{#1}
% Soulignement (ulem non chargé) : \uline, \ul
\providecommand{\uline}[1]{\underline{#1}}\providecommand{\ul}[1]{\underline{#1}}
% \glossaire{\item ...} inventée par les modèles : liste de glossaire
\providecommand{\glossaire}[1]{\begin{itemize}#1\end{itemize}}
% \alert (beamer) inventée par les modèles : texte mis en évidence
\providecommand{\alert}[1]{\textbf{\textcolor{poppink}{#1}}}
% variantes ASCII d'ordinaux écrites en macro
\providecommand{\iere}{\textsuperscript{re}}\providecommand{\ieme}{\textsuperscript{e}}\providecommand{\ere}{\textsuperscript{re}}\providecommand{\eme}{\textsuperscript{e}}\providecommand{\er}{\textsuperscript{er}}\providecommand{\ie}{\textsuperscript{e}}
% Ordinaux français écrits en macro par les modèles : 1\ière, 3\ième
\newcommand{\ière}{\textsuperscript{re}}\newcommand{\ième}{\textsuperscript{e}}
% \exemple (inventée par les modèles) : étiquette « Exemple : »
\providecommand{\exemple}{\textbf{Exemple~:}\ }
% \square utilisable aussi hors mode math (cases à cocher)
\AtBeginDocument{\let\squareMath\square\renewcommand{\square}{\ensuremath{\squareMath}}}
% Mot ou phrase en chinois à l'intérieur d'un paragraphe français
\newcommand{\zh}[1]{\ifmmode\text{#1}\else{#1}\fi}
\let\eng\zh \let\esp\zh \let\all\zh % alias tolérés
% flèches d'intonation inventées par les modèles
\providecommand{\textsearrow}{}\renewcommand{\textsearrow}{\ensuremath{\searrow}}\providecommand{\textnearrow}{}\renewcommand{\textnearrow}{\ensuremath{\nearrow}}
% \fr{..} : mot français (inventé par les modèles)
\providecommand{\fr}[1]{\textit{#1}}
% zhbox : encadré de texte chinois inventé par les modèles
\newenvironment{zhbox}{\begin{quote}}{\end{quote}}
% \vs : « versus » inventé par les modèles
\providecommand{\vs}{\textit{vs}\ }
% \blank : case à compléter inventée par les modèles
\providecommand{\blank}[1][]{\underline{\hspace{1.6em}}}
\newtcolorbox{exemplebox}[1][]{popbase,colframe=poporange,before upper={\popboxhead{Exemple}{poporange}{#1}}}
\newtcolorbox{definition}[1][]{popbase,colframe=popblue,before upper={\popboxhead{Définition}{popblue}{#1}}}
\newtcolorbox{propriete}[1][]{popbase,colframe=poppurple,before upper={\popboxhead{Propriété}{poppurple}{#1}}}
\newtcolorbox{methode}[1][]{popbase,colframe=popgold,before upper={\popboxhead{Méthode}{popgold}{#1}}}
\newtcolorbox{attention}[1][]{popbase,colframe=poppink,before upper={\popboxhead{Attention}{poppink}{#1}}}
\newtcolorbox{experience}[1][]{popbase,colframe=popblue,colback=popblueL,before upper={\popboxhead{Expérience}{popblue}{#1}}}
% « Le savais-tu ? » : carte violette pleine (contexte, culture, Cameroun)
\newtcolorbox{savaistu}[1][]{popbase,colback=poppurple,colframe=popink,
  fontupper=\popsemi\fontsize{10.4}{14.2}\selectfont\color{white},
  before upper={\poppill{Le savais-tu\,?}{white}{poppurple}\ifx\relax#1\relax\else\hspace{6pt}{\popxbold\color{white}#1}\fi\par\vspace{7pt}}}
% Exercice résolu : énoncé en haut, \tcblower puis la solution sur fond crème
\newcounter{popexo}\newcounter{popexr}
\newtcolorbox{exoresolu}[1][]{popbase,colframe=poporange,colbacklower=popcream,
  segmentation style={popink,line width=1.2pt,dashed},
  before upper={\stepcounter{popexr}\popboxhead{Exercice résolu \thepopexr}{poporange}{#1}},
  before lower={\poppill{Solution}{popink}{popcream}\par\vspace{6pt}}}
% Exercice (non résolu en place) — la correction est dans la partie Corrigés
\newtcolorbox{exercice}[1][]{popbase,colframe=popink,
  before upper={\stepcounter{popexo}\popboxhead{Exercice \thepopexo}{popink}{#1}}}
% Corrigé
\newtcolorbox{corrige}[1][]{popbase,colframe=popgreen,colback=popgreenL,
  before upper={\popboxhead{Corrigé}{popgreen}{#1}}}
% Objectifs / prérequis / bilan
\newtcolorbox{objectifs}{popbase,colframe=popink,colback=poporangeL,
  before upper={\popboxhead{Objectifs de la leçon}{popink}{}}}
\newtcolorbox{prerequis}{popbase,colframe=popink,colback=popblueL,
  before upper={\popboxhead{Prérequis}{popblue}{}}}
\newtcolorbox{bilan}{popbase,colframe=popink,colback=white,boxrule=2.8pt,
  before upper={\popboxhead{Fiche bilan}{poporange}{L'essentiel à connaître pour l'examen}}}
% Figure encadrée avec légende
\newtcolorbox{popfigure}[1][]{enhanced,breakable=false,colback=white,colframe=popline,boxrule=1.4pt,arc=8pt,
  left=10pt,right=10pt,top=10pt,bottom=8pt,before skip=10pt,after skip=12pt,halign=center,
  lower separated=false,halign lower=center,
  fontlower=\popsemi\fontsize{9}{11.5}\selectfont\color{popmuted},
  #1}
\newcommand{\legende}[1]{\par\vspace{6pt}{\popsemi\fontsize{9}{11.5}\selectfont\color{popmuted}#1}}

% ============================================================
% Signature OnBuch+
% ============================================================
\newcommand{\poplogoinline}[1]{{\poplogo\fontsize{#1}{#1}\selectfont\color{popink}OnBuch\color{poporange}+}}
\newcommand{\signatureonbuch}{%
  \begin{tcolorbox}[enhanced,colback=popink,colframe=popink,boxrule=2.2pt,arc=10pt,
      drop shadow=poporange,left=14pt,right=14pt,top=10pt,bottom=10pt,before skip=0pt,after skip=0pt]
    \tikz[baseline=-0.7ex]\node[circle,fill=popgreen,draw=popcream,line width=1.3pt,minimum size=22pt,inner sep=0]{\popcheck{white}};%
    \hspace{9pt}\begin{minipage}[c]{0.62\linewidth}
      {\popxbold\fontsize{12}{13}\selectfont\color{popcream}Signé\ \textbullet\ {\poplogo OnBuch\color{poporange}+}}\par\vspace{2pt}
      {\raggedright\popsemi\fontsize{8}{10.5}\selectfont\color{popline}\MakeUppercase{Équipe pédagogique OnBuch+}\\\MakeUppercase{Conforme au programme officiel MINESEC}\par}
    \end{minipage}\hfill
    {\poplogo\fontsize{22}{22}\selectfont\color{popcream}OnBuch\color{poporange}+}
  \end{tcolorbox}%
}

% ============================================================
% Page de garde
% #1 titre · #2 matière · #3 niveau · #4 module · #5 n° leçon · #6 durée ·
% #7 description / objectifs (texte ou itemize)
% ============================================================
\newcommand{\pagegarde}[7]{%
  \thispagestyle{empty}
  \newgeometry{a4paper,margin=1.7cm,top=1.6cm,bottom=1.4cm}%
  \begin{tikzpicture}[remember picture,overlay]
    \fill[poporange!18] ([xshift=-3.2cm,yshift=-3.6cm]current page.north east) circle (4.6cm);
    \fill[poppurple!14] ([xshift=2.4cm,yshift=7.6cm]current page.south west) circle (3.8cm);
  \end{tikzpicture}%
  {\poplogo\fontsize{30}{30}\selectfont\color{popink}OnBuch\color{poporange}+}\hfill
  \raisebox{0.6ex}{\poppill{Cours signé \textbullet\ Premium}{popink}{popcream}}\par
  \vspace{1pt}\popsquiggle{poppurple}\par
  \vspace{8pt}
  \begin{tcolorbox}[enhanced,colback=poporange,colframe=popink,boxrule=2.8pt,arc=13pt,
      drop shadow=popink,left=18pt,right=18pt,top=12pt,bottom=13pt,after skip=10pt]
    \poppill{#2 \textbullet\ #3}{popink}{popcream}\hspace{5pt}\poppill{#4}{white}{popink}\par\vspace{8pt}
    {\popdisplay\fontsize{14}{14}\selectfont\color{popink}LEÇON #5}\par\vspace{4pt}
    {\raggedright\hyphenpenalty=10000\exhyphenpenalty=10000\popdisplay\fontsize{25}{27}\selectfont\color{popcream}\MakeUppercase{#1}\par}
    \vspace{8pt}
    {\popxbold\fontsize{9.5}{9.5}\selectfont\color{popink}\MakeUppercase{Durée conseillée : #6}}
  \end{tcolorbox}
  \begin{tcolorbox}[enhanced,colback=white,colframe=popink,boxrule=2.2pt,arc=10pt,
      drop shadow=popink,left=16pt,right=16pt,top=10pt,bottom=10pt,after skip=8pt]
    \poppill{Dans cette leçon}{poppurple}{white}\par\vspace{5pt}
    \popsemi\fontsize{9.3}{12.6}\selectfont\color{popink}#7
  \end{tcolorbox}
  \noindent\begin{minipage}[t]{0.487\linewidth}
    \begin{tcolorbox}[enhanced,colback=popgreen,colframe=popink,boxrule=2.2pt,arc=10pt,
        drop shadow=popink,left=11pt,right=11pt,top=10pt,bottom=10pt,height=3.1cm,valign=center]
      \tcbox[colback=white,colframe=popink,boxrule=1.3pt,arc=5pt,boxsep=0pt,nobeforeafter,
        left=3pt,right=3pt,top=3pt,bottom=3pt,valign=center]{\qrcode[height=1.9cm]{https://play.google.com/store/apps/details?id=com.luvvix.onbuch&hl=fr}}%
      \hspace{8pt}\begin{minipage}[c]{0.5\linewidth}
        {\popdisplay\fontsize{11}{12.5}\selectfont\color{white}\MakeUppercase{Télécharge OnBuch}}\par\vspace{4pt}
        {\raggedright\popsemi\fontsize{8.4}{11}\selectfont\color{white}Gratuit \textbullet\ Google Play\\Tuteur IA, annales, concours, résultats\par}
      \end{minipage}
    \end{tcolorbox}
  \end{minipage}\hfill
  \begin{minipage}[t]{0.487\linewidth}
    \begin{tcolorbox}[enhanced,colback=poppurple,colframe=popink,boxrule=2.2pt,arc=10pt,
        drop shadow=popink,left=11pt,right=11pt,top=10pt,bottom=10pt,height=3.1cm,valign=center]
      \tikz[baseline=-0.7ex]\node[circle,fill=white,draw=popink,line width=1.3pt,minimum size=1.6cm,inner sep=0]{\popdisplay\fontsize{24}{24}\selectfont\color{poppurple}+};%
      \hspace{8pt}\begin{minipage}[c]{0.6\linewidth}
        {\popdisplay\fontsize{11}{12.5}\selectfont\color{white}\MakeUppercase{Passe à OnBuch+}}\par\vspace{4pt}
        {\raggedright\popsemi\fontsize{8.4}{11}\selectfont\color{white}Tous les cours signés, Tuteur IA illimité, fiches PDF — onglet OnBuch+ de l'app\par}
      \end{minipage}
    \end{tcolorbox}
  \end{minipage}
  \vspace{8pt}
  \popcontenu
  \vfill
  \signatureonbuch
  \restoregeometry
  \newpage
}

% Rangée « Ce cours contient » (page de garde)
\newcommand{\poptile}[3]{% #1 couleur #2 gros texte #3 légende
  \begin{tcolorbox}[enhanced,colback=#1,colframe=popink,boxrule=2pt,arc=9pt,shadow={2.5pt}{-2.5pt}{0pt}{popink},
      left=6pt,right=6pt,top=9pt,bottom=9pt,halign=center,valign=center,height=1.75cm,width=\linewidth,nobeforeafter]
    {\popdisplay\fontsize{12.5}{13}\selectfont\color{white}\MakeUppercase{#2}}\par\vspace{3pt}
    {\centering\popsemi\fontsize{7.8}{9.5}\selectfont\color{white}#3\par}
  \end{tcolorbox}}
\newcommand{\popcontenu}{%
  \par\noindent{\popxboldls\fontsize{7.5}{7.5}\selectfont\color{popmuted}CE COURS SIGNÉ CONTIENT}\par\vspace{7pt}
  \noindent\begin{minipage}[t]{0.235\linewidth}\poptile{popblue}{Cours}{complet, illustré, pas à pas}\end{minipage}\hfill
  \begin{minipage}[t]{0.235\linewidth}\poptile{poporange}{Méthodes}{et exercices résolus}\end{minipage}\hfill
  \begin{minipage}[t]{0.235\linewidth}\poptile{poppink}{Exercices}{type examen + corrigés}\end{minipage}\hfill
  \begin{minipage}[t]{0.235\linewidth}\poptile{popgreen}{Fiche bilan}{l'essentiel à retenir}\end{minipage}\par}

% Sommaire
\newtcolorbox{sommaire}{enhanced,colback=white,colframe=popink,boxrule=2.2pt,arc=10pt,
  drop shadow=popink,left=18pt,right=18pt,top=13pt,bottom=13pt,before skip=6pt,after skip=20pt,
  before upper={\poppill{Sommaire}{poppurple}{white}\par\vspace{9pt}\popsemi\fontsize{10.6}{14.5}\selectfont\color{popink}}}

% Signature de fin de cours — poussée tout en bas de la dernière page
\newcommand{\findecours}{%
  \par\vfill\nopagebreak
  \begin{center}\popsquiggle{poporange}\end{center}\vspace{4pt}
  \signatureonbuch
}
\AtBeginDocument{\def\llbracket{\mathopen{[\mkern-3.5mu[}}\def\rrbracket{\mathclose{]\mkern-3.5mu]}}} % intervalles d'entiers (glyphe ⟦ absent de la police)
% Glyphes absents de Fira Math : redéfinis à partir de glyphes disponibles
\AtBeginDocument{%
  \def\setminus{\mathbin{\backslash}}\let\smallsetminus\setminus
  \def\vdots{\mathord{\vbox{\baselineskip3.2pt\lineskiplimit0pt\kern4pt\hbox{.}\hbox{.}\hbox{.}}}}%
  \def\ddots{\mathinner{\mkern1mu\raise6.5pt\hbox{.}\mkern2mu\raise3.6pt\hbox{.}\mkern2mu\raise0.7pt\hbox{.}\mkern1mu}}%
  \def\triangle{\mathord{\scalebox{0.9}{$\symup{\Delta}$}}}%
  \def\nabla{\mathord{\raisebox{\depth}{\scalebox{1}[-1]{$\symup{\Delta}$}}}}%
  \def\checkmark{\popcheck{popdark}}%
  \def\star{\mathbin{\ast}}\def\bigstar{\mathord{\ast}}%
  \def\diamond{\mathbin{\scalebox{0.6}{$\Diamond$}}}%
  \def\bigcup{\mathop{\mathchoice{\raisebox{-0.3em}{\scalebox{1.7}{$\cup$}}}{\scalebox{1.25}{$\cup$}}{\cup}{\cup}}\displaylimits}%
  \def\bigcap{\mathop{\mathchoice{\raisebox{-0.3em}{\scalebox{1.7}{$\cap$}}}{\scalebox{1.25}{$\cap$}}{\cap}{\cap}}\displaylimits}%
  \def\lesseqqgtr{\mathrel{\lessgtr}}\def\bigoplus{\mathop{\oplus}\displaylimits}%
}
\providecommand{\ssi}{si et seulement si} % abréviation fréquente
\AtBeginDocument{\providecommand{\wideparen}{\overparen}} % arc de cercle

\begin{document}

\pagegarde{Leçon 2 — Grammaire : 第一……第二…. ; 不但……而且….}{Chinois}{1ère A}{Module 1 : Vie familiale et intégration sociale}{ZHA02}{2 h}{Cette leçon de grammaire présente deux structures essentielles du chinois : 第一……第二…… pour énumérer des arguments et 不但……而且…… pour exprimer la gradation « non seulement... mais aussi... ». À travers l'observation d'exemples, des schémas, la comparaison avec le français et des exercices progressifs, les élèves apprennent à construire des phrases argumentées et nuancées.
\vspace{4pt}
\begin{multicols}{2}\small
\begin{itemize}
\item À la fin de cette leçon, je sais reconnaître les structures 第一……第二…… et 不但……而且…… dans un texte
\item À la fin de cette leçon, je sais construire une énumération ordonnée avec 第一、第二、第三
\item À la fin de cette leçon, je sais exprimer la gradation avec 不但……而且……
\item À la fin de cette leçon, je sais placer correctement 不但 et 而且 selon que les deux propositions ont le même sujet ou des sujets différents
\item À la fin de cette leçon, je sais éviter les erreurs fréquentes des francophones (ordre des mots, oubli de 而且, calque du français)
\item À la fin de cette leçon, je sais transformer des phrases simples en phrases à gradation
\end{itemize}\end{multicols}}

\begin{sommaire}
\begin{multicols}{2}
\begin{enumerate}
\item Observation guidée d'exemples
\item Structure 1 : 第一……第二…… (énumération)
\item Structure 2 : 不但……而且…… (gradation)
\item Comparaison avec le français et pièges des francophones
\item Exercices d'application et de transformation
\item Synthèse et prolongement vers l'examen
\item Activité d'intégration
\item Méthodes et exercices résolus
\item Exercices
\item Corrigés des exercices
\item Fiche bilan
\end{enumerate}
\end{multicols}
\end{sommaire}

\begin{objectifs}
\begin{itemize}
\item À la fin de cette leçon, je sais reconnaître les structures 第一……第二…… et 不但……而且…… dans un texte
\item À la fin de cette leçon, je sais construire une énumération ordonnée avec 第一、第二、第三
\item À la fin de cette leçon, je sais exprimer la gradation avec 不但……而且……
\item À la fin de cette leçon, je sais placer correctement 不但 et 而且 selon que les deux propositions ont le même sujet ou des sujets différents
\item À la fin de cette leçon, je sais éviter les erreurs fréquentes des francophones (ordre des mots, oubli de 而且, calque du français)
\item À la fin de cette leçon, je sais transformer des phrases simples en phrases à gradation
\item À la fin de cette leçon, je sais utiliser ces deux structures dans un court texte argumentatif (présentation, lettre, exposé)
\item À la fin de cette leçon, je sais traduire correctement du français vers le chinois des phrases contenant « premièrement/deuxièmement » et « non seulement/mais aussi »
\end{itemize}
\end{objectifs}

\begin{prerequis}
\begin{itemize}
\item Structure de base de la phrase chinoise Sujet-Verbe-Objet (Seconde)
\item Adverbes de liaison simples : 也, 还, 都 (Seconde)
\item Nombres et numérotation : 一、二、三 (Seconde)
\item Pronoms personnels et vocabulaire de la vie familiale et scolaire (Module 1)
\end{itemize}
\end{prerequis}

\newpage

\coursec{Observation guidée d'exemples}

\courssub{Situation d'accroche : la rentrée au Lycée de Yaoundé}

C'est la rentrée au Lycée de Yaoundé. Dans la classe de Première A, le professeur principal fait le tour des élèves et pose à chacun la même question : « Pourquoi apprends-tu le chinois ? » Amina lève la main et répond avec assurance : « Premièrement, j'aime la culture chinoise ; deuxièmement, le chinois est utile pour le commerce. » Son camarade Rodrigue ajoute : « Le chinois est non seulement intéressant, mais aussi très utile. »

Deux idées, deux façons de les organiser : Amina \emph{énumère} ses raisons (premièrement, deuxièmement), Rodrigue \emph{gradue} son argument (non seulement\ldots mais aussi). Comment dire cela en chinois ? C'est exactement ce que nous allons apprendre avec les deux structures \cle{第一\ldots 第二\ldots} et \cle{不但\ldots 而且\ldots}.

\textbf{Problématique de la leçon :} comment énumérer des arguments et exprimer une gradation en chinois, sans calquer le français mot à mot ?

\courssub{Texte support : la réponse d'Amina}

Lisons d'abord le petit texte suivant, qui reprend la scène de la rentrée. Lis-le deux fois : une première fois en regardant les caractères, une deuxième fois en suivant le pinyin à voix haute.

\begin{textebox}[À la rentrée, pourquoi le chinois ?]
我学中文，第一，我喜欢中国文化；第二，中文对工作很有用。他不但会说汉语，而且会写汉字。\\
Wǒ xué Zhōngwén, dì-yī, wǒ xǐhuan Zhōngguó wénhuà; dì-èr, Zhōngwén duì gōngzuò hěn yǒuyòng. Tā bùdàn huì shuō Hànyǔ, érqiě huì xiě Hànzì.\\
« J'apprends le chinois : premièrement, j'aime la culture chinoise ; deuxièmement, le chinois est très utile pour le travail. Non seulement il sait parler chinois, mais il sait aussi écrire les caractères. »
\end{textebox}

\textbf{Mini-glossaire du texte}

\begin{center}
\begin{tabularx}{\linewidth}{|X|X|X|X|}
\hline
\rowcolor{popblueL} Caractères & Pinyin & Nature & Traduction \\
\hline
学中文 & xué Zhōngwén & verbe + nom & apprendre le chinois \\
\hline
文化 & wénhuà & nom & culture \\
\hline
对\ldots 有用 & duì\ldots yǒuyòng & structure & être utile pour\ldots \\
\hline
工作 & gōngzuò & nom & travail \\
\hline
会说 & huì shuō & auxiliaire + verbe & savoir parler \\
\hline
写汉字 & xiě Hànzì & verbe + nom & écrire les caractères \\
\hline
\end{tabularx}
\end{center}

\courssub{Les quatre exemples en contexte}

Voici maintenant les quatre phrases de départ de la leçon. Chacune est tirée de la vie scolaire ou familiale camerounaise. Pour chacune, lis le chinois, puis le pinyin, puis la traduction, et repère les mots qui se répètent.

\begin{exemplebox}[Exemple 1 — Énumérer ses raisons]
\zh{我学中文，第一，我喜欢中国文化；第二，中文对工作很有用。}\\
\emph{Wǒ xué Zhōngwén, dì-yī, wǒ xǐhuan Zhōngguó wénhuà; dì-èr, Zhōngwén duì gōngzuò hěn yǒuyòng.}\\
« J'apprends le chinois : premièrement, j'aime la culture chinoise ; deuxièmement, le chinois est très utile pour le travail. »
\end{exemplebox}

\begin{exemplebox}[Exemple 2 — Graduer : un seul sujet]
\zh{他不但会说汉语，而且会写汉字。}\\
\emph{Tā bùdàn huì shuō Hànyǔ, érqiě huì xiě Hànzì.}\\
« Non seulement il sait parler chinois, mais il sait aussi écrire les caractères. »
\end{exemplebox}

\begin{exemplebox}[Exemple 3 — Graduer : deux sujets différents]
\zh{不但我喜欢中国菜，而且我弟弟也喜欢。}\\
\emph{Bùdàn wǒ xǐhuan Zhōngguó cài, érqiě wǒ dìdi yě xǐhuan.}\\
« Non seulement moi j'aime la cuisine chinoise, mais mon petit frère l'aime aussi. »
\end{exemplebox}

\begin{exemplebox}[Exemple 4 — Contexte familial à Douala]
\zh{我们家第一离学校很近，第二离市场也不远，所以很方便。}\\
\emph{Wǒmen jiā dì-yī lí xuéxiào hěn jìn, dì-èr lí shìchǎng yě bù yuǎn, suǒyǐ hěn fāngbiàn.}\\
« Notre maison, premièrement, est très proche de l'école ; deuxièmement, elle n'est pas loin du marché non plus, donc c'est très pratique. »
\end{exemplebox}

\courssub{Question guide : observe la place de 不但}

\begin{experience}[Observation en binôme (5 minutes)]
Par groupes de deux, répondez aux questions suivantes en regardant les exemples 2 et 3 :
\begin{enumerate}
\item Dans l'exemple 2, où se trouve \zh{不但} par rapport au sujet \zh{他} ?
\item Dans l'exemple 3, où se trouve \zh{不但} par rapport au sujet \zh{我} ?
\item Quelle différence y a-t-il entre les sujets des deux phrases ? (Combien de personnes différentes sont concernées ?)
\item Dans l'exemple 1, quels mots-outils introduisent chaque argument ? Ont-ils un équivalent exact en français ?
\end{enumerate}
\end{experience}

\textbf{Réponses attendues (dégagement collectif).} Dans l'exemple 2, il n'y a qu'\emph{un seul sujet} (\zh{他}) : \zh{不但} se place \emph{après} le sujet. Dans l'exemple 3, il y a \emph{deux sujets différents} (\zh{我} et \zh{我弟弟}) : \zh{不但} se place \emph{avant} le premier sujet. C'est la règle d'or de la place de \zh{不但}, que nous étudierons en détail dans la section 3.

\courssub{Tableau d'observation : chinois et français face à face}

\begin{popfigure}
\begin{tikzpicture}[
  cell/.style={draw=popline, rounded corners=2pt, align=left, inner sep=6pt},
  head/.style={fill=popblue, text=white, font=\bfseries, align=center, inner sep=6pt, rounded corners=2pt},
  tool/.style={text=poporange, font=\bfseries}
]
\node[head, minimum width=7.2cm] (h1) at (0,0) {Phrase en chinois};
\node[head, minimum width=6.2cm] (h2) at (6.9,0) {Traduction française};
\node[cell, minimum width=7.2cm, anchor=north west, text width=6.8cm] (a1) at (-3.6,-0.5)
  {我学中文，\textcolor{poporange}{\textbf{第一}}，我喜欢中国文化；\textcolor{poporange}{\textbf{第二}}，中文对工作很有用。};
\node[cell, minimum width=6.2cm, anchor=north west, text width=5.8cm] (a2) at (3.8,-0.5)
  {J'apprends le chinois : \emph{premièrement}, j'aime la culture chinoise ; \emph{deuxièmement}, le chinois est utile pour le travail.};
\node[cell, minimum width=7.2cm, anchor=north west, text width=6.8cm] (b1) at (-3.6,-3.1)
  {他\textcolor{poporange}{\textbf{不但}}会说汉语，\textcolor{poporange}{\textbf{而且}}会写汉字。};
\node[cell, minimum width=6.2cm, anchor=north west, text width=5.8cm] (b2) at (3.8,-3.1)
  {\emph{Non seulement} il sait parler chinois, \emph{mais} il sait \emph{aussi} écrire les caractères.};
\node[cell, minimum width=7.2cm, anchor=north west, text width=6.8cm] (c1) at (-3.6,-4.9)
  {\textcolor{poporange}{\textbf{不但}}我喜欢中国菜，\textcolor{poporange}{\textbf{而且}}我弟弟也喜欢。};
\node[cell, minimum width=6.2cm, anchor=north west, text width=5.8cm] (c2) at (3.8,-4.9)
  {\emph{Non seulement} moi j'aime la cuisine chinoise, \emph{mais} mon petit frère l'aime \emph{aussi}.};
\end{tikzpicture}
\legende{Les mots-outils 第一、第二、不但、而且 surlignés dans les phrases d'observation.}
\end{popfigure}

\courssub{Dégagement collectif des deux structures}

À partir des exemples, la classe dégage les deux schémas suivants. Ne les mémorise pas encore : contente-toi de vérifier qu'ils correspondent bien aux quatre exemples.

\begin{definition}[Structure 1 : l'énumération 第一……第二……]
\zh{第一} (dì-yī) et \zh{第二} (dì-èr) introduisent des arguments numérotés : « premièrement\ldots deuxièmement\ldots ». Chaque mot-outil se place \emph{en tête de sa proposition}, suivi d'une virgule chinoise \zh{，}. On peut continuer : \zh{第三} (troisièmement), \zh{第四} (quatrièmement)\ldots\\
Schéma : \textbf{énoncé général，第一 + proposition 1；第二 + proposition 2。}
\end{definition}

\begin{definition}[Structure 2 : la gradation 不但……而且……]
\zh{不但} (bùdàn) « non seulement » et \zh{而且} (érqiě) « mais aussi » relient deux propositions dont la seconde ajoute une information plus forte.\\
Schéma A (même sujet) : \textbf{Sujet + 不但 + prédicat 1，而且 + prédicat 2。}\\
Schéma B (sujets différents) : \textbf{不但 + Sujet 1 + prédicat 1，而且 + Sujet 2 + (也) + prédicat 2。}
\end{definition}

\begin{methode}[Comment observer une structure grammaticale chinoise]
Face à des exemples nouveaux, procède toujours en quatre étapes :
\begin{enumerate}
\item \textbf{Repérer les mots-outils} : souligne les mots qui se répètent d'un exemple à l'autre (ici \zh{第一}、\zh{第二}、\zh{不但}、\zh{而且}).
\item \textbf{Entourer ce qui change} : le reste de la phrase (sujets, verbes, compléments) est libre ; c'est toi qui le choisiras.
\item \textbf{Comparer les positions du sujet} : le mot-outil est-il avant ou après le sujet ? La position change-t-elle selon les exemples ? (Compare les exemples 2 et 3.)
\item \textbf{Formuler le schéma} : écris la structure avec des cases vides, par exemple « Sujet + 不但 + \ldots ，而且 + \ldots », puis teste-la avec tes propres phrases.
\end{enumerate}
\end{methode}

\begin{attention}[Ne traduis pas « premièrement » mot à mot]
En français, on dit « premièrement » ou « d'abord ». En chinois, l'outil de base de l'énumération, à l'oral comme à l'écrit, est \cle{第一} (dì-yī), littéralement « numéro un ». Le mot \zh{首先} (shǒuxiān, « d'abord, en premier lieu ») existe, mais il s'emploie seul, dans un style plus soutenu, et il ne forme pas une série numérotée : on ne dit pas \zh{首先……第二}. Pour énumérer tes arguments à l'examen, utilise la paire \zh{第一……第二……}, c'est la structure attendue. Autre piège : n'écris jamais \zh{一第} ni \zh{二第} — le préfixe \zh{第} précède toujours le chiffre.
\end{attention}

\courssub{Comparaison rapide avec le français}

\begin{center}
\begin{tabularx}{\linewidth}{|X|X|X|}
\hline
\rowcolor{popblueL} Français & Chinois & Remarque \\
\hline
premièrement\ldots deuxièmement\ldots & 第一\ldots 第二\ldots & correspondance directe, structure très régulière \\
\hline
non seulement\ldots mais aussi\ldots & 不但\ldots 而且\ldots & les deux mots-outils sont obligatoires en chinois standard \\
\hline
d'abord / ensuite & 首先 / 然后 & autre famille de mots, ne pas mélanger avec 第一、第二 \\
\hline
\end{tabularx}
\end{center}

Bonne nouvelle pour les francophones : contrairement à beaucoup de structures chinoises, ces deux-là se traduisent presque mot à mot. La seule vraie difficulté est la \emph{place} de \zh{不但} par rapport au sujet — c'est précisément ce que ta question guide t'a fait découvrir.

\begin{exoresolu}[Repérer la structure]
Énoncé : dans la phrase \zh{我不但喜欢足球，而且喜欢篮球。}(Wǒ bùdàn xǐhuan zúqiú, érqiě xǐhuan lánqiú.), indique : (a) les mots-outils ; (b) le nombre de sujets ; (c) la place de \zh{不但} ; (d) la traduction.
\tcblower
Solution détaillée : (a) les mots-outils sont \zh{不但} et \zh{而且}. (b) Il n'y a qu'un seul sujet, \zh{我}, qui vaut pour les deux propositions. (c) Comme le sujet est unique, \zh{不但} se place \emph{après} le sujet \zh{我} : c'est le schéma A. (d) Traduction : « Non seulement j'aime le football, mais j'aime aussi le basket-ball. » Vérification : si l'on voulait deux sujets (« moi » et « mon frère »), il faudrait dire \zh{不但我喜欢足球，而且我哥哥也喜欢。}
\end{exoresolu}

\begin{aretenir}[Ce qu'il faut retenir de cette observation]
\begin{itemize}
\item \zh{第一……第二……} énumère des arguments : chaque mot-outil ouvre sa proposition, suivi d'une virgule.
\item \zh{不但……而且……} exprime la gradation : « non seulement\ldots mais aussi\ldots ».
\item Règle d'or découverte aujourd'hui : \textbf{un seul sujet} $\Rightarrow$ \zh{不但} \emph{après} le sujet ; \textbf{deux sujets} $\Rightarrow$ \zh{不但} \emph{avant} le premier sujet.
\item Ces deux structures rendent ton expression écrite et orale plus organisée : elles sont très appréciées à l'examen.
\end{itemize}
\end{aretenir}

Dans la prochaine section, nous étudierons en détail la structure \zh{第一……第二……} : sa formation avec le préfixe \zh{第}, ses prolongements (\zh{第三}、\zh{第四}) et ses emplois dans les exposés.

\coursec{Structure 1 : 第一……第二…… (énumération)}

Dans la section précédente, nous avons observé plusieurs phrases authentiques où le locuteur organisait ses idées en les numérotant : \zh{第一}… \zh{第二}… Nous avons remarqué que cette structure permet de présenter des raisons ou des arguments de façon claire et ordonnée, exactement comme le français « premièrement… deuxièmement… troisièmement… ». Nous allons maintenant étudier cette structure de manière rigoureuse : sa formation, son emploi, sa ponctuation, puis ses pièges pour les francophones.

\courssub{Observation : un exemple de la vie scolaire}

Reprenons d'abord un exemple concret, celui d'un élève camerounais qui explique pourquoi il a choisi sa filière.

\begin{exemplebox}[Observer avant de comprendre]
\zh{我选择这个专业，第一，我喜欢它；第二，它很有用；第三，老师很好。}\\
\emph{Wǒ xuǎnzé zhège zhuānyè, dì-yī, wǒ xǐhuan tā; dì-èr, tā hěn yǒuyòng; dì-sān, lǎoshī hěn hǎo.}\\
« J'ai choisi cette filière : premièrement, je l'aime ; deuxièmement, elle est utile ; troisièmement, le professeur est bon. »
\end{exemplebox}

Observons attentivement cette phrase. Que remarque-t-on ?

\begin{itemize}
\item Chaque argument est introduit par \zh{第} suivi d'un nombre : \zh{第一}, \zh{第二}, \zh{第三}.
\item Après \zh{第一} et \zh{第二}, il y a une virgule chinoise \zh{，}.
\item Entre les propositions, on trouve un point-virgule chinois \zh{；}, et non un point.
\item L'ordre des arguments est logique : on ne peut pas dire \zh{第二} avant \zh{第一}.
\end{itemize}

C'est exactement le mécanisme que nous allons formaliser.

\courssub{Formation : 第 + nombre = ordinal}

\begin{propriete}[La marque de l'ordinal : 第]
Le préfixe \cle{第} (dì) placé \textbf{devant un nombre} transforme ce nombre cardinal en \textbf{nombre ordinal} :
\begin{center}
\zh{第} + 一 (yī, un) = \zh{第一} (dì-yī) « premier / premièrement »\\
\zh{第} + 二 (èr, deux) = \zh{第二} (dì-èr) « deuxième / deuxièmement »\\
\zh{第} + 三 (sān, trois) = \zh{第三} (dì-sān) « troisième / troisièmement »
\end{center}
Ce mécanisme est régulier et sans exception : \zh{第四} (quatrième), \zh{第五} (cinquième), \zh{第十} (dixième), etc. On le retrouve partout : \zh{第一课} (leçon 1), \zh{第二名} (deuxième au classement), \zh{三楼} se dit aussi \zh{第三层} (troisième étage).
\end{propriete}

\begin{attention}[二 ou 两 ? Dans l'énumération, toujours 第二 !]
Les francophones connaissent deux mots pour « deux » : \zh{二} (èr) et \zh{两} (liǎng). Règle absolue : après \zh{第}, on utilise \textbf{toujours} \zh{二}, jamais \zh{两}.\\
Correct : \zh{第二} (dì-èr) « deuxième ».\\
Incorrect : \zh{第两} — cette forme n'existe pas.\\
De même : \zh{第二个学生} (dì-èr ge xuésheng) « le deuxième élève », jamais \zh{第两个学生}.
\end{attention}

\courssub{Schéma général de la structure}

\begin{propriete}[Schéma de l'énumération avec 第一……第二……]
\textbf{第一，+ proposition 1 ；第二，+ proposition 2 ；（第三，+ proposition 3。）}
\begin{itemize}
\item \zh{第一} et \zh{第二} sont placés \textbf{en tête de proposition}, suivis d'une virgule \zh{，}.
\item Les propositions sont séparées par un \textbf{point-virgule chinois} \zh{；}.
\item La dernière proposition se termine par un point \zh{。}.
\item On peut prolonger avec \zh{第三}, \zh{第四}… et conclure avec \zh{最后} (zuìhòu, « enfin, pour finir »).
\end{itemize}
\end{propriete}

\begin{popfigure}
\begin{tikzpicture}[
  block/.style={rectangle, rounded corners, draw=popblue, fill=popblueL, align=center, minimum height=0.9cm, inner sep=5pt},
  lab/.style={font=\small\itshape, text=popmuted}
]
\node[block, align=center] (a) at (0,0){第 + 一/二/三\\ \emph{dì + nombre}};
\node[block, align=center] (b) at (3.6,0){，\\ virgule};
\node[block] (c) at (6.4,0) {proposition 1};
\node[block, align=center] (d) at (9.2,0){；\\ point-virgule};
\node[block] (e) at (12.2,0) {第二，proposition 2};
\draw[-{Stealth}, thick, poporange] (a) -- (b);
\draw[-{Stealth}, thick, poporange] (b) -- (c);
\draw[-{Stealth}, thick, poporange] (c) -- (d);
\draw[-{Stealth}, thick, poporange] (d) -- (e);
\draw[-{Stealth}, very thick, popgreen] (0,-1.1) -- (12.2,-1.1) node[midway, below, lab]{ordre logique : du premier argument au dernier};
\end{tikzpicture}
\legende{Schéma de la structure 第一……第二…… : chaque bloc suit un ordre logique fixe.}
\end{popfigure}

\begin{popfigure}
\begin{tikzpicture}[mindmap, grow cyclic, every node/.style={concept}, concept color=popblueL, text=popdark,
  level 1/.append style={level distance=3.2cm, sibling angle=90},
  level 2/.append style={level distance=2.2cm}]
\node[align=center]{第 \emph{dì}\\ (préfixe ordinal)}
  child[concept color=poporangeL] { node[align=center]{第一\\ \emph{dì-yī}\\ 1er / premièrement} }
  child[concept color=popgreenL] { node[align=center]{第二\\ \emph{dì-èr}\\ 2e / deuxièmement} }
  child[concept color=poppurpleL] { node[align=center]{第三\\ \emph{dì-sān}\\ 3e / troisièmement} }
  child[concept color=popgoldL] { node[align=center]{最后\\ \emph{zuìhòu}\\ enfin (clôture)} };
\end{tikzpicture}
\legende{Carte mentale : le préfixe 第 forme les ordinaux ; 最后 clôt l'énumération.}
\end{popfigure}

\courssub{Emplois : énumérer raisons, arguments, étapes, qualités}

La structure \zh{第一……第二……} sert à organiser un discours ou un texte en points numérotés. Elle est très fréquente :

\begin{itemize}
\item pour donner des \textbf{raisons} : pourquoi j'apprends le chinois, pourquoi je suis en retard ;
\item pour présenter des \textbf{arguments} dans un débat ou un exposé ;
\item pour décrire des \textbf{étapes} : les étapes d'une recette, d'un trajet ;
\item pour énumérer des \textbf{qualités} ou des défauts d'une personne, d'un lieu.
\end{itemize}

Elle s'emploie \textbf{à l'oral comme à l'écrit}. À l'écrit, elle donne au texte une organisation claire, très appréciée à l'examen.

\begin{exemplebox}[Exemples traduits et variés]
1. \zh{第一，我家离学校很远；第二，我没有自行车。}\\
\emph{Dì-yī, wǒ jiā lí xuéxiào hěn yuǎn; dì-èr, wǒ méiyǒu zìxíngchē.}\\
« Premièrement, ma maison est loin de l'école ; deuxièmement, je n'ai pas de vélo. »\\[4pt]
2. \zh{我想学汉语，第一，汉语很有意思；第二，中国很大；第三，我想去中国工作。}\\
\emph{Wǒ xiǎng xué Hànyǔ, dì-yī, Hànyǔ hěn yǒu yìsi; dì-èr, Zhōngguó hěn dà; dì-sān, wǒ xiǎng qù Zhōngguó gōngzuò.}\\
« Je veux apprendre le chinois : premièrement, le chinois est intéressant ; deuxièmement, la Chine est un grand pays ; troisièmement, je veux aller travailler en Chine. »\\[4pt]
3. \zh{做这道菜，第一，先洗西红柿；第二，切洋葱；最后，加一点盐。}\\
\emph{Zuò zhè dào cài, dì-yī, xiān xǐ xīhóngshì; dì-èr, qiē yángcōng; zuìhòu, jiā yìdiǎn yán.}\\
« Pour faire ce plat : premièrement, laver d'abord les tomates ; deuxièmement, couper les oignons ; enfin, ajouter un peu de sel. »\\[4pt]
4. \zh{雅温得是一个好城市：第一，人很友好；第二，天气不太热；最后，学校很多。}\\
\emph{Yǎwēndé shì yí ge hǎo chéngshì: dì-yī, rén hěn yǒuhǎo; dì-èr, tiānqì bú tài rè; zuìhòu, xuéxiào hěn duō.}\\
« Yaoundé est une bonne ville : premièrement, les gens sont amicaux ; deuxièmement, il ne fait pas trop chaud ; enfin, il y a beaucoup d'écoles. »
\end{exemplebox}

\courssub{Tableau récapitulatif et variante soutenue}

\begin{center}
\begin{tabularx}{\linewidth}{|X|X|X|X|}
\hline
\rowcolor{popblueL} Caractères & Pinyin & Nature & Traduction \\
\hline
第 & dì & préfixe & (marque de l'ordinal) \\
\hline
第一 & dì-yī & ordinal / connecteur & premier ; premièrement \\
\hline
第二 & dì-èr & ordinal / connecteur & deuxième ; deuxièmement \\
\hline
第三 & dì-sān & ordinal / connecteur & troisième ; troisièmement \\
\hline
最后 & zuìhòu & connecteur & enfin, pour finir \\
\hline
首先 & shǒuxiān & connecteur & d'abord, en premier lieu \\
\hline
其次 & qícì & connecteur & ensuite, en second lieu \\
\hline
\end{tabularx}
\end{center}

\begin{aretenir}[Variante plus soutenue : 首先……其次……]
À l'oral comme dans un texte soigné, on peut remplacer \zh{第一……第二……} par \zh{首先……其次……（最后……）}, qui correspond au français « d'abord… ensuite… enfin… ». Les deux structures ont le même sens ; \zh{第一……第二……} est la plus courante au niveau HSK 1-3.\\
Exemple : \zh{首先，我要感谢老师；其次，我要感谢我的父母。}\\
\emph{Shǒuxiān, wǒ yào gǎnxiè lǎoshī; qícì, wǒ yào gǎnxiè wǒ de fùmǔ.}\\
« D'abord, je veux remercier le professeur ; ensuite, je veux remercier mes parents. »
\end{aretenir}

\courssub{Comparaison avec le français et pièges des francophones}

\begin{center}
\begin{tabularx}{\linewidth}{|X|X|X|}
\hline
\rowcolor{popblueL} Point de comparaison & Chinois & Français \\
\hline
Marque de l'ordinal & 第 + nombre (第一) & suffixe (-ième) ou « premier » \\
\hline
Place du connecteur & en tête de proposition, + virgule & en tête de phrase, + virgule \\
\hline
Séparation des arguments & point-virgule ； & point-virgule ou point \\
\hline
Clôture & 最后 (zuìhòu) & « enfin, pour finir » \\
\hline
\end{tabularx}
\end{center}

\begin{attention}[Erreurs fréquentes des francophones]
\begin{enumerate}
\item \textbf{Oublier la virgule} après \zh{第一} : on écrit \zh{第一，…} et non \zh{第一…}.
\item \textbf{Mettre un point} entre les propositions au lieu du point-virgule \zh{；} : l'énumération forme une seule phrase en chinois.
\item \textbf{Dire 第两} au lieu de \zh{第二} : après \zh{第}, toujours \zh{二}.
\item \textbf{Placer 第一 après le sujet} : le connecteur se met \textbf{avant} la proposition entière, pas au milieu. On dit \zh{第一，我喜欢它}, jamais \zh{我第一喜欢它} (qui signifierait « je l'aime en premier »).
\item \textbf{Mélanger les systèmes} : ne pas alterner \zh{第一} puis \zh{其次} dans la même énumération ; rester cohérent.
\end{enumerate}
\end{attention}

\begin{exoresolu}[Transformation : construire une énumération]
Énoncé. Un élève donne deux raisons pour lesquelles il aime sa ville, Douala : (a) il y a beaucoup de magasins ; (b) le port est très grand. Construisez une phrase complète avec \zh{第一……第二……}.\\
Vocabulaire : 商店 (shāngdiàn, magasin), 港口 (gǎngkǒu, port).
\tcblower
Solution détaillée.\\
Étape 1 : identifier le thème général : \zh{我喜欢杜阿拉。} (Wǒ xǐhuan Dù'ālā.) « J'aime Douala. »\\
Étape 2 : placer \zh{第一，} en tête du premier argument : \zh{第一，商店很多} (dì-yī, shāngdiàn hěn duō) « premièrement, il y a beaucoup de magasins ».\\
Étape 3 : séparer par le point-virgule \zh{；} et introduire le second argument par \zh{第二，} : \zh{第二，港口很大} (dì-èr, gǎngkǒu hěn dà) « deuxièmement, le port est très grand ».\\
Étape 4 : terminer par le point \zh{。}.\\
Phrase finale : \zh{我喜欢杜阿拉：第一，商店很多；第二，港口很大。}\\
\emph{Wǒ xǐhuan Dù'ālā: dì-yī, shāngdiàn hěn duō; dì-èr, gǎngkǒu hěn dà.}\\
« J'aime Douala : premièrement, il y a beaucoup de magasins ; deuxièmement, le port est très grand. »
\end{exoresolu}

\begin{savaistu}[第 dans la vie scolaire chinoise… et camerounaise]
Le préfixe \zh{第} est partout dans la vie de classe : \zh{第二课} (dì-èr kè) signifie « leçon 2 » — c'est justement la leçon que tu étudies ! — et \zh{第一名} (dì-yī míng) désigne le « premier » au classement, une fierté immense pour les familles chinoises. En Chine comme au Cameroun, les résultats scolaires sont affichés avec un rang : être \zh{第一名} à Yaoundé ou à Pékin, c'est le même honneur. On dit aussi \zh{第一次} (dì-yī cì) « la première fois » : \zh{这是我第一次学汉语。} (Zhè shì wǒ dì-yī cì xué Hànyǔ.) « C'est la première fois que j'apprends le chinois. »
\end{savaistu}

\begin{exercice}[À toi de jouer]
1. Traduis en chinois avec \zh{第一……第二……} : « J'aime le football : premièrement, c'est un sport d'équipe ; deuxièmement, les Camerounais jouent très bien. » (équipe : 队 duì ; sport : 运动 yùndòng).\\
2. Corrige la phrase fautive : \zh{我第两喜欢数学。}\\
3. Complète avec \zh{第一}, \zh{第二} ou \zh{最后} : \zh{……，我们先吃饭；……，我们做练习；……，我们回家。}
\end{exercice}

\begin{corrige}[Exercice « À toi de jouer »]
1. \zh{我喜欢足球：第一，足球是团队运动；第二，喀麦隆人踢得很好。} \emph{Wǒ xǐhuan zúqiú: dì-yī, zúqiú shì tuánduì yùndòng; dì-èr, Kāmàilóng rén tī de hěn hǎo.}\\
2. La forme correcte est \zh{第二} : \zh{我第二喜欢数学。} signifie « les mathématiques sont ma deuxième préférence » ; si l'on voulait énumérer, il faudrait \zh{第二，我喜欢数学。}\\
3. \zh{第一，我们先吃饭；第二，我们做练习；最后，我们回家。} « Premièrement, nous mangeons d'abord ; deuxièmement, nous faisons les exercices ; enfin, nous rentrons à la maison. »
\end{corrige}

Nous savons désormais énumérer des arguments dans un ordre logique avec \zh{第一……第二……（第三……最后）}. Mais comment faire pour ajouter une information qui \emph{renforce} la précédente, comme le français « non seulement… mais aussi… » ? C'est l'objet de la section suivante : la structure de gradation \zh{不但……而且……}.

\coursec{Structure 2 : 不但……而且…… (gradation)}

Dans la section précédente, nous avons appris à \cle{énumérer} des arguments avec \zh{第一……第二……}, une structure qui met des idées \emph{côte à côte}, au même niveau. Mais dans la vie, toutes les idées n'ont pas le même poids : parfois, on veut dire qu'une chose est vraie, \emph{et qu'en plus} une autre chose, plus forte encore, est vraie aussi. En français, on dit : « non seulement... mais aussi... ». En chinois, c'est exactement le rôle de la structure \zh{不但……而且……}, que nous étudions maintenant.

\courssub{Observation guidée : lire avant de comprendre}

\begin{textebox}[Notre école à Yaoundé]
我们的学校不但很大，而且很干净。\\
Wǒmen de xuéxiào bùdàn hěn dà, érqiě hěn gānjìng.\\
« Notre école est non seulement grande, mais aussi très propre. »\\[4pt]
不但老师喜欢足球，而且学生也喜欢。\\
Bùdàn lǎoshī xǐhuan zúqiú, érqiě xuésheng yě xǐhuan.\\
« Non seulement les professeurs aiment le football, mais les élèves l'aiment aussi. »\\[4pt]
她不但会说汉语，而且会写汉字。\\
Tā bùdàn huì shuō Hànyǔ, érqiě huì xiě Hànzì.\\
« Elle sait non seulement parler chinois, mais aussi écrire les caractères. »
\end{textebox}

\textbf{Questions d'observation :}
\begin{enumerate}
\item Dans la première phrase, qui est le sujet des deux verbes (大 et 干净) ? Où se trouve 不但 par rapport au sujet ?
\item Dans la deuxième phrase, les deux propositions ont-elles le même sujet ? Où se trouve 不但 ? Quel mot nouveau apparaît dans la deuxième proposition ?
\item Que remarquez-vous sur la place de 而且 dans les trois phrases ?
\end{enumerate}

Vous l'avez remarqué : la structure se comporte différemment selon que les deux propositions ont \cle{le même sujet} ou \cle{deux sujets différents}. C'est la clé de toute cette leçon.

\courssub{Le sens : la gradation}

\begin{definition}[La gradation]
La \cle{gradation} est une figure (et une structure logique) qui \textbf{renforce le propos en ajoutant un élément plus important, plus fort ou complémentaire} au premier. La deuxième information « monte d'un cran » : elle surprend, elle ajoute, elle dépasse la première. En français : « non seulement... mais aussi... », « ... et même... ». En chinois : 不但……而且…… (\emph{bùdàn... érqiě...}).
\end{definition}

\begin{exemplebox}[Sentir la gradation]
\zh{他不但通过了考试，而且得了第一名。}\\
\emph{Tā bùdàn tōngguò le kǎoshì, érqiě dé le dì-yī míng.}\\
« Non seulement il a réussi l'examen, mais en plus il a eu la première place. »\\[4pt]
La deuxième information (得第一名) est \emph{plus forte} que la première (通过考试) : c'est cela, la gradation. Si l'on inversait les deux propositions, la phrase perdrait sa logique.
\end{exemplebox}

\courssub{Le schéma général}

\begin{propriete}[Schéma de la structure]
\textbf{不但 + proposition 1，而且 + proposition 2}\\
= « non seulement [proposition 1], mais aussi [proposition 2] ».
\begin{itemize}
\item 不但 (\emph{bùdàn}) introduit la première information ;
\item 而且 (\emph{érqiě}) introduit la deuxième information, plus forte ;
\item les deux mots sont \textbf{inséparables} dans cette structure : 而且 est obligatoire pour « fermer » la phrase ;
\item une virgule chinoise ，sépare les deux propositions.
\end{itemize}
\end{propriete}

\courssub{Cas A — Les deux propositions ont le même sujet}

\begin{propriete}[Cas A : même sujet]
Quand les deux propositions parlent de la \textbf{même personne ou de la même chose}, le sujet se place \textbf{avant} 不但 et n'est \textbf{pas répété} dans la deuxième proposition :\\[4pt]
\textbf{Sujet + 不但 + Verbe/Adjectif，而且 + Verbe/Adjectif}
\end{propriete}

\begin{exemplebox}[Cas A en action]
\zh{她不但聪明，而且很努力。}\\
\emph{Tā bùdàn cōngming, érqiě hěn nǔlì.}\\
« Elle est non seulement intelligente, mais aussi travailleuse. »\\[6pt]
\zh{这所学校不但很大，而且很漂亮。}\\
\emph{Zhè suǒ xuéxiào bùdàn hěn dà, érqiě hěn piàoliang.}\\
« Cette école est non seulement grande, mais aussi belle. »\\[6pt]
\zh{杜阿拉不但人口多，而且经济很发达。}\\
\emph{Dù'ālā bùdàn rénkǒu duō, érqiě jīngjì hěn fādá.}\\
« Douala a non seulement une grande population, mais aussi une économie très développée. »\\[6pt]
\zh{他不但会说汉语，而且会写汉字。}\\
\emph{Tā bùdàn huì shuō Hànyǔ, érqiě huì xiě Hànzì.}\\
« Il sait non seulement parler chinois, mais aussi écrire les caractères. »
\end{exemplebox}

Notez que dans la dernière phrase, le verbe 会 est répété dans la deuxième proposition : c'est normal, car chaque proposition a besoin de son propre verbe. C'est le \emph{sujet} (他) qui n'est pas répété.

\courssub{Cas B — Les deux propositions ont des sujets différents}

\begin{propriete}[Cas B : sujets différents]
Quand les deux propositions ont des \textbf{sujets différents}, 不但 se place \textbf{avant le premier sujet}, et l'adverbe \cle{也} (\emph{yě}, « aussi ») apparaît dans la deuxième proposition, juste après le deuxième sujet :\\[4pt]
\textbf{不但 + Sujet 1 + Verbe/Adj，而且 + Sujet 2 + 也 + Verbe/Adj}
\end{propriete}

\begin{exemplebox}[Cas B en action]
\zh{不但老师喜欢这首歌，而且学生也喜欢。}\\
\emph{Bùdàn lǎoshī xǐhuan zhè shǒu gē, érqiě xuésheng yě xǐhuan.}\\
« Non seulement le professeur aime cette chanson, mais les élèves l'aiment aussi. »\\[6pt]
\zh{不但中国人喝茶，而且喀麦隆人也喝茶。}\\
\emph{Bùdàn Zhōngguórén hē chá, érqiě Kāmàilóngrén yě hē chá.}\\
« Non seulement les Chinois boivent du thé, mais les Camerounais en boivent aussi. »\\[6pt]
\zh{不但哥哥学汉语，而且我也学汉语。}\\
\emph{Bùdàn gēge xué Hànyǔ, érqiě wǒ yě xué Hànyǔ.}\\
« Non seulement mon grand frère apprend le chinois, mais moi aussi je l'apprends. »
\end{exemplebox}

\begin{attention}[也 est obligatoire dans le cas B]
Quand les sujets sont différents, n'oubliez jamais \cle{也} dans la deuxième proposition : il porte à lui seul l'idée de « aussi ». Sans 也, la phrase sonne incomplète, voire fautive. Comparez :\\
\zh{不但老师喜欢这首歌，而且学生也喜欢。} (correct)\\
*不但老师喜欢这首歌，而且学生喜欢。(incorrect : il manque 也)
\end{attention}

\begin{popfigure}
\begin{tikzpicture}[font=\small]
\node[draw=popblue, fill=popblueL, rounded corners, align=center, text width=6.2cm, minimum height=2.6cm] (A) at (0,0)
{\textbf{Cas A : même sujet}\\[2pt]
Sujet + \textcolor{popink}{\textbf{不但}} + X，\\
\textcolor{popink}{\textbf{而且}} + Y\\[2pt]
她不但聪明，而且很努力。};
\node[draw=poporange, fill=poporangeL, rounded corners, align=center, text width=6.2cm, minimum height=2.6cm] (B) at (8.2,0)
{\textbf{Cas B : sujets différents}\\[2pt]
\textcolor{popink}{\textbf{不但}} + Sujet 1 + X，\\
\textcolor{popink}{\textbf{而且}} + Sujet 2 + \textcolor{popink}{\textbf{也}} + Y\\[2pt]
不但老师喜欢，而且学生也喜欢。};
\node[align=center, font=\bfseries] at (4.1,2.2) {不但……而且……};
\end{tikzpicture}
\legende{Les deux constructions de la gradation : la place de 不但 dépend du sujet.}
\end{popfigure}

\courssub{Comment choisir : l'organigramme de décision}

\begin{popfigure}
\begin{tikzpicture}[font=\small, node distance=0.9cm]
\node[draw=popdark, fill=popcream, rounded corners, align=center, text width=4.5cm] (start) {Je veux dire « non seulement... mais aussi... »};
\node[draw=poppurple, fill=poppurpleL, diamond, aspect=2.4, align=center, below=of start, text width=3.2cm, inner sep=1pt] (q) {Les deux propositions ont-elles le même sujet ?};
\node[draw=popblue, fill=popblueL, rounded corners, align=center, below left=1.1cm and -0.6cm of q, text width=5.2cm] (oui) {\textbf{Oui} : le sujet va \textbf{avant} 不但\\ Sujet + 不但 + X，而且 + Y};
\node[draw=poporange, fill=poporangeL, rounded corners, align=center, below right=1.1cm and -0.6cm of q, text width=5.2cm] (non) {\textbf{Non} : 不但 va \textbf{avant} le sujet 1\\ 不但 + S1 + X，而且 + S2 + 也 + Y};
\draw[-{Stealth}, thick, popdark] (start) -- (q);
\draw[-{Stealth}, thick, popblue] (q) -| (oui) node[midway, above left, font=\bfseries\small] {oui};
\draw[-{Stealth}, thick, poporange] (q) -| (non) node[midway, above right, font=\bfseries\small] {non};
\end{tikzpicture}
\legende{Organigramme de décision : une seule question à se poser — « même sujet ? ».}
\end{popfigure}

\begin{methode}[Construire une phrase avec 不但……而且……]
\begin{enumerate}
\item Identifiez les deux informations à relier et vérifiez que la deuxième est bien \emph{plus forte} ou complémentaire.
\item Posez-vous la question : les deux propositions ont-elles le même sujet ?
\item Si oui : placez le sujet une seule fois, avant 不但, puis 不但 + X，而且 + Y.
\item Si non : placez 不但 avant le premier sujet, puis 而且 + deuxième sujet + 也 + Y.
\item Relisez : 而且 est-il présent ? 也 est-il présent (cas B) ? La virgule ，est-elle placée ?
\end{enumerate}
\end{methode}

\courssub{Variantes et registre soutenu}

\begin{propriete}[Variantes de la structure]
\begin{itemize}
\item 而且 peut être \textbf{renforcé par 还} (\emph{hái}) ou \textbf{也} (\emph{yě}) dans la deuxième proposition, même sujet compris : \zh{他不但聪明，还很努力。} (\emph{Tā bùdàn cōngming, hái hěn nǔlì.}) « Il est non seulement intelligent, mais encore très travailleur. »
\item Au registre \textbf{soutenu et écrit}, on remplace 不但 par \cle{不仅} (\emph{bùjǐn}) : 不仅……而且……, très fréquent dans les textes d'examen et les articles.
\end{itemize}
\end{propriete}

\begin{exemplebox}[Registre soutenu avec 不仅]
\zh{汉语不仅有用，而且很有意思。}\\
\emph{Hànyǔ bùjǐn yǒuyòng, érqiě hěn yǒu yìsi.}\\
« Le chinois est non seulement utile, mais aussi très intéressant. »\\[6pt]
\zh{中国不仅是一个大国，而且是一个文明古国。}\\
\emph{Zhōngguó bùjǐn shì yí ge dàguó, érqiě shì yí ge wénmíng gǔguó.}\\
« La Chine est non seulement un grand pays, mais aussi un pays de civilisation ancienne. »
\end{exemplebox}

\begin{savaistu}[不仅, la variante des textes d'examen]
Dans la presse chinoise, les discours officiels et les textes littéraires, c'est presque toujours 不仅……而且…… que l'on rencontre, et non 不但. Les deux structures ont exactement le même sens et les mêmes règles de place (même sujet / sujets différents). À l'examen, si un texte de compréhension contient 不仅, ne soyez pas surpris : lisez-le comme 不但. À l'écrit, utiliser 不仅 à la place de 不但 donne immédiatement à votre production un ton plus soigné — les correcteurs apprécient !
\end{savaistu}

\courssub{Comparaison avec le français et tableau récapitulatif}

En français, « non seulement » se place de façon assez libre et « mais aussi » peut être remplacé par « mais également », « mais même ». En chinois, la structure est beaucoup plus \textbf{mécanique} : la place de 不但 dépend uniquement de l'identité des sujets, et 而且 (avec 也 si nécessaire) est incontournable. C'est une bonne nouvelle : une fois la règle apprise, il n'y a presque pas d'exception.

\begin{center}
\begin{tabularx}{\linewidth}{|X|X|X|}
\hline
\rowcolor{popblueL} \textbf{Situation} & \textbf{Structure chinoise} & \textbf{Français} \\
\hline
Même sujet & Sujet + 不但 + X，而且 + Y & Sujet + non seulement X, mais aussi Y \\
\hline
Sujets différents & 不但 + S1 + X，而且 + S2 + 也 + Y & Non seulement S1 + X, mais S2 + Y aussi \\
\hline
Registre soutenu & 不仅……而且…… & non seulement... mais aussi... (soutenu) \\
\hline
Renforcement & ……，而且 还/也 + Y & ... et même / et aussi \\
\hline
\end{tabularx}
\end{center}

\begin{center}
\begin{tabularx}{\linewidth}{|X|X|X|X|}
\hline
\rowcolor{popblueL} Caractères & Pinyin & Nature & Traduction \\
\hline
不但 & bùdàn & conjonction & non seulement \\
\hline
而且 & érqiě & conjonction & mais aussi, et de plus \\
\hline
不仅 & bùjǐn & conjonction (soutenue) & non seulement \\
\hline
也 & yě & adverbe & aussi \\
\hline
还 & hái & adverbe & encore, en plus \\
\hline
聪明 & cōngming & adjectif & intelligent \\
\hline
努力 & nǔlì & adjectif / verbe & travailleur, faire des efforts \\
\hline
发达 & fādá & adjectif & développé \\
\hline
\end{tabularx}
\end{center}

\courssub{Exercice résolu et pièges à éviter}

\begin{exoresolu}[Transformer et construire]
\textbf{Énoncé.} 1) Combinez avec 不但……而且…… : 他会说法语。他会说汉语。(même sujet) \quad 2) Combinez : 我喜欢足球。我弟弟喜欢足球。(sujets différents) \quad 3) Traduisez : « Non seulement Yaoundé est grande, mais elle est aussi très belle. »
\tcblower
\textbf{Solution détaillée.}\\
1) Même sujet (他) : le sujet passe devant 不但 et n'est pas répété.\\
\zh{他不但会说法语，而且会说汉语。} \emph{Tā bùdàn huì shuō Fǎyǔ, érqiě huì shuō Hànyǔ.} « Il sait parler non seulement français, mais aussi chinois. »\\[4pt]
2) Sujets différents (我 / 我弟弟) : 不但 avant le premier sujet, et 也 dans la deuxième proposition.\\
\zh{不但我喜欢足球，而且我弟弟也喜欢。} \emph{Bùdàn wǒ xǐhuan zúqiú, érqiě wǒ dìdi yě xǐhuan.} « Non seulement j'aime le football, mais mon petit frère l'aime aussi. »\\[4pt]
3) Même sujet (Yaoundé) : \zh{雅温得不但很大，而且很漂亮。} \emph{Yǎwēndé bùdàn hěn dà, érqiě hěn piàoliang.}
\end{exoresolu}

\begin{attention}[Les trois erreurs classiques des francophones]
\begin{enumerate}
\item \textbf{Oublier 而且.} Sous l'influence du français, où « mais aussi » peut être sous-entendu, beaucoup d'élèves ferment la phrase trop tôt : *\zh{他不但会说汉语，会写汉字。} est \textbf{faux}. 而且 est \textbf{obligatoire} pour fermer la structure : \zh{他不但会说汉语，而且会写汉字。}
\item \textbf{Placer 不但 avant le sujet quand il n'y a qu'un seul sujet} : *\zh{不但她聪明，而且很努力。} est faux ; dites \zh{她不但聪明，而且很努力。}
\item \textbf{Utiliser 而且 seul} en début de phrase pour dire « mais aussi » : 而且 ne s'emploie jamais seul dans ce sens ; il \textbf{répond toujours} à 不但 ou 不仅 dans cette structure. Pour ajouter une idée dans une nouvelle phrase, le français « de plus / en outre » se dit plutôt 另外 (\emph{lìngwài}) ou 而且 uniquement si un 不但 a été posé avant.
\end{enumerate}
\end{attention}

\courssub{Exercices d'application}

\begin{exercice}[Entraînement : Construire et traduire]
\begin{enumerate}
\item \textbf{Associez les deux phrases avec 不但……而且…… (même sujet).}\\
a) 这道菜很好吃。这道菜不贵。\\
b) 我的手机坏了。我的电脑也坏了。 (Attention : sujets différents !)\\
c) 小王会开车。小王会修车。\\
d) 今天很热。今天很湿 (shī, humide).
\item \textbf{Traduisez en chinois (caractères + pinyin).}\\
a) Non seulement le professeur est sévère, mais il est aussi très gentil.\\
b) Non seulement les élèves de Première, mais aussi les élèves de Terminale aiment ce film.\\
c) Le chinois n'est pas seulement difficile, il est aussi très utile. (Utilisez 不仅).
\item \textbf{Transformez les phrases suivantes en utilisant la structure demandée.}\\
a) 我喜欢唱歌。我喜欢跳舞。 (不但……而且……)\\
b) 他不但高，而且壮。 (改用 不仅……而且……)\\
c) 爸爸会做饭。妈妈也会做饭。 (不但……而且……)
\end{enumerate}
\end{exercice}

\begin{corrige}[Exercices d'application]
\begin{enumerate}
\item a) \zh{这道菜不但很好吃，而且不贵。} (\emph{Zhè dào cài bùdàn hěn hǎochī, érqiě bù guì.})\\
b) \zh{不但我的手机坏了，而且我的电脑也坏了。} (\emph{Bùdàn wǒ de shǒujī huài le, érqiě wǒ de diànnǎo yě huài le.})\\
c) \zh{小王不但会开车，而且会修车。} (\emph{Xiǎo Wáng bùdàn huì kāichē, érqiě huì xiūchē.})\\
d) \zh{今天不但很热，而且很湿。} (\emph{Jīntiān bùdàn hěn rè, érqiě hěn shī.})
\item a) \zh{老师不但严厉，而且很和蔼。} (\emph{Lǎoshī bùdàn yánlì, érqiě hěn hé'ǎi.}) — \emph{Note : même sujet (老师), pas de 也.}\\
b) \zh{不但高一学生喜欢这部电影，而且高三学生也喜欢。} (\emph{Bùdàn Gāoyī xuésheng xǐhuan zhè bù diànyǐng, érqiě Gāosān xuésheng yě xǐhuan.}) — \emph{Sujets différents, 也 obligatoire.}\\
c) \zh{汉语不仅不难，而且很有用。} (\emph{Hànyǔ bùjǐn bù nán, érqiě hěn yǒuyòng.}) — \emph{Registre soutenu avec 不仅.}
\item a) \zh{我不但喜欢唱歌，而且喜欢跳舞。} (\emph{Wǒ bùdàn xǐhuan chànggē, érqiě xǐhuan tiàowǔ.})\\
b) \zh{他不仅高，而且壮。} (\emph{Tā bùjǐn gāo, érqiě zhuàng.})\\
c) \zh{不但爸爸会做饭，而且妈妈也会做饭。} (\emph{Bùdàn bàba huì zuòfàn, érqiě māma yě huì zuòfàn.})
\end{enumerate}
\end{corrige}

\begin{aretenir}[L'essentiel de la structure 2]
\begin{itemize}
\item 不但……而且…… = « non seulement... mais aussi... » : la deuxième information est \textbf{plus forte}.
\item \textbf{Même sujet} : Sujet + 不但 + X，而且 + Y (le sujet précède 不但).
\item \textbf{Sujets différents} : 不但 + S1 + X，而且 + S2 + \textbf{也} + Y (不但 précède le sujet, 也 apparaît).
\item 而且 est \textbf{toujours obligatoire} ; il ne s'emploie jamais seul pour « mais aussi ».
\item Registre soutenu : \textbf{那些}……那些…… ; renforcement possible avec 还 ou 也.
\item \textbf{Pour l'examen} : savoir passer du Cas A au Cas B, ne pas oublier 也 (Cas B) ni 那些 (les deux cas), et reconnaître 不仅 comme variante de 不但.
\end{itemize}
\end{aretenir}

\coursec{Leçon 2 — Grammaire : 第一……第二…… ; 不但……而且……}

\courssub{Observation guidée d'exemples}

\begin{textebox}[Situation de communication : Présenter le Cameroun à un ami chinois]
\zh{王朋：喀麦隆怎么样？}\\
\zh{恩波贝尔：喀麦隆很好。第一，风景很美；第二，菜很好吃，比如恩东菜；第三，人很热情。最后，足球很厉害。}\\
\zh{王朋：听起来很棒。不但我想去旅游，而且我想学法语。}\\
\zh{恩波贝尔：太好了！不但你会说法语，而且你弟弟也会说。我们可以一起交流。}\\[4pt]
\emph{Wáng Péng: Kāmàilóng zěnmeyàng?}\\
\emph{Ēnbōbèi'ěr: Kāmàilóng hěn hǎo. Dìyī, fēngjǐng hěn měi; dìèr, cài hěn hǎochī, bǐrú Ēndōngcài; dìsān, rén hěn rèqíng. Zuìhòu, zúqiú hěn lìhai.}\\
\emph{Wáng Péng: Tīng qǐlái hěn bàng. Bùdàn wǒ xiǎng qù lǚyóu, érqiě wǒ xiǎng xué Fǎyǔ.}\\
\emph{Ēnbōbèi'ěr: Tài hǎo le! Bùdàn nǐ huì shuō Fǎyǔ, érqiě nǐ dìdi yě huì shuō. Wǒmen kěyǐ yìqǐ jiāoliú.}\\[4pt]
« Wang Peng : Comment est le Cameroun ? \\
Ngo Bell : Le Cameroun est super. Premièrement, les paysages sont beaux ; deuxièmement, la cuisine est délicieuse, par exemple le ndolé ; troisièmement, les gens sont chaleureux. Enfin, le football est fort. \\
Wang Peng : Ça a l'air super. Non seulement je veux voyager, mais je veux aussi apprendre le français. \\
Ngo Bell : Super ! Non seulement tu parles français, mais ton petit frère aussi. On peut échanger ensemble. »
\end{textebox}

\begin{methode}[Comment observer ?]
\begin{enumerate}
\item Repère les deux structures cibles : \zh{第一……第二……第三……最后} et \zh{不但……而且……}.
\item Identifie le sujet de chaque phrase : est-il unique ou y en a-t-il deux différents ?
\item Note la place du sujet par rapport à \zh{不但} et la présence de \zh{也} dans la seconde partie.
\item Observe la ponctuation chinoise (， ；) qui structure l'énumération.
\end{enumerate}
\end{methode}

\courssub{Structure 1 : 第一……第二…… (Énumération ordonnée)}

\begin{propriete}[Schéma et emplois de l'énumération]
\begin{enumerate}
\item \textbf{Schéma de base} : \zh{第一，[Phrase 1]；第二，[Phrase 2]；第三，[Phrase 3]……最后，[Phrase finale]。}
\item \textbf{Emploi} : Énumérer des arguments, des raisons, des étapes ou des points de vue de façon logique et formelle. Chaque item est une phrase complète.
\item \textbf{Marqueur de fin} : \zh{最后} (\emph{zuìhòu}) « enfin, finalement » remplace \zh{第} + nombre pour clore la liste.
\item \textbf{Variante chronologique} : Pour des actions successives dans le temps, on préfère \zh{先……然后……} (\emph{xiān... ránhòu...}) « d'abord... ensuite... ».
\end{enumerate}
\end{propriete}

\begin{exemplebox}[Exemples d'énumération (contexte camerounais)]
\zh{学习汉字有三个步骤：第一，看笔顺；第二，练书写；第三，背意思。最后，每天复习。}\\
\emph{Xuéxí Hànzì yǒu sān ge bùzhòu: dìyī, kàn bǐshùn; dìèr, liàn shūxiě; dìsān, bèi yìsi. Zuìhòu, měitiān fùxí.}\\
« L'apprentissage des caractères a trois étapes : premièrement, observer l'ordre des traits ; deuxièmement, s'entraîner à l'écriture ; troisièmement, mémoriser le sens. Enfin, réviser chaque jour. »\\[4pt]
\zh{我喜欢雅温得，第一，气候舒服；第二，交通方便。}\\
\emph{Wǒ xǐhuan Yǎwēndé, dìyī, qìhòu shūfu; dìèr, jiāotōng fāngbiàn.}\\
« J'aime Yaoundé : premièrement, le climat est agréable ; deuxièmement, les transports sont pratiques. »
\end{exemplebox}

\begin{center}
\begin{tabularx}{\linewidth}{|X|X|X|X|}
\hline
\rowcolor{popblueL} Caractères & Pinyin & Nature & Traduction \\
\hline
第一 & dìyī & nombre ordinal & premièrement, en premier \\
\hline
第二 & dìèr & nombre ordinal & deuxièmement, en second \\
\hline
第三 & dìsān & nombre ordinal & troisièmement, en troisième \\
\hline
最后 & zuìhòu & adv. & enfin, finalement, en dernier \\
\hline
先……然后…… & xiān... ránhòu... & loc. & d'abord... ensuite... (chronologie) \\
\hline
\end{tabularx}
\end{center}

\begin{attention}[Piège majeur : \textbf{第两} n'existe pas]
Le chiffre « deux » s'écrit \zh{二} (\emph{èr}) dans les nombres ordinaux (rang) et \zh{两} (\emph{liǎng}) devant un classificateur (quantité).\\
\zh{第二} \emph{dìèr} « deuxième / second » \textbf{✓} \hspace{1cm} \zh{*第两} \textbf{✗}\\
\zh{第二个} \emph{dìèr ge} « le deuxième » \textbf{✓} \hspace{1cm} \zh{*第两个} \textbf{✗}\\
\zh{两个人} \emph{liǎng ge rén} « deux personnes » \textbf{✓} (quantité).
\end{attention}

\courssub{Structure 2 : 不但……而且…… (Gradation / Progression)}

\begin{propriete}[La règle d'or : un sujet ou deux sujets ?]
La structure \zh{不但……而且……} (\emph{bùdàn... érqiě...}) exprime une \textbf{gradation} : le second élément renforce le premier (« non seulement... mais en plus... »). La place du sujet est \textbf{mécanique} :
\begin{enumerate}
\item \textbf{Un seul sujet (commun aux deux verbes)} : Le sujet se place \textbf{une seule fois, en tête}, \textbf{devant} \zh{不但}.\\
\zh{Sujet + 不但 + Verbe/Adj 1，而且 + Verbe/Adj 2。}\\
\zh{喀麦隆不但美丽，而且很丰富。} \emph{Kāmàilóng bùdàn měilì, érqiě hěn fēngfù.} « Le Cameroun n'est pas seulement beau, il est aussi riche. »
\item \textbf{Deux sujets différents} : \zh{不但} précède le \textbf{premier sujet} ; le second sujet suit \zh{而且} et \textbf{doit être suivi de \zh{也} (ou \zh{还})}.\\
\zh{不但 + Sujet 1 + Verbe 1，而且 + Sujet 2 + 也/还 + Verbe 2。}\\
\zh{不但老师喜欢足球，而且学生也喜欢。} \emph{Bùdàn lǎoshī xǐhuan zúqiú, érqiě xuésheng yě xǐhuan.} « Non seulement le prof aime le foot, mais les élèves aussi. »
\end{enumerate}
\end{propriete}

\begin{exemplebox}[Cas 1 : Sujet unique (le plus fréquent en HSK 1-3)]
\zh{玛丽不但会说英语，而且会说法语。}\\
\emph{Mǎlì bùdàn huì shuō Yīngyǔ, érqiě huì shuō Fǎyǔ.}\\
« Non seulement Marie parle anglais, mais elle parle aussi français. »\\[4pt]
\zh{这道菜不但好吃，而且不贵。}\\
\emph{Zhè dào cài bùdàn hǎochī, érqiě bù guì.}\\
« Ce plat n'est pas seulement bon, en plus il n'est pas cher. »
\end{exemplebox}

\begin{exemplebox}[Cas 2 : Deux sujets différents (attention à \zh{也})]
\zh{不但我爸爸喜欢恩东菜，而且我妈妈也喜欢。}\\
\emph{Bùdàn wǒ bàba xǐhuan Ēndōngcài, érqiě wǒ māma yě xǐhuan.}\\
« Non seulement mon père aime le ndolé, mais ma mère aussi. »\\[4pt]
\zh{不但中国学生学法语，而且喀麦隆学生也学汉语。}\\
\emph{Bùdàn Zhōngguó xuésheng xué Fǎyǔ, érqiě Kāmàilóng xuésheng yě xué Hànyǔ.}\\
« Non seulement les étudiants chinois apprennent le français, mais les étudiants camerounais apprennent aussi le chinois. »
\end{exemplebox}

\begin{center}
\begin{tabularx}{\linewidth}{|X|X|X|}
\hline
\rowcolor{poppurpleL} Structure & Exemple (Caractères + Pinyin) & Traduction \\
\hline
Sujet unique : S + 不但 + A，而且 + B & 他不但高，而且壮 \emph{Tā bùdàn gāo, érqiě zhuàng} & Non seulement il est grand, mais il est costaud \\
\hline
Deux sujets : 不但 + S1 + A，而且 + S2 + 也 + B & 不但哥哥去，而且弟弟也去 \emph{Bùdàn gēge qù, érqiě dìdi yě qù} & Non seulement le grand frère y va, mais le petit aussi \\
\hline
\end{tabularx}
\end{center}

\begin{attention}[Les 4 pièges grammaticaux des francophones]
\begin{enumerate}
\item \textbf{Oublier \zh{而且}} : \zh{*他不但聪明，很努力。} ✗ $\rightarrow$ \zh{他不但聪明，而且很努力。} ✓ (\emph{Érqiě} est le pivot obligatoire).
\item \textbf{Mal placer \zh{不但} (sujet unique)} : \zh{*不但他聪明，而且很努力。} ✗ $\rightarrow$ \zh{他不但聪明，而且很努力。} ✓ (Sujet unique = en tête).
\item \textbf{Oublier \zh{也} (deux sujets)} : \zh{*不但我喜欢，而且他喜欢。} ✗ $\rightarrow$ \zh{不但我喜欢，而且他也喜欢。} ✓ (\emph{Yě} porte le « aussi »).
\item \textbf{Confondre \zh{也} et la gradation} : \zh{他也聪明，也努力} = « Il est aussi intelligent, aussi travailleur » (simple addition). Pour la gradation : \zh{他不但聪明，而且努力} ✓.
\end{enumerate}
\end{attention}

\courssub{Comparaison avec le français et pièges des francophones}

\begin{exemplebox}[Observation : deux langues, deux logiques]
\textbf{Énumération :}\\
\zh{我喜欢汉语，第一，汉语很有意思；第二，汉语很有用。}\\
\emph{Wǒ xǐhuan Hànyǔ, dìyī, Hànyǔ hěn yǒu yìsi; dìèr, Hànyǔ hěn yǒuyòng.}\\
« J'aime le chinois : premièrement, le chinois est intéressant ; deuxièmement, il est utile. »\\[4pt]
\textbf{Gradation :}\\
\zh{他不但会说汉语，而且会写汉字。}\\
\emph{Tā bùdàn huì shuō Hànyǔ, érqiě huì xiě Hànzì.}\\
« Non seulement il sait parler chinois, mais il sait aussi écrire les caractères. »
\end{exemplebox}

\courssub{« Premièrement, deuxièmement » = 第一、第二 : une correspondance presque parfaite}

Bonne nouvelle : l'énumération chinoise fonctionne presque comme la française. \zh{第一} correspond à « premièrement », \zh{第二} à « deuxièmement », \zh{第三} à « troisièmement ». Mais deux différences doivent être notées.

\begin{propriete}[Correspondances et différences]
\begin{enumerate}
\item Le français utilise un adverbe en \emph{-ment} (« premièrement ») ou une locution avec préposition (« en premier lieu », « d'abord », « ensuite »). Le chinois n'a \textbf{ni suffixe ni préposition} : il suffit du préfixe ordinal \zh{第} + nombre : \zh{第一}、\zh{第二}、\zh{第三}.
\item En chinois, chaque point de l'énumération est une \textbf{phrase complète} introduite directement par \zh{第} + nombre, souvent séparée par une virgule chinoise « ，» ou un point-virgule « ；». On ne dit jamais \zh{*第一地} ni \zh{*在第一}.
\item Pour conclure l'énumération, le français dit « enfin » ou « finalement » ; le chinois dit \zh{最后} (\emph{zuìhòu}, « en dernier »), sans \zh{第}.
\end{enumerate}
\end{propriete}

\begin{center}
\begin{tabularx}{\linewidth}{|X|X|X|}
\hline
\rowcolor{popblueL} Français & Chinois (caractères + pinyin) & Remarque \\
\hline
premièrement & 第一 \emph{dìyī} & préfixe 第 + 一 \\
\hline
deuxièmement & 第二 \emph{dìèr} & jamais 第两 \\
\hline
troisièmement & 第三 \emph{dìsān} & régulier \\
\hline
enfin, finalement & 最后 \emph{zuìhòu} & pas de 第 \\
\hline
d'abord... ensuite... & 先……然后…… \emph{xiān... ránhòu...} & autre système (chronologie) \\
\hline
\end{tabularx}
\end{center}

\courssub{« Non seulement... mais aussi... » : la règle stricte du sujet}

C'est avec la gradation que les francophones trébuchent le plus. En français, la place de « non seulement » est assez libre : « Non seulement il parle chinois, mais il l'écrit » ou « Il parle non seulement le chinois, mais aussi l'anglais ». Le chinois, lui, obéit à une \cle{règle mécanique} fondée sur une seule question : \textbf{y a-t-il un seul sujet ou deux sujets différents ?}

\begin{propriete}[Rappel de la règle d'or de 不但……而且……]
\begin{enumerate}
\item \textbf{Un seul sujet} : le sujet se place \textbf{une seule fois, en tête}, devant \zh{不但}. Schéma : \zh{Sujet + 不但 + A，而且 + B}。
\item \textbf{Deux sujets différents} : \zh{不但} se place \textbf{avant le premier sujet}, et le second sujet est suivi de \zh{也} (ou \zh{还}). Schéma : \zh{不但 + Sujet 1 + A，而且 + Sujet 2 + 也 + B}。
\end{enumerate}
\end{propriete}

\begin{attention}[En français « non seulement il parle, mais il écrit »]
Le français répète le pronom (« il... il... ») et peut même inverser le sujet. Le chinois interdit les deux : avec un seul sujet, celui-ci apparaît \textbf{une seule fois, au tout début} :\\
\zh{他不但会说，而且会写。} ✓ \emph{Tā bùdàn huì shuō, érqiě huì xiě.}\\
\zh{*不但他会说，而且他会写。} ✗ (sujet répété et mal placé)\\
Retenez : en chinois, on ne « dédouble » jamais le sujet autour de \zh{不但}.
\end{attention}

\courssub{Les cinq pièges classiques des francophones (synthèse)}

\textbf{Piège 1 : oublier \zh{而且}.} En français, « mais aussi » peut être sous-entendu ; en chinois, la structure exige \textbf{les deux éléments}.\\
\zh{*他不但聪明，很努力。} ✗ \hspace{0.8cm} \zh{他不但聪明，而且很努力。} ✓\\
\emph{Tā bùdàn cōngmíng, érqiě hěn nǔlì.} « Non seulement il est intelligent, mais en plus il travaille dur. »

\textbf{Piège 2 : mal placer \zh{不但} quand le sujet est unique.} Le francophone calque « non seulement lui... » et place \zh{不但} en tête :\\
\zh{*不但他聪明，而且很努力。} ✗ \hspace{0.8cm} \zh{他不但聪明，而且很努力。} ✓\\
Test rapide : si le sujet ne change pas entre les deux moitiés, il passe \textbf{devant} \zh{不但}.

\textbf{Piège 3 : oublier \zh{也} avec deux sujets différents.}\\
\zh{*不但我喜欢音乐，而且我弟弟喜欢。} ✗\\
\zh{不但我喜欢音乐，而且我弟弟也喜欢。} ✓\\
\emph{Bùdàn wǒ xǐhuan yīnyuè, érqiě wǒ dìdi yě xǐhuan.} « Non seulement j'aime la musique, mais mon petit frère aussi. » Le \zh{也} est obligatoire : c'est lui qui porte l'idée du « aussi » français.

\textbf{Piège 4 : calquer « et... et... » ou « aussi » par \zh{也} seul.} \zh{也} signifie seulement « aussi, également » ; il n'exprime \textbf{pas la gradation} (« non seulement... mais en plus... »).\\
\zh{*他也聪明，也努力。} signifie simplement « Il est aussi intelligent, et aussi travailleur » — sans idée de progression. Pour la gradation, il faut la paire complète \zh{不但……而且……}.

\textbf{Piège 5 : écrire \zh{第两} au lieu de \zh{第二}} (voir la boîte \emph{Attention} section 2).

\begin{exemplebox}[Une phrase correcte à analyser]
\zh{我不但喜欢音乐，而且喜欢运动。} ✓\\
\emph{Wǒ bùdàn xǐhuan yīnyuè, érqiě xǐhuan yùndòng.}\\
« Non seulement j'aime la musique, mais j'aime aussi le sport. »\\
Cette phrase est \textbf{correcte} : le sujet \zh{我} est unique et précède \zh{不但} ; \zh{而且} est présent ; le verbe \zh{喜欢} est simplement répété dans la seconde partie, ce qui est normal en chinois. Inutile d'ajouter \zh{也} ici, car il n'y a pas de second sujet.
\end{exemplebox}

\begin{methode}[Vérifier sa phrase en trois questions]
Avant de rendre votre copie, interrogez chaque phrase en \zh{不但……而且……} :
\begin{enumerate}
\item \textbf{Même sujet dans les deux moitiés ?} $\rightarrow$ le sujet unique va en tête, devant \zh{不但}, et n'apparaît qu'une fois. \textbf{Sujets différents ?} $\rightarrow$ \zh{不但} + sujet 1, puis \zh{而且} + sujet 2.
\item \textbf{而且 est-il présent ?} Sans lui, la phrase est amputée : \zh{*他不但聪明，很努力} ✗.
\item \textbf{也 est-il présent si les sujets sont différents ?} $\rightarrow$ \zh{而且我弟弟也喜欢} ✓, jamais \zh{*而且我弟弟喜欢} ✗.
\end{enumerate}
\end{methode}

\begin{popfigure}
\begin{center}
\begin{tabularx}{\linewidth}{|X|X|X|}
\hline
\rowcolor{poppinkL} Phrase fautive ✗ & Correction ✓ & Règle rappelée \\
\hline
他不但聪明，很努力 & 他不但聪明，而且很努力 & 而且 obligatoire \\
\hline
不但他聪明，而且很努力 & 他不但聪明，而且很努力 & sujet unique en tête \\
\hline
不但我喜欢音乐，而且我弟弟喜欢 & 不但我喜欢音乐，而且我弟弟也喜欢 & 也 avec 2 sujets \\
\hline
他也聪明，也努力 & 他不但聪明，而且努力 & 也 seul ≠ gradation \\
\hline
第两，汉语很有用 & 第二，汉语很有用 & 第 + 二 (jamais 两) \\
\hline
\end{tabularx}
\end{center}
\legende{Tableau des cinq erreurs typiques des francophones et leurs corrections.}
\end{popfigure}

\courssub{Exercices d'application et de transformation}

\begin{exoresolu}[Corrige les fautes d'un élève de Première A]
Énoncé. Voici la copie de Ngo Bell, élève à Yaoundé. Trois phrases contiennent une faute ; une est correcte. Corrige et justifie.\\
1) \zh{不但他会说法语，而且会说汉语。}\\
2) \zh{喀麦隆不但美丽，而且很丰富。}\\
3) \zh{不但老师喜欢足球，而且学生喜欢。}\\
4) \zh{第两，我们要学习汉字。}
\tcblower
Solution détaillée.\\
1) ✗ $\rightarrow$ \zh{他不但会说法语，而且会说汉语。} \emph{Tā bùdàn huì shuō Fǎyǔ, érqiě huì shuō Hànyǔ.} « Non seulement il parle français, mais il parle aussi chinois. » Le sujet \zh{他} est unique : il doit précéder \zh{不但} et n'apparaître qu'une fois.\\
2) ✓ Correcte : sujet unique \zh{喀麦隆} en tête, \zh{而且} présent. \emph{Kāmàilóng bùdàn měilì, érqiě hěn fēngfù.} « Non seulement le Cameroun est beau, mais il est aussi riche. »\\
3) ✗ $\rightarrow$ \zh{不但老师喜欢足球，而且学生也喜欢。} \emph{Bùdàn lǎoshī xǐhuan zúqiú, érqiě xuésheng yě xǐhuan.} « Non seulement le professeur aime le football, mais les élèves aussi. » Deux sujets différents (\zh{老师} / \zh{学生}) : il faut \zh{也} après le second sujet.\\
4) ✗ $\rightarrow$ \zh{第二，我们要学习汉字。} \emph{Dìèr, wǒmen yào xuéxí Hànzì.} « Deuxièmement, nous devons apprendre les caractères. » Après \zh{第}, on emploie \zh{二}, jamais \zh{两}.
\end{exoresolu}

\begin{exercice}[À toi de jouer : Traduction et transformation]
\begin{enumerate}
\item Traduis en chinois en appliquant la méthode des trois questions :\\
« Non seulement Marie parle anglais, mais elle parle aussi chinois. »
\item Traduis en chinois :\\
« Non seulement mon père aime le ndolé, mais ma mère l'aime aussi. »
\item Traduis en chinois (énumération) :\\
« Premièrement, Douala est grande ; deuxièmement, Douala est intéressante. »
\item Transforme la phrase à sujet unique en phrase à deux sujets (ajoute \zh{我弟弟} comme second sujet) :\\
\zh{我不但喜欢看书，而且喜欢听音乐。}
\end{enumerate}
\end{exercice}

\begin{corrige}[Exercice « À toi de jouer »]
1) Sujet unique (\zh{玛丽}) $\rightarrow$ \zh{玛丽不但会说英语，而且会说汉语。} \emph{Mǎlì bùdàn huì shuō Yīngyǔ, érqiě huì shuō Hànyǔ.}\\
2) Deux sujets (\zh{我爸爸} / \zh{我妈妈}) $\rightarrow$ \zh{不但我爸爸喜欢恩东菜，而且我妈妈也喜欢。} \emph{Bùdàn wǒ bàba xǐhuan Ēndōngcài, érqiě wǒ māma yě xǐhuan.}\\
3) Énumération $\rightarrow$ \zh{第一，杜阿拉很大；第二，杜阿拉很有意思。} \emph{Dìyī, Dù'ālā hěn dà; dìèr, Dù'ālā hěn yǒu yìsi.}\\
4) Transformation (deux sujets) $\rightarrow$ \zh{不但我喜欢看书，而且我弟弟也喜欢听音乐。} \emph{Bùdàn wǒ xǐhuan kànshū, érqiě wǒ dìdi yě xǐhuan tīng yīnyuè.} (Le verbe change, \zh{也} obligatoire).
\end{corrige}

\courssub{Synthèse et prolongement vers l'examen}

\begin{aretenir}[L'essentiel de la leçon ZHA02]
\begin{itemize}
\item \textbf{Énumération} : \zh{第一、第二、第三……最后}. Correspondance directe avec « premièrement, deuxièmement... enfin ». \textbf{Attention :} \zh{第二} (jamais \zh{第两}). Ponctuation : ， ou ； entre les items.
\item \textbf{Gradation} : \zh{不但……而且……} = « non seulement... mais aussi... ». \textbf{Les deux marqueurs sont obligatoires}.
\item \textbf{Règle du sujet} :
      \begin{itemize}
      \item Sujet unique $\rightarrow$ Sujet + \zh{不但}... \zh{而且}... (sujet \textbf{une seule fois} en tête).
      \item Deux sujets $\rightarrow$ \zh{不但} + S1... \zh{而且} + S2 + \zh{也/还}... (\zh{也} \textbf{obligatoire}).
      \end{itemize}
\item \textbf{Distinction} : \zh{也} = « aussi » (addition) ; \zh{不但……而且……} = « non seulement... mais en plus » (gradation/progression).
\item \textbf{Vocabulaire clé} : \zh{第一} (dìyī), \zh{第二} (dìèr), \zh{最后} (zuìhòu), \zh{不但} (bùdàn), \zh{而且} (érqiě), \zh{也} (yě), \zh{恩东菜} (Ēndōngcài, ndolé), \zh{杜阿拉} (Dù'ālā), \zh{雅温得} (Yǎwēndé).
\end{itemize}
\end{aretenir}

\begin{methode}[Stratégies pour l'examen (Bac / Contrôle continu)]
\begin{enumerate}
\item \textbf{Traduction (Thème) :} Avant d'écrire, entourez les sujets dans la phrase française. Décidez : \textit{Un ou deux sujets ?} $\rightarrow$ Placez \zh{不但} et \zh{也} en conséquence.
\item \textbf{Expression écrite} : Utilisez \zh{第一……第二……最后} pour structurer un paragraphe argumentatif (ex: « Pourquoi apprendre le chinois ? »). Cela montre une organisation logique claire, très valorisée.
\item \textbf{Compréhension orale/écrite} : Repérez \zh{不但} pour anticiper \zh{而且} et l'idée de renforcement. Repérez \zh{第一} pour structurer la prise de notes (Point 1, Point 2...).
\item \textbf{Expression orale} : Entraînez-vous à lister 3 atouts de votre ville/école avec \zh{第一、第二、第三、最后}. Enchaînez avec une phrase en \zh{不但……而且……} pour donner votre avis personnel.
\end{enumerate}
\end{methode}

\begin{experience}[Activité orale : Débat « Vivre à Yaoundé vs Vivre à Douala »]
\begin{enumerate}
\item \textbf{Préparation (5 min) :} Par binôme. L'élève A note 3 arguments pour Yaoundé avec \zh{第一、第二、最后}. L'élève B note 3 arguments pour Douala.
\item \textbf{Échange (5 min) :} A présente : \zh{第一……第二……最后……} B répond avec \zh{不但……而且……} pour rebondir (ex: \zh{不但雅温得气候好，而且杜阿拉海鲜也好吃。}).
\item \textbf{Synthèse (3 min) :} Chaque binôme produit une phrase complexe combinant les deux structures : \zh{第一，雅温得气候舒服；第二，杜阿拉经济发达。不但两个城市都美，而且人都热情。}
\end{enumerate}
\textbf{Critères d'évaluation :} Exactitude des structures (ordre sujets, \zh{而且}, \zh{也}), prononciation des tons (bùdàn, érqiě, dìyī, dìèr), richesse lexicale (adjectifs qualitatifs).
\end{experience}

\begin{savaistu}[Le saviez-vous ? La logique de l'argumentation en Chine]
Dans la rhétorique chinoise classique (\zh{议论文}), la structure \zh{第一……第二……第三……} est le standard absolu pour développer une thèse (\zh{论点}). On exige \zh{条理清晰} (structure claire). L'usage de \zh{不但……而且……} permet d'opérer une \zh{层层递进} (progression par paliers), figure de style très prisée pour montrer la profondeur de la réflexion. Au Cameroun, dans les dissertations de philo ou de français, on utilise « Premièrement, Deuxièmement, Enfin » + « Non seulement... mais encore... ». Les deux cultures partagent le goût de l'ordre logique, mais le chinois le grammaticalise de façon plus rigide (marqueurs obligatoires, place figée du sujet).
\end{savaistu}

\coursec{Exercices d'application et de transformation}

Après avoir analysé les mécanismes de l'énumération (\zh{第一……第二……}) et de la gradation (\zh{不但……而且……}), et identifié les pièges classiques pour un francophone (ordre des mots, omission de \zh{而且}, confusion avec \zh{也}), nous passons maintenant à la mise en œuvre active. Cette section propose un parcours progressif : de la reconnaissance guidée à la production autonome, en passant par la correction réflexive et la traduction thématique. Chaque exercice est suivi de son corrigé commenté pour ancrer les règles.

\courssub{Exercice 1 : Complétion à trous (Reconnaissance et choix du connecteur)}

\begin{exercice}[Compléter avec \zh{第一}、\zh{第二}、\zh{不但}、\zh{而且}、\zh{也}]
Complétez les phrases suivantes en choisissant le mot approprié parmi la liste. Chaque mot peut être utilisé une ou plusieurs fois. Écrivez les caractères chinois, le pinyin et la traduction française.
\begin{enumerate}
    \item 我喜欢喀麦隆的菜，\_\_\_\_\_\_ 我喜欢中国菜。
    \item 学习汉语 \_\_\_\_\_\_ 要认真听课，\_\_\_\_\_\_ 要多练习写字。
    \item 我哥哥 \_\_\_\_\_\_ 会说法语，\_\_\_\_\_\_ 会说英语。
    \item 今天天气 \_\_\_\_\_\_ 热，\_\_\_\_\_\_ 下雨。
    \item 去年我去了北京，\_\_\_\_\_\_ 我去了上海，\_\_\_\_\_\_ 我还去了西安。
\end{enumerate}
\end{exercice}

\begin{corrige}[Exercice 1]
\begin{enumerate}
    \item \zh{而且} (érqiě) — 我喜欢喀麦隆的菜，\textbf{而且} 我喜欢中国菜。\\
    (J'aime la cuisine camerounaise, \textbf{de plus / en plus} j'aime la cuisine chinoise.) \textit{Note : énumération de deux faits positifs, sujet identique, structure A 而且 B.}
    \item \zh{第一} (dì yī) ; \zh{第二} (dì èr) — 学习汉语 \textbf{第一} 要认真听课，\textbf{第二} 要多练习写字。\\
    (Apprendre le chinois : \textbf{premièment} il faut écouter attentivement en cours, \textbf{deuxièmement} il faut beaucoup s'entraîner à écrire les caractères.)
    \item \zh{不但} (bùdàn) ; \zh{而且} (érqiě) — 我哥哥 \textbf{不但} 会说法语，\textbf{而且} 会说英语。\\
    (Mon grand frère \textbf{non seulement} parle français, \textbf{mais aussi} l'anglais.) \textit{Structure de gradation : sujet + 不但 + V1, 而且 + V2.}
    \item \zh{不但} (bùdàn) ; \zh{而且} (érqiě) — 今天天气 \textbf{不但} 热，\textbf{而且} 下雨。\\
    (Aujourd'hui il fait \textbf{non seulement} chaud, \textbf{mais en plus} il pleut.) \textit{Gradation d'intensité (chaleur + pluie).}
    \item \zh{第一} (dì yī) ; \zh{第二} (dì èr) — 去年我去了北京，\textbf{第一} 我去了上海，\textbf{第二} 我还去了西安。\\
    (L'an dernier je suis allé à Pékin, \textbf{premièrement} à Shanghai, \textbf{deuxièmement} à Xi'an.) \textit{Énumération chronologique ou logique d'événements. Le mot \zh{也} n'est pas nécessaire ici car \zh{还} (hái, « en plus / aussi ») est déjà présent dans la troisième proposition.}
\end{enumerate}
\end{corrige}

\courssub{Exercice 2 : Transformation — Fusion de deux phrases en gradation (Manipulation syntaxique)}

\begin{methode}[Transformer deux phrases simples en structure \zh{不但……而且……}]
Suivez ces 4 étapes pour réussir la transformation sans erreur :
\begin{enumerate}
    \item \textbf{Identifier le(s) sujet(s).} Si les deux phrases ont le \textbf{même sujet} (ex. : \zh{他}、\zh{我妈妈}), on le place une seule fois en début de phrase. Si les sujets \textbf{sont différents}, on place \zh{不但} devant le premier sujet et \zh{而且} devant le second.
    \item \textbf{Placer \zh{不但} (bùdàn) immédiatement devant le premier prédicat} (verbe ou adjectif) \textbf{ou} devant le premier sujet (si sujets différents).
    \item \textbf{Relier la seconde proposition par \zh{而且} (érqiě)} placé devant le second prédicat \textbf{ou} devant le second sujet.
    \item \textbf{N'oubliez jamais \zh{而且} !} La corrélation est indissociable.
\end{enumerate}
\textbf{Modèle 1 (Sujet unique) :} \zh{他会说汉语。他会写汉字。} $\rightarrow$ \zh{他不但会说汉语，而且会写汉字。}\\
\textbf{Modèle 2 (Sujets différents) :} \zh{我爸爸喜欢喝茶。我妈妈喜欢喝咖啡。} $\rightarrow$ \zh{不但我爸爸喜欢喝茶，而且我妈妈喜欢喝咖啡。}
\end{methode}

\begin{exercice}[Fusionner les phrases avec \zh{不但……而且……}]
Transformez chaque paire de phrases en une seule phrase utilisant la structure de gradation. Respectez l'ordre des idées (du moins fort au plus fort, ou ordre chronologique/logique).
\begin{enumerate}
    \item 我喜欢踢足球。我喜欢打篮球。
    \item 这件衣服很贵。这件衣服很漂亮。
    \item 我爸爸会做中国菜。我爸爸会做喀麦隆菜。
    \item 中文语法不难。中文发音很难。
    \item 昨天我去了市场。昨天我去了图书馆。
\end{enumerate}
\end{exercice}

\begin{corrige}[Exercice 2]
\begin{enumerate}
    \item \zh{我不但喜欢踢足球，而且喜欢打篮球。} (Wǒ bùdàn xǐhuan tī zúqiú, érqiě xǐhuan dǎ lánqiú.) — Je n'aime pas seulement jouer au foot, mais aussi au basket.
    \item \zh{这件衣服不但很贵，而且很漂亮。} (Zhè jiàn yīfu bùdàn hěn guì, érqiě hěn piàoliang.) — Ce vêtement n'est pas seulement cher, il est en plus très beau.
    \item \zh{我爸爸不但会做中国菜，而且会做喀麦隆菜。} (Wǒ bàba bùdàn huì zuò Zhōngguó cài, érqiě huì zuò Kāmàilóng cài.) — Mon père ne sait pas seulement cuisiner chinois, mais aussi camerounais.
    \item \zh{不但中文语法不难，而且中文发音很难。} (Bùdàn Zhōngwén yǔfǎ bù nán, érqiě Zhōngwén fāyīn hěn nán.) — \textbf{Non seulement} la grammaire du chinois n'est pas difficile, \textbf{mais en plus} la prononciation est difficile. \textit{Sujets différents (\zh{语法} vs \zh{发音}) $\rightarrow$ \zh{不但} devant Suj1, \zh{而且} devant Suj2.}
    \item \zh{昨天我不但去了市场，而且去了图书馆。} (Zuótiān wǒ bùdàn qù le shìchǎng, érqiě qù le túshūguǎn.) — Hier je ne suis pas seulement allé au marché, mais aussi à la bibliothèque.
\end{enumerate}
\end{corrige}

\begin{exoresolu}[Analyse d'une transformation piège : sujets différents]
\textbf{Énoncé :} Transformer : « Mon père aime le thé. Ma mère aime le café. »
\textbf{Tentative erronée :} \zh{我爸爸不但喜欢喝茶，而且我妈妈喜欢喝咖啡。} $\leftarrow$ \textbf{Incorrect} (structure bancale : \zh{不但} après Suj1 mais \zh{而且} devant Suj2 ; \zh{不但} doit être devant Suj1 si sujets différents).
\textbf{Solution correcte :} \zh{不但我爸爸喜欢喝茶，而且我妈妈喜欢喝咖啡。} (Bùdàn wǒ bàba xǐhuan hē chá, érqiě wǒ māma xǐhuan hē kāfēi.) — \textbf{Non seulement} mon père aime le thé, \textbf{mais aussi} ma mère aime le café.
\textbf{Règle :} Si les sujets sont différents, placer \zh{不但} \textbf{devant le premier sujet} et \zh{而且} \textbf{devant le second sujet}.
\end{exoresolu}

\courssub{Exercice 3 : Correction d'erreurs fréquentes (Réflexion métalinguistique)}

\begin{exercice}[Corriger et justifier]
Les phrases ci-dessous contiennent chacune une erreur typique d'apprenant francophone. Corrigez la phrase (caractères + pinyin) et justifiez l'erreur en une phrase en français.
\begin{enumerate}
    \item \zh{我一喜欢吃米饭，二喜欢吃面条。}
    \item \zh{他不但会说汉语，也会写汉字。}
    \item \zh{不但今天很冷，而且下雪了。}
    \item \zh{学习中文第一要背单词，第二要语法。}
\end{enumerate}
\end{exercice}

\begin{corrige}[Exercice 3]
\begin{enumerate}
    \item \textbf{Correction :} \zh{我\textbf{第一}喜欢吃米饭，\textbf{第二}喜欢吃面条。} (Wǒ \textbf{dì yī} xǐhuan chī mǐfàn, \textbf{dì èr} xǐhuan chī miàntiáo.) \\
    \textbf{Justification :} Les ordinaux exigent le préfixe \zh{第} (dì) : \zh{第一}、\zh{第二}. *\zh{一}、*\zh{二} sont des cardinaux (nombres), pas des ordinaux (rangs).
    \item \textbf{Correction :} \zh{他不但会说汉语，\textbf{而且}会写汉字。} (Tā bùdàn huì shuō Hànyǔ, \textbf{érqiě} huì xiě Hànzì.) \\
    \textbf{Justification :} \zh{不但} appelle obligatoirement \zh{而且} (érqiě) pour former la corrélation de gradation. \zh{也} (yě, « aussi ») ne peut pas remplacer \zh{而且} après \zh{不但}. C'est une erreur de corrélation (paire fixe).
    \item \textbf{Correction :} \zh{\textbf{今天不但很冷，而且下雪了。}} (Jīntiān \textbf{bùdàn} hěn lěng, érqiě xià xuě le.) \\
    \textbf{Justification :} Lorsque le sujet est commun (\zh{今天}), \zh{不但} doit être placé \textbf{après le sujet} et non en tout début de phrase (sauf si sujets différents, cf. Exercice 2 résolu). Ici, « Aujourd'hui » est le sujet/thème.
    \item \textbf{Correction :} \zh{学习中文\textbf{第一}要背单词，\textbf{第二}要\textbf{学}语法。} (Xuéxí Zhōngwén \textbf{dì yī} yào bèi dāncí, \textbf{dì èr} yào \textbf{xué} yǔfǎ.) \\
    \textbf{Justification :} 1. Manque de \zh{第} devant \zh{一} et \zh{二}. 2. Ellipse du verbe dans le second membre : \zh{要语法} est incomplet (verbe manquant), il faut \zh{要学语法} (apprendre la grammaire) ou \zh{要懂语法}. Parallélisme structurel obligatoire : Verbe + Objet / Verbe + Objet.
\end{enumerate}
\end{corrige}

\courssub{Exercice 4 : Thème — Traduction dirigée vers le chinois}

\begin{exercice}[Traduire en chinois (caractères + pinyin)]
Traduisez les deux phrases suivantes en utilisant obligatoirement les structures étudiées.
\begin{enumerate}
    \item « J'apprends le chinois : premièrement, il est intéressant ; deuxièmement, il est utile. »
    \item « Non seulement mon père aime la Chine, mais ma mère l'aime aussi. »
\end{enumerate}
\end{exercice}

\begin{corrige}[Exercice 4]
\begin{enumerate}
    \item \zh{我学中文，第一很有意思，第二很有用。} (Wǒ xué Zhōngwén, dì yī hěn yǒu yìsi, dì èr hěn yǒuyòng.) \\
    \textit{Analyse :} Sujet \zh{我学中文} (thème) + énumération des raisons via \zh{第一/第二} + adjectifs prédicatifs (\zh{很有意思}、\zh{很有用}). Attention : \zh{有意思} (intéressant) et \zh{有用} (utile) sont des adjectifs statifs, pas de \zh{是} (shì) devant.
    \item \zh{不但我爸爸喜欢中国，而且我妈妈也喜欢中国。} (Bùdàn wǒ bàba xǐhuan Zhōngguó, érqiě wǒ māma yě xǐhuan Zhōngguó.) \\
    \textit{Analyse :} Sujets différents (\zh{我爸爸} vs \zh{我妈妈}) $\rightarrow$ \zh{不但} placé \textbf{devant le premier sujet}, \zh{而且} placé \textbf{devant le second sujet}. L'ajout de \zh{也} (yě) dans le second membre renforce l'analogie (« aussi »), c'est correct et naturel.
\end{enumerate}
\end{corrige}

\courssub{Exercice 5 : Réordonnancement (Sensibilisation à l'ordre des mots)}

\begin{exercice}[Remettre les mots dans l'ordre correct]
Réorganisez les éléments pour former une phrase grammaticale correcte. Écrivez la phrase en caractères, pinyin et français.
\begin{enumerate}
    \item [ \zh{而且} / \zh{很聪明} / \zh{不但} / \zh{他} / \zh{很努力} ]
    \item [ \zh{第一} / \zh{复习课文} / \zh{第二} / \zh{我们} / \zh{预习生词} ]
    \item [ \zh{不但} / \zh{下雨} / \zh{而且} / \zh{今天} / \zh{刮风} ]
    \item [ \zh{而且} / \zh{会唱歌} / \zh{我妹妹} / \zh{不但} / \zh{会跳舞} ]
    \item [ \zh{第三} / \zh{去图书馆} / \zh{第一} / \zh{回家} / \zh{第二} / \zh{吃饭} / \zh{我} ]
\end{enumerate}
\end{exercice}

\begin{corrige}[Exercice 5]
\begin{enumerate}
    \item \zh{他不但很努力，而且很聪明。} (Tā bùdàn hěn nǔlì, érqiě hěn cōngmíng.) — Il est non seulement travailleur, mais en plus intelligent.
    \item \zh{我们第一预习生词，第二复习课文。} (Wǒmen dì yī yùxí shēngcí, dì èr fùxí kèwén.) — Nous, premièrement prévisons les mots nouveaux, deuxièmement révisons le texte. \textit{Ordre logique d'étude.}
    \item \zh{今天不但下雨，而且刮风。} (Jīntiān bùdàn xià yǔ, érqiě guā fēng.) — Aujourd'hui il ne pleut pas seulement, mais en plus il y a du vent. \textit{Sujet \zh{今天} avant \zh{不但}.}
    \item \zh{我妹妹不但会跳舞，而且会唱歌。} (Wǒ mèimei bùdàn huì tiàowǔ, érqiě huì chànggē.) — Ma petite sœur ne sait pas seulement danser, mais aussi chanter.
    \item \zh{我第一回家，第二吃饭，第三去图书馆。} (Wǒ dì yī huí jiā, dì èr chī fàn, dì sān qù túshūguǎn.) — Je rentre d'abord chez moi, puis je mange, ensuite je vais à la bibliothèque. \textit{Énumération chronologique d'actions (série verbale).}
\end{enumerate}
\end{corrige}

\courssub{Exercice 6 : Production guidée — Expression écrite (Transfert et créativité)}

\begin{exercice}[Rédiger un court paragraphe : « Pourquoi j'apprends le chinois »]
Rédigez un texte de 6 à 8 lignes (environ 60-80 caractères) en chinois simplifié.
\textbf{Contraintes obligatoires :}
\begin{itemize}
    \item Utilisez \zh{第一}、\zh{第二}、\zh{第三} pour structurer trois raisons distinctes.
    \item Utilisez \textbf{une fois} la structure \zh{不但……而且……} à l'intérieur du texte (pour développer une raison ou en ajouter une quatrième).
    \item Utilisez du vocabulaire du module (famille, école, Cameroun, Chine, avenir, travail, culture).
    \item Soignez l'ordre des mots (Sujet + Connecteur + Prédicat) et la ponctuation chinoise.
\end{itemize}
\end{exercice}

\begin{corrige}[Exercice 6 — Proposition de production type (Modèle pour l'enseignant)]
\textbf{Texte élève (exemple) :}
\begin{textebox}[Production élève type : 为什么我学中文]
我学中文，\textbf{第一}，中文不但很有意思，\textbf{而且}很有用。\\
\textbf{第二}，我想将来去中国留学，在北京大学学习国际贸易。\\
\textbf{第三}，我爸爸在杜阿拉做生意，很多中国客户说法语不流利，我要帮他翻译。\\
所以，我每天都努力练习听说读写。
\end{textebox}

\textbf{Pinyin :}
Wǒ xué Zhōngwén, \textbf{dì yī}, Zhōngwén \textbf{bùdàn} hěn yǒu yìsi, \textbf{érqiě} hěn yǒuyòng.
\textbf{Dì èr}, wǒ xiǎng jiānglái qù Zhōngguó liúxué, zài Běijīng Dàxué xuéxí guójì màoyì.
\textbf{Dì sān}, wǒ bàba zài Dū'ālā zuò shēngyì, hěn duō Zhōngguó kèhù shuō Fǎyǔ bù liúlì, wǒ yào bāng tā fānyì.
Suǒyǐ, wǒ měitiān dōu nǔlì liànxī tīng shuō dú xiě.

\textbf{Traduction :}
J'apprends le chinois, \textbf{premièrement}, le chinois est non seulement intéressant, \textbf{mais aussi} très utile. \textbf{Deuxièmement}, je veux aller étudier en Chine plus tard, à l'Université de Pékin, pour le commerce international. \textbf{Troisièmement}, mon père fait du commerce à Douala, beaucoup de clients chinois ne parlent pas couramment français, je dois l'aider à traduire. C'est pourquoi, chaque jour, je m'entraîne dur à l'écoute, l'oral, la lecture et l'écriture.

\textbf{Analyse didactique du modèle :}
\begin{itemize}
    \item Respect strict des 3 \zh{第一/二/三} + 1 \zh{不但……而且……} (imbriqué dans la 1ère raison).
    \item Vocabulaire HSK 2-3 : \zh{将来} (futur), \zh{留学} (étudier à l'étranger), \zh{国际贸易} (commerce international), \zh{做生意} (faire du commerce), \zh{客户} (client), \zh{流利} (courant), \zh{翻译} (traduire), \zh{所以} (donc), \zh{每天} (chaque jour).
    \item Structure \zh{想……} (vouloir), \zh{要……} (devoir/avoir l'intention), \zh{练习……} (s'entraîner à).
\end{itemize}
\end{corrige}

\courssub{Lexique de la leçon (Mots nouveaux et révisés)}

\begin{definition}[Vocabulaire clé des exercices]
\begin{center}
\begin{tabularx}{\linewidth}{|X|X|X|X|}
\hline
\rowcolor{popblueL}
\textbf{Caractères} & \textbf{Pinyin} & \textbf{Nature} & \textbf{Traduction} \\
\hline
踢足球 & tī zúqiú & VO & jouer au football \\
打篮球 & dǎ lánqiú & VO & jouer au basket-ball \\
贵 & guì & Adj & cher \\
漂亮 & piàoliang & Adj & beau, joli \\
做生意 & zuò shēngyì & VO & faire du commerce / des affaires \\
客户 & kèhù & N & client \\
流利 & liúlì & Adj & courant (langue) \\
翻译 & fānyì & V/N & traduire / traduction \\
预习 & yùxí & V & préparer (une leçon), prévoir \\
复习 & fùxí & V & réviser \\
生词 & shēngcí & N & mot nouveau, vocabulaire nouveau \\
课文 & kèwén & N & texte de leçon \\
下雨 & xià yǔ & VO & pleuvoir \\
刮风 & guā fēng & VO & venter (il y a du vent) \\
跳舞 & tiàowǔ & VO & danser \\
唱歌 & chànggē & VO & chanter \\
回家 & huí jiā & VO & rentrer à la maison \\
去图书馆 & qù túshūguǎn & VO & aller à la bibliothèque \\
有意思 & yǒu yìsi & Adj & intéressant \\
有用 & yǒuyòng & Adj & utile \\
将来 & jiānglái & Adv/N & futur, plus tard \\
留学 & liúxué & V & étudier à l'étranger \\
国际贸易 & guójì màoyì & N & commerce international \\
努力 & nǔlì & Adj/V & travailleur / faire des efforts \\
\hline
\end{tabularx}
\end{center}
\end{definition}

\courssub{Synthèse et prolongement vers l'examen}

\begin{aretenir}[Tableau récapitulatif comparatif : \zh{第一……第二……} vs \zh{不但……而且……}]
\begin{center}
\begin{tabularx}{\linewidth}{|X|X|X|X|}
\hline
\rowcolor{popblueL}
\textbf{Critère} & \textbf{\zh{第一……第二…… (Énumération)}} & \textbf{\zh{不但……而且…… (Gradation / Progression)}} & \textbf{Point de vigilance (Francophone)} \\
\hline
\textbf{Fonction} & Lister des éléments sur un plan \textbf{égal} (raisons, étapes, actions successives). & Mettre en relief un \textbf{deuxième élément plus fort, inattendu ou additionnel} par rapport au premier. & Ne pas confondre « liste » et « insistance ». \\
\hline
\textbf{Structure type} & Sujet + \zh{第一} + Pred\_1 ， \zh{第二} + Pred\_2 (, \zh{第三} + Pred\_3...). & Sujet + \zh{不但} + Pred\_1 ， \zh{而且} + Pred\_2. & \zh{不但} JAMAIS sans \zh{而且}. \\
\hline
\textbf{Sujets différents} & Rare / Maladroit (préférer \zh{一个……一个……}). & Possible : \zh{不但} + Suj\_1 + Pred\_1， \zh{而且} + Suj\_2 + Pred\_2. & Placer \zh{不但} \textbf{devant} le 1er sujet. \\
\hline
\textbf{Négation} & \zh{不第一……} (rare). & \zh{不但不……而且……} (ex: \zh{不但不贵，而且很好}). & Ne pas mettre \zh{不} devant \zh{不但}. \\
\hline
\textbf{Exemple clé} & \zh{买菜第一看价钱，第二看新鲜。} & \zh{这道菜不但好吃，而且不贵。} & \textit{Vérifier la paire corrélative à l'écrit.} \\
\hline
\end{tabularx}
\end{center}
\end{aretenir}

\begin{attention}[Pièges « Zéro tolérance » à l'examen officiel (MINESEC / Baccalauréat)]
\begin{enumerate}
    \item \textbf{Omission de \zh{而且} :} Écrire \zh{他不但会说汉语，也会写汉字。} $\rightarrow$ \textbf{0 point} sur la phrase. La corrélation est \textbf{indissociable} : \zh{不但……而且……}.
    \item \textbf{Mauvaise place de \zh{不但} (sujet commun) :} Écrire \zh{不但他会说汉语，而且会写汉字。} $\rightarrow$ \textbf{Erreur de syntaxe}. Avec un sujet unique : \zh{Sujet + 不但 + V1, 而且 + V2}.
    \item \textbf{Oubli de \zh{第} pour les ordinaux :} Écrire \zh{一、二、三} au lieu de \zh{第一、第二、第三}. $\rightarrow$ \textbf{Erreur de morphologie} (cardinal vs ordinal).
    \item \textbf{Parallélisme brisé :} \zh{第一要背单词，第二语法。} (Verbe manquant dans 2e membre). $\rightarrow$ \textbf{Erreur de structure}. Exiger \zh{第二要学语法} ou \zh{第二重语法}.
    \item \textbf{Ponctuation :} Utiliser la virgule française (,) au lieu de la virgule chinoise (，) ou point-virgule (;) au lieu de ；. $\rightarrow$ Pénalité présentation.
\end{enumerate}
\end{attention}

\begin{popfigure}
\begin{tikzpicture}[
    node distance=1.2cm and 1.5cm,
    every node/.style={font=\small, align=center},
    stage/.style={rectangle, rounded corners, draw=popblue, thick, fill=popblueL, minimum width=2.8cm, minimum height=1.1cm},
    arrow/.style={-Stealth, thick, draw=popdark}
]
% Nœuds
\node[stage, align=center] (rec){Reconnaissance\\ (Ex. 1 : Complétion)};
\node[stage, right=of rec, align=center] (manip){Manipulation\\ (Ex. 2 \& 5 : Transformation,\\ Réordonnancement)};
\node[stage, right=of manip, align=center] (corr){Correction Réflexive\\ (Ex. 3 : Erreurs francophones)};
\node[stage, right=of corr, align=center] (trad){Transfert / Thème\\ (Ex. 4 : Traduction FR$\rightarrow$ZH)};
\node[stage, right=of trad, align=center] (prod){Production Autonome\\ (Ex. 6 : Rédaction guidée)};

% Flèches
\draw[arrow] (rec) -- (manip) node[midway, above, font=\tiny, text=popblue] {Identifier};
\draw[arrow] (manip) -- (corr) node[midway, above, font=\tiny, text=popblue] {Analyser};
\draw[arrow] (corr) -- (trad) node[midway, above, font=\tiny, text=popblue] {Traduire};
\draw[arrow] (trad) -- (prod) node[midway, above, font=\tiny, text=popblue] {Créer};

% Flèche globale de progression (sous les boîtes)
\draw[popline, very thick, decorate, decoration={snake, amplitude=1pt, segment length=6pt}] 
    ($(rec.south west) + (-0.3, -0.6)$) -- ($(prod.south east) + (0.3, -0.6)$)
    node[midway, below=0.4cm, font=\tiny\bfseries, text=poporange] {Complexité cognitive croissante $\rightarrow$};
\end{tikzpicture}
\legende{Frise de progression pédagogique de la leçon : du guidé vers l'autonome (Taxonomie de Bloom adaptée).}
\end{popfigure}

\begin{methode}[Stratégie « Bac » : Réussir l'exercice de transformation / thème]
\begin{enumerate}
    \item \textbf{Lire} la consigne et identifier la structure cible (\zh{第一/二} ou \zh{不但/而且}).
    \item \textbf{Surligner} le(s) sujet(s) et les prédicats (verbes/adjectifs) dans le texte source (français ou chinois).
    \item \textbf{Construire} l'ossature chinoise \textbf{avant} d'écrire les caractères :
        \begin{itemize}
            \item \textit{Énumération :} [Sujet] \zh{第一} [Pred1] ， \zh{第二} [Pred2] 。
            \item \textit{Gradation (1 sujet) :} [Sujet] \zh{不但} [Pred1] ， \zh{而且} [Pred2] 。
            \item \textit{Gradation (2 sujets) :} \zh{不但} [Suj1] [Pred1] ， \zh{而且} [Suj2] [Pred2] 。
        \end{itemize}
    \item \textbf{Remplir} le vocabulaire précis (pinyin $\rightarrow$ caractères).
    \item \textbf{Relire} : Vérifier la paire \zh{不但/而且}, le \zh{第}, la ponctuation ， 。, l'ordre Sujet-Connecteur-Prédicat.
\end{enumerate}
\end{methode}

\coursec{Synthèse et prolongement vers l'examen}

Après les exercices d'application et de transformation de la section précédente, tu sais maintenant manipuler les deux structures séparément. Cette dernière section a trois objectifs : \textbf{fixer} les deux structures dans une carte de synthèse, \textbf{enrichir} ton expression avec des variantes soutenues valorisées à l'écrit, et \textbf{t'entraîner} à combiner les deux structures dans un mini-texte argumentatif, exactement comme on te le demandera à l'examen.

\courssub{Carte de synthèse des deux structures}

\begin{aretenir}[Les deux structures en un coup d'œil]
\begin{itemize}
\item \cle{第一……第二……} sert à \textbf{énumérer} des arguments, des raisons ou des étapes, dans un ordre clair. On peut continuer avec \zh{第三} et conclure avec \zh{最后} (« enfin »).
\item \cle{不但……而且……} sert à \textbf{ajouter} une information plus forte que la première (gradation) : « non seulement... mais aussi... ». Si le sujet est le même, il se place \textbf{avant} \zh{不但} ; si les deux sujets sont différents, chaque sujet se place \textbf{après} son connecteur.
\end{itemize}
\end{aretenir}

\begin{exemplebox}[Exemples types à retenir par cœur]
\zh{我喜欢这所学校，第一，老师很好；第二，校园很大。}\\
\emph{Wǒ xǐhuan zhè suǒ xuéxiào, dì-yī, lǎoshī hěn hǎo; dì-èr, xiàoyuán hěn dà.}\\
« J'aime cette école : premièrement, les professeurs sont bons ; deuxièmement, le campus est grand. »

\medskip
\zh{他不但会说英语，而且会说中文。}\\
\emph{Tā bùdàn huì shuō Yīngyǔ, érqiě huì shuō Zhōngwén.}\\
« Non seulement il sait parler anglais, mais il sait aussi parler chinois. » (même sujet : sujet avant \zh{不但})

\medskip
\zh{不但老师会说法语，而且学生也会说。}\\
\emph{Bùdàn lǎoshī huì shuō Fǎyǔ, érqiě xuésheng yě huì shuō.}\\
« Non seulement le professeur sait parler français, mais les élèves aussi. » (sujets différents : chaque sujet après son connecteur)
\end{exemplebox}

\begin{popfigure}
\begin{tikzpicture}[mindmap, grow cyclic, every node/.style={concept, align=center, font=\small},
root concept/.append style={concept color=popblue, font=\bfseries},
level 1/.append style={level distance=3.4cm, sibling angle=60},
level 2/.append style={level distance=2.4cm, sibling angle=45}]
\node[root concept, align=center]{Argumenter\\en chinois}
  child[concept color=poporange] { node {énumérer}
    child[concept color=poporangeL] { node[align=center]{\zh{第一}\\premièrement} }
    child[concept color=poporangeL] { node[align=center]{\zh{第二}\\deuxièmement} }
    child[concept color=poporangeL] { node[align=center]{\zh{第三}\\troisièmement} }
    child[concept color=poporangeL] { node[align=center]{\zh{最后}\\enfin} }
  }
  child[concept color=popgreen] { node {gradation}
    child[concept color=popgreenL] { node[align=center]{cas A : même sujet\\Sujet + \zh{不但}... \zh{而且}...} }
    child[concept color=popgreenL] { node[align=center]{cas B : sujets différents\\\zh{不但} + S1, \zh{而且} + S2} }
  };
\end{tikzpicture}
\legende{Carte mentale finale : deux outils pour argumenter en chinois.}
\end{popfigure}

\courssub{Tableau contrastif final}

\begin{center}
\begin{tabularx}{\linewidth}{|X|X|X|}
\hline
\rowcolor{popblueL} Critère & \zh{第一……第二……} (énumération) & \zh{不但……而且……} (gradation) \\
\hline
Fonction & Lister des raisons ou des étapes & Ajouter une information plus forte \\
\hline
Équivalent français & « premièrement... deuxièmement... » & « non seulement... mais aussi... » \\
\hline
Place du sujet & Sujet unique en tête de phrase & Avant \zh{不但} si même sujet ; après chaque connecteur si sujets différents \\
\hline
Ponctuation & Virgule \zh{，} ou point-virgule \zh{；} entre les points & Virgule \zh{，} entre les deux parties \\
\hline
Exemple & \zh{第一，学校很大；第二，老师很好。} \emph{Dì-yī, xuéxiào hěn dà; dì-èr, lǎoshī hěn hǎo.} & \zh{我不但学中文，而且学英文。} \emph{Wǒ bùdàn xué Zhōngwén, érqiě xué Yīngwén.} \\
\hline
Traduction & « Premièrement, l'école est grande ; deuxièmement, les professeurs sont bons. » & « Je n'apprends pas seulement le chinois, mais aussi l'anglais. » \\
\hline
\end{tabularx}
\end{center}

\begin{attention}[Piège n° 1 des francophones à l'examen]
Ne confonds pas les deux logiques. \zh{第一……第二……} met les arguments \textbf{côte à côte} (ils ont le même poids) ; \zh{不但……而且……} \textbf{monte en puissance} (le deuxième élément est plus fort). Dire « \zh{我不但第一学中文，而且第二学英文} » est un mélange incorrect des deux structures. Choisis l'une ou l'autre, jamais les deux dans la même paire de connecteurs.
\end{attention}

\courssub{Complément utile à l'examen : les variantes soutenues}

\begin{propriete}[Variantes soutenues à l'écrit]
\begin{itemize}
\item \cle{首先……其次……} (\emph{shǒuxiān... qícì...}) = variante soutenue de \zh{第一……第二……}, idéale dans une rédaction : « d'abord... ensuite... ».
\item \cle{不仅……而且……} (\emph{bùjǐn... érqiě...}) = variante soutenue de \zh{不但……而且……}, très appréciée dans les textes écrits.
\end{itemize}
Les règles de place du sujet sont exactement les mêmes que pour les structures de base.
\end{propriete}

\begin{exemplebox}[Les variantes en action]
\zh{首先，我们要好好学习；其次，我们要帮助家人。}\\
\emph{Shǒuxiān, wǒmen yào hǎohao xuéxí; qícì, wǒmen yào bāngzhù jiārén.}\\
« D'abord, nous devons bien étudier ; ensuite, nous devons aider notre famille. »

\medskip
\zh{雅温得不仅很大，而且很漂亮。}\\
\emph{Yǎwēndé bùjǐn hěn dà, érqiě hěn piàoliang.}\\
« Yaoundé est non seulement grande, mais aussi très belle. »
\end{exemplebox}

\begin{savaistu}[La marque d'un propos bien structuré]
En Chine, dans les dissertations et les examens, énumérer ses arguments avec \zh{第一、第二、第三} est la marque d'un propos clair et bien structuré : les correcteurs y sont très sensibles. Un élève camerounais qui organise sa rédaction chinoise avec \zh{首先……其次……最后……} donne immédiatement l'impression d'un candidat sérieux et méthodique. C'est un réflexe facile à prendre, et il rapporte des points !
\end{savaistu}

\courssub{Combiner les deux structures : le modèle de mini-texte}

Observe ce texte modèle : il combine la gradation et l'énumération, exactement le type de paragraphe attendu à l'examen.

\begin{textebox}[Mon école — texte modèle]
\zh{我的学校不但很大，而且很干净。我喜欢这里，第一，老师很好；第二，同学很友好；第三，我们可以学中文。}\\
\emph{Wǒ de xuéxiào bùdàn hěn dà, érqiě hěn gānjìng. Wǒ xǐhuan zhèlǐ, dì-yī, lǎoshī hěn hǎo; dì-èr, tóngxué hěn yǒuhǎo; dì-sān, wǒmen kěyǐ xué Zhōngwén.}\\
« Mon école est non seulement grande, mais aussi propre. J'aime cet endroit : premièrement, les professeurs sont bons ; deuxièmement, les camarades sont amicaux ; troisièmement, on peut y apprendre le chinois. »
\end{textebox}

\begin{methode}[Construire un mini-texte argumentatif en 4-5 phrases]
\begin{enumerate}
\item \textbf{Phrase d'ouverture avec gradation} : présente ton sujet avec \zh{不但……而且……} (ou \zh{不仅……而且……}) pour donner deux qualités fortes.
\item \textbf{Phrase de transition} : annonce ton point de vue (\zh{我喜欢这里。} / \zh{我觉得……}).
\item \textbf{Énumération des arguments} : développe avec \zh{第一……第二……第三……}, chaque argument séparé par un point-virgule \zh{；}.
\item \textbf{Conclusion éventuelle} : termine avec \zh{最后} ou une phrase de synthèse courte.
\item \textbf{Vérification} : contrôle la place des sujets autour de \zh{不但} et la ponctuation chinoise pleine chasse.
\end{enumerate}
\end{methode}

\begin{exoresolu}[Rédaction guidée : « Ma ville »]
Énoncé : écris un mini-texte de 4 phrases sur ta ville en utilisant \zh{不但……而且……} puis \zh{第一……第二……}.
\tcblower
Solution détaillée :
\begin{enumerate}
\item Ouverture avec gradation (même sujet, donc sujet avant \zh{不但}) : \zh{杜阿拉不但很大，而且很热闹。} \emph{Dù'ālā bùdàn hěn dà, érqiě hěn rènao.} « Douala est non seulement grande, mais aussi très animée. »
\item Transition : \zh{我很喜欢这个城市。} \emph{Wǒ hěn xǐhuan zhè ge chéngshì.} « J'aime beaucoup cette ville. »
\item Énumération : \zh{第一，这里的人很友好；第二，市场很大，东西很多。} \emph{Dì-yī, zhèlǐ de rén hěn yǒuhǎo; dì-èr, shìchǎng hěn dà, dōngxi hěn duō.} « Premièrement, les gens d'ici sont amicaux ; deuxièmement, les marchés sont grands et il y a beaucoup de choses. »
\item Vérification : le sujet \zh{杜阿拉} précède \zh{不但} (même sujet) ; les arguments sont séparés par \zh{；} ; la ponctuation est chinoise. Le texte est correct.
\end{enumerate}
\end{exoresolu}

\begin{attention}[Dernières erreurs à éliminer avant l'examen]
\begin{itemize}
\item Ne dis jamais \zh{很不但} : \zh{不但} se place avant le verbe ou l'adjectif, jamais après \zh{很}.
\item Dans la gradation avec sujets différents, n'oublie pas \zh{也} dans la deuxième partie : \zh{不但老师来，而且学生也来。}
\item N'utilise pas la virgule française « , » : en chinois, on écrit \zh{，} et le point-virgule \zh{；}.
\item \zh{第一} s'écrit avec le préfixe ordinal \zh{第} : sans lui, \zh{一} signifie simplement « un ».
\end{itemize}
\end{attention}

\courssub{Auto-évaluation : je sais faire}

Coche mentalement chaque ligne. Si une case reste vide, revois la section correspondante.

\begin{center}
\begin{tabularx}{\linewidth}{|X|X|}
\hline
\rowcolor{popblueL} Je sais... & Oui / À revoir \\
\hline
Énumérer des arguments avec \zh{第一……第二……第三……} & \\
\hline
Conclure une énumération avec \zh{最后} & \\
\hline
Placer correctement le sujet avec \zh{不但……而且……} (même sujet) & \\
\hline
Construire la gradation avec deux sujets différents (\zh{不但} + S1, \zh{而且} + S2 + \zh{也}) & \\
\hline
Employer les variantes soutenues \zh{首先……其次……} et \zh{不仅……而且……} & \\
\hline
Combiner les deux structures dans un mini-texte argumentatif & \\
\hline
Traduire ces structures du chinois vers le français et inversement & \\
\hline
\end{tabularx}
\end{center}

\courssub{Devoir à la maison}

\begin{exercice}[Rédaction : « Mon école »]
Rédige un paragraphe de 5 à 6 phrases intitulé \zh{我的学校} (\emph{Wǒ de xuéxiào}, « Mon école ») en respectant les consignes suivantes :
\begin{enumerate}
\item une phrase d'ouverture avec \zh{不但……而且……} (ou la variante \zh{不仅……而且……}) ;
\item une énumération d'au moins trois arguments avec \zh{第一……第二……第三……} (ou \zh{首先……其次……}) ;
\item une conclusion courte ;
\item ponctuation chinoise obligatoire (\zh{，。；}) ;
\item écris le texte en caractères, puis ajoute le pinyin et la traduction française.
\end{enumerate}
Tu peux t'inspirer du texte modèle de cette leçon, mais personnalise-le : parle de ton lycée, de tes professeurs, de tes camarades, des matières que tu apprends.
\end{exercice}

\begin{corrige}[Devoir à la maison — éléments de correction]
Un bon devoir contient : (1) une gradation correcte avec le sujet avant \zh{不但} si le sujet est unique ; (2) au moins trois arguments numérotés séparés par \zh{；} ; (3) une conclusion avec \zh{最后} ou une phrase de synthèse ; (4) aucune virgule française. Exemple de réponse attendue : \zh{我的学校不但很漂亮，而且很有名。我非常喜欢我的学校。第一，老师很认真；第二，同学很友好；第三，我们可以学中文、英文和法文。最后，我们的校园也很大。欢迎来到我的学校！} \emph{Wǒ de xuéxiào bùdàn hěn piàoliang, érqiě hěn yǒumíng. Wǒ fēicháng xǐhuan wǒ de xuéxiào. Dì-yī, lǎoshī hěn rènzhēn; dì-èr, tóngxué hěn yǒuhǎo; dì-sān, wǒmen kěyǐ xué Zhōngwén, Yīngwén hé Fǎwén. Zuìhòu, wǒmen de xiàoyuán yě hěn dà. Huānyíng lái dào wǒ de xuéxiào!} « Mon école est non seulement belle, mais aussi réputée. J'aime beaucoup mon école. Premièrement, les professeurs sont sérieux ; deuxièmement, les camarades sont amicaux ; troisièmement, on peut y apprendre le chinois, l'anglais et le français. Enfin, notre campus est aussi très grand. Bienvenue dans mon école ! »
\end{corrige}

\newpage

\coursec{Activité d'intégration}

\coursec{Leçon 2 — Grammaire : \zh{第一……第二……} ; \zh{不但……而且……}}

\courssub{Objectifs de l'activité}
\noindent
À l'issue de cette activité d'intégration, l'élève sera capable de :
\begin{itemize}
    \item \textbf{Mobiliser} les structures \cle{第一……第二……} (énumération ordonnée) et \cle{不但……而且……} (gradation additive) dans une production orale et écrite cohérente.
    \item \textbf{Argumenter} de manière structurée sur un sujet de société familier : l'apprentissage du chinois au Cameroun.
    \item \textbf{Coopérer} en petit groupe pour préparer une prise de parole courte (6 à 8 phrases), respecter une contrainte syntaxique forte et évaluer les productions des pairs via une grille d'observation.
    \item \textbf{Rédiger} un paragraphe argumentatif individuel reprenant les idées du débat, en soignant l'enchaînement logique et la correction grammaticale (HSK 2-3).
\end{itemize}

\courssub{Situation}
\noindent
\textbf{Contexte :} Le lycée organise la \emph{Semaine des Langues Vivantes}. Le club de chinois « \zh{龙的传人} » (Lóng de chuánrén — Descendants du Dragon) tient un stand d'information dans le hall principal. Un groupe d'élèves de Seconde hésite à choisir le chinois en LV2 pour l'année prochaine. Ils demandent : « \zh{为什么要学中文？法语和英语不够吗？} » (Pourquoi apprendre le chinois ? Le français et l'anglais ne suffisent-ils pas ?).

\noindent
Votre groupe (4 élèves) représente le club. Vous devez convaincre les indécis en présentant \textbf{au moins trois arguments solides}, structurés et illustrés par des exemples concrets (commerce sino-camerounais, bourses d'études, culture, tourisme, opportunités professionnelles à Douala/Yaoundé/Buea).

\begin{textebox}[Document déclencheur : Témoignage d'un ancien élève]
\zh{我叫保罗，我在雅温得上大学。} \\
\zh{第一，学中文帮我找到了在华为喀麦隆的实习。} \\
\zh{第二，我获得了中国政府奖学金，去北京留学一年。} \\
\zh{第三，我不但会说中文，而且认识很多汉字。} \\
\zh{我的中国同学不但热情，而且很乐意帮我复习功课。} \\
\zh{所以，学中文不但有用，而且很有意思！} \\
\noindent
\textbf{Pinyin :} Wǒ jiào Bǎoluó, wǒ zài Yǎwēndé shàng dàxué. Dì yī, xué Zhōngwén bāng wǒ zhǎodào le zài Huáwéi Kāmàilóng de shíxí. Dì èr, wǒ huòdé le Zhōngguó zhèngfǔ jiǎngxuéjīn, qù Běijīng liúxué yī nián. Dì sān, wǒ búdàn huì shuō Zhōngwén, érqiě rènshi hěnduō hànzì. Wǒ de Zhōngguó tóngxué búdàn rèqíng, érqiě hěn lèyì bāng wǒ fùxí gōngkè. Suǒyǐ, xué Zhōngwén búdàn yǒuyòng, érqiě hěn yǒu yìsi! \\
\noindent
\textbf{Traduction :} Je m'appelle Paul, j'étudie à l'université à Yaoundé. Premièrement, apprendre le chinois m'a permis de trouver un stage chez Huawei Cameroun. Deuxièmement, j'ai obtenu une bourse du gouvernement chinois pour aller étudier à Pékin un an. Troisièmement, je ne sais pas seulement parler chinois, mais en plus je connais beaucoup de caractères. Mes camarades chinois ne sont pas seulement chaleureux, mais en plus ils sont très disposés à m'aider à réviser mes cours. Donc, apprendre le chinois n'est pas seulement utile, c'est en plus très intéressant !
\end{textebox}

\courssub{Consignes}
\noindent
\emph{Travail en groupe (4 élèves) — Durée : 45 min préparation + 15 min présentations + 10 min bilan.}

\begin{enumerate}
    \item \textbf{Analyse du document déclencheur (5 min) :} Repérez dans le témoignage de Paul les deux structures cibles. Combien d'arguments énumérés ? Combien de phrases avec \cle{不但……而且……} ? Y a-t-il une phrase avec deux sujets différents ? (Réponse attendue : 3 arguments \cle{第一/第二/第三} ; 3 phrases \cle{不但……而且……} dont une avec sujets différents : \zh{我的中国同学……}).
    
    \item \textbf{Brainstorming \& Sélection des arguments (10 min) :} En groupe, listez 5 raisons d'apprendre le chinois au Cameroun aujourd'hui. Sélectionnez les 3 meilleures. Associez à chaque raison un exemple concret (nom d'entreprise, ville, plat, fête, bourse).
    \begin{itemize}
        \item \textit{Pistes :} Marché central de Douala / commerçants chinois ; Bourses CSC / Institut Confucius Yaoundé ; Fête du Printemps à l'ambassade ; Projets routiers/barrages (ex: barrage de Nachtigal) ; Tourisme (Grande Muraille, pandas) ; Gastronomie (canard laqué vs ndolé).
    \end{itemize}
    
    \item \textbf{Rédaction collaborative du script oral (20 min) :} Rédigez un texte commun de \textbf{6 à 8 phrases}. Contraintes obligatoires :
    \begin{itemize}
        \item Utiliser \cle{第一……第二……第三……} pour énumérer les 3 arguments principaux (une phrase par argument ou phrase complexe).
        \item Inclure \textbf{au moins deux phrases} avec \cle{不但……而且……}.
        \item Parmi ces phrases, \textbf{au moins une} doit avoir \textbf{deux sujets différents} (ex: \zh{老师不但严厉，而且学生也很努力} — structure avancée : Sujet A + 不但……, Sujet B + 而且……).
        \item Vocabulaire HSK 1-3, phrases courtes, tons vérifiés.
    \end{itemize}
    
    \item \textbf{Répétition \& Distribution des rôles (10 min) :} Chaque élève prononce 1 à 2 phrases. Travaillez l'intonation (montée sur \zh{第一}、\zh{不但}, chute sur \zh{而且}、\zh{第三}). Vérifiez l'ordre des mots : \cle{Sujet + 不但 + V1, 而且 + V2}.
    
    \item \textbf{Présentation devant la classe (15 min) :} Chaque groupe présente. Les autres groupes remplissent la \textbf{Grille d'observation} (distribuée par l'enseignant) :
    \begin{center}
    \begin{tabularx}{\linewidth}{|X|X|X|X|X|}
    \hline
    \rowcolor{popblueL} Groupe & Nb \cle{第一/第二} corrects & Nb \cle{不但……而且……} corrects & Phrase à 2 sujets ? (Oui/Non) & Erreur(s) repérée(s) \\ \hline
    1 & & & & \\ \hline
    2 & & & & \\ \hline
    3 & & & & \\ \hline
    \end{tabularx}
    \end{center}
    
    \item \textbf{Vote \& Débat (5 min) :} La classe vote pour le groupe le plus convaincant (contenu + correction + aisance). L'enseignant corrige les erreurs majeures au tableau.
    
    \item \textbf{Trace écrite individuelle (Devoir / 20 min en classe) :} Chaque élève rédige \textbf{son propre paragraphe argumentatif} (80-100 caractères chinois) en reprenant les idées du groupe mais en reformulant personnellement. Contraintes identiques (énumération + 2× gradation dont 1 à 2 sujets). Remise à la fin de l'heure.
\end{enumerate}

\courssub{Ressources à mobiliser}
\noindent
\emph{Rappel synthétique pour l'élève (à garder sous les yeux pendant la préparation).}

\begin{propriete}[Structure 1 : Énumération ordonnée — \zh{第一……第二……(第三……)}]
\noindent
\textbf{Schéma :} \zh{第一，[Argument 1] 。第二，[Argument 2] 。(第三，[Argument 3] 。)}
\begin{itemize}
    \item \cle{第一} (dì yī) = « Premièrement / En premier » ; \cle{第二} (dì èr) = « Deuxièmement » ; \cle{第三} (dì sān) = « Troisièmement ».
    \item Elles fonctionnent comme \textbf{adverbiaux de séquence} placés \textbf{en tête de proposition} (souvent suivis d'une virgule ，).
    \item Pas de \cle{的} après \cle{第一/第二} ici (différent de \zh{第一个} « le premier [nom] »).
    \item Registre : oral et écrit, formel ou courant. Indispensable pour structurer un exposé, une dissertation, un argumentaire.
\end{itemize}
\begin{center}
\begin{tabularx}{\linewidth}{|X|X|X|}
\hline
\rowcolor{popblueL} Structure & Exemple (Caractères + Pinyin) & Traduction \\ \hline
\zh{第一，……} & \zh{第一，学中文能找好工作。} (Dì yī, xué Zhōngwén néng zhǎo hǎo gōngzuò.) & Premièrement, apprendre le chinois permet de trouver un bon travail. \\ \hline
\zh{第二，……} & \zh{第二，可以去中国旅游。} (Dì èr, kěyǐ qù Zhōngguó lǚyóu.) & Deuxièmement, on peut voyager en Chine. \\ \hline
\zh{第三，……} & \zh{第三，了解中国文化很有意思。} (Dì sān, liǎojiě Zhōngguó wénhuà hěn yǒu yìsi.) & Troisièmement, comprendre la culture chinoise est très intéressant. \\ \hline
\end{tabularx}
\end{center}
\end{propriete}

\begin{propriete}[Structure 2 : Gradation additive — \zh{不但……而且……}]
\noindent
\textbf{Schéma A (Sujet unique) :} \zh{Sujet + 不但 + V / Adj 1 ，而且 + V / Adj 2 。} \\
\textbf{Schéma B (Deux sujets différents — accent sur l'ampleur) :} \zh{Sujet A + 不但 + V / Adj 1 ，而且 + Sujet B + 也/还 + V / Adj 2 。}
\begin{itemize}
    \item \cle{不但} (búdàn) = « non seulement » ; \cle{而且} (érqiě) = « mais encore / en plus / de plus ».
    \item \textbf{Ordre strict :} \cle{不但} précède le premier élément, \cle{而且} précède le second. \textbf{Jamais} \zh{而且……不但……}.
    \item \cle{而且} attire souvent \cle{也} (yě) ou \cle{还} (hái) devant le V/Adj suivant pour renforcer l'addition.
    \item Négation : on nie généralement le second terme (\zh{不但不……而且不……}) ou on utilise \zh{既不……也不……}.
    \item \textbf{Piège francophone :} Ne pas confondre avec \zh{不但……但是……} (incorrect). \cle{而且} exprime l'addition, \cle{但是} (dànshì) l'opposition.
\end{itemize}
\begin{center}
\begin{tabularx}{\linewidth}{|X|X|X|}
\hline
\rowcolor{poporangeL} Structure & Exemple (Caractères + Pinyin) & Traduction \\ \hline
Sujet unique & \zh{中文不但有用，而且好玩。} (Zhōngwén búdàn yǒuyòng, érqiě hǎowán.) & Le chinois n'est pas seulement utile, il est en plus amusant. \\ \hline
Sujet unique + \zh{也} & \zh{他不但会说中文，而且也会写汉字。} (Tā búdàn huì shuō Zhōngwén, érqiě yě huì xiě hànzì.) & Il ne sait pas seulement parler chinois, mais en plus il sait écrire des caractères. \\ \hline
\cle{2 Sujets} & \zh{老师不但严厉，而且学生也很用功。} (Lǎoshī búdàn yánlì, érqiě xuésheng yě hěn yònggōng.) & Le prof n'est pas seulement strict, mais en plus les élèves sont très travailleurs. \\ \hline
\cle{2 Sujets} & \zh{中国企业不但投资喀麦隆，而且中国游客也越来越多。} (Zhōngguó qǐyè búdàn tóuzī Kāmàilóng, érqiě Zhōngguó yóukè yě yuè lái yuè duō.) & Les entreprises chinoises n'investissent pas seulement au Cameroun, mais en plus les touristes chinois sont de plus en plus nombreux. \\ \hline
\end{tabularx}
\end{center}
\end{propriete}

\begin{attention}[Erreurs fréquentes des francophones — Pièges à éviter]
\begin{enumerate}
    \item \textbf{Ordre des mots :} \emph{*« Je ne suis pas seulement fatigué, mais aussi affamé »} $\rightarrow$ \zh{我不但累，而且饿。} (Sujet \zh{我} une seule fois devant \zh{不但}). \textbf{Interdit :} \zh{*我不但累，我而且饿。}
    \item \textbf{Oubli de \cle{而且} :} \emph{*« Non seulement c'est cher »} $\rightarrow$ Phrase incomplète. \cle{不但} appelle \cle{而且}. \textbf{Correction :} \zh{这不但贵，而且不好吃。}
    \item \textbf{Confusion \cle{而且} / \cle{但是} :} \emph{*« Il est non seulement grand mais il est mince »} (opposition) $\rightarrow$ \zh{他不但高，但是瘦。} Si addition (grand ET mince) : \zh{他不但高，而且瘦。}
    \item \textbf{Place du sujet (2 sujets) :} \emph{*« Le prof est严厉, et les élèves sont aussi travailleurs »} $\rightarrow$ \zh{老师不但严厉，而且学生也很用功。} Le 2\textsuperscript{e} sujet (\zh{学生}) se place \textbf{après} \zh{而且}, souvent suivi de \zh{也/还}.
    \item \textbf{\cle{第一} vs \cle{第一个} :} \emph{*« Le premier, je veux dire… »} $\rightarrow$ \zh{第一，我想说……} (adverbe). \emph{« Le premier élève »} $\rightarrow$ \zh{第一个学生} (déterminant + classifieur).
\end{enumerate}
\end{attention}

\begin{definition}[Lexique utile pour le débat (HSK 1-3)]
\begin{center}
\begin{tabularx}{\linewidth}{|X|X|X|X|}
\hline
\rowcolor{popgreenL} Caractères & Pinyin & Nature & Traduction \\ \hline
\zh{机会} & jīhuì & n. & occasion, opportunité \\ \hline
\zh{奖学金} & jiǎngxuéjīn & n. & bourse d'études \\ \hline
\zh{贸易} & màoyì & n./v. & commerce, commercer \\ \hline
\zh{投资} & tóuzī & v./n. & investir, investissement \\ \hline
\zh{合作} & hézuò & v./n. & coopérer, coopération \\ \hline
\zh{文化} & wénhuà & n. & culture \\ \hline
\zh{交流} & jiāoliú & v./n. & échanger, échange \\ \hline
\zh{桥梁} & qiáoliáng & n. & pont (fig. lien) \\ \hline
\zh{发展} & fāzhǎn & v. & se développer \\ \hline
\zh{未来} & wèilái & n. & avenir, futur \\ \hline
\zh{热情} & rèqíng & adj. & chaleureux, enthousiaste \\ \hline
\zh{乐意} & lèyì & adj./v. & disposé à, volontaire \\ \hline
\zh{有用} & yǒuyòng & adj. & utile \\ \hline
\zh{有意思} & yǒu yìsi & adj. & intéressant, amusant \\ \hline
\zh{越来越多} & yuè lái yuè duō & loc. & de plus en plus nombreux \\ \hline
\end{tabularx}
\end{center}
\end{definition}

\courssub{Proposition de production et corrigé commenté}
\noindent
Voici une production modèle respectant toutes les contraintes (niveau HSK 3, contexte camerounais authentique).

\begin{exoresolu}[Production orale modèle — Groupe « Les Ambassadeurs du Dragon »]
\textbf{Énoncé :} Présenter 3 arguments avec \cle{第一/第二/第三} et 3 phrases \cle{不但……而且……} (dont 1 à 2 sujets) pour convaincre des élèves de Seconde de choisir le chinois LV2.

\tcblower

\textbf{Production chinoise (Texte intégral) :}
\zh{同学们好！欢迎来到中文角。} \\
\zh{第一，学中文不但能帮我们在雅温得和杜阿拉找到好工作，而且还能获得中国政府奖学金去北京、上海留学。} \\
\zh{第二，中国企业不但在喀麦隆投资修路、建桥、建大坝，而且中国商人也在中央市场、马卡市场做生意，会说中文沟通更方便。} \\
\zh{第三，中国文化不但有功夫、书法、熊猫，而且中餐也很好吃，比如北京烤鸭、四川火锅，我们喀麦隆的学生不但爱吃，而且也想学做菜。} \\
\zh{所以，选择中文，不但对未来有用，而且让生活更精彩！谢谢大家。}

\noindent
\textbf{Pinyin :}
Tóngxuémen hǎo! Huānyíng láidào Zhōngwén jiǎo. \\
Dì yī, xué Zhōngwén búdàn néng bāng wǒmen zài Yǎwēndé hé Dù'ālā zhǎodào hǎo gōngzuò, érqiě hái néng huòdé Zhōngguó zhèngfǔ jiǎngxuéjīn qù Běijīng, Shànghǎi liúxué. \\
Dì èr, Zhōngguó qǐyè búdàn zài Kāmàilóng tóuzī xiūlù, jiàn qiáo, jiàn dàbà, érqiě Zhōngguó shāngrén yě zài Zhōngyāng shìchǎng, Mǎkǎ shìchǎng zuò shēngyì, huì shuō Zhōngwén gōutōng gèng fāngbiàn. \\
Dì sān, Zhōngguó wénhuà búdàn yǒu gōngfu, shūfǎ, xióngmāo, érqiě Zhōngcān yě hěn hǎochī, bǐrú Běijīng kǎoyā, Sìchuān huǒguō, wǒmen Kāmàilóng de xuésheng búdàn ài chī, érqiě yě xiǎng xué zuò cài. \\
Suǒyǐ, xuǎnzé Zhōngwén, búdàn duì wèilái yǒuyòng, érqiě ràng shēnghuó gèng jīngcǎi! Xièxiè dàjiā.

\noindent
\textbf{Traduction française :}
Bonjour les camarades ! Bienvenue au Coin Chinois.
Premièrement, apprendre le chinois ne nous aide pas seulement à trouver du bon travail à Yaoundé et Douala, mais en plus cela permet d'obtenir des bourses du gouvernement chinois pour aller étudier à Pékin, Shanghai.
Deuxièmement, les entreprises chinoises n'investissent pas seulement au Cameroun pour construire des routes, des ponts, des barrages, mais en plus les commerçants chinois font aussi des affaires au Marché Central, au Marché Maka, parler chinois rend la communication plus facile.
Troisièmement, la culture chinoise n'a pas seulement le kung-fu, la calligraphie, les pandas, mais en plus la cuisine chinoise est très bonne, comme le canard laqué de Pékin, le fondue sichuanaise, nos élèves camerounais n'aiment pas seulement manger, mais veulent aussi apprendre à cuisiner.
Donc, choisir le chinois, n'est pas seulement utile pour l'avenir, mais en plus cela rend la vie plus merveilleuse ! Merci à tous.

\noindent
\textbf{Commentaire des choix de langue (Corrigé commenté) :}
\begin{enumerate}
    \item \textbf{Structure \cle{第一/第二/第三} :} Utilisées en tête de phrase pour structurer l'exposé en 3 parties logiques (Emploi/Études $\rightarrow$ Économie/Quotidien $\rightarrow$ Culture/Loisirs). Respect de la ponctuation \zh{第一，……} .
    \item \textbf{Phrase 1 (Argument 1) — \cle{不但……而且……} (Sujet unique \zh{学中文}) :} \zh{学中文不但能帮我们……找到好工作，而且还能获得……奖学金……}. Double verbe modal \zh{能} / \zh{还能}. \cle{而且} renforcé par \cle{还} (hái) « en plus / également ». Structure parfaite.
    \item \textbf{Phrase 2 (Argument 2) — \cle{不但……而且……} (\textbf{Deux sujets différents}) :} Sujet A = \zh{中国企业} (Entreprises chinoises) ; Sujet B = \zh{中国商人} (Commerçants chinois). \zh{中国企业不但在喀麦隆投资……，而且中国商人也在……做生意}. Respect du schéma B : \zh{而且 + Sujet B + 也 + Verbe}. Vocabulaire contextuel : \zh{修路、建桥、建大坝} (routes, ponts, barrages — ref. Nachtigal, Kribi), \zh{中央市场、马卡市场} (Marché Central Yaoundé, Marché Maka Douala).
    \item \textbf{Phrase 3 (Argument 3) — Première \cle{不但……而且……} (\textbf{Deux sujets} : \zh{中国文化} $\rightarrow$ \zh{中餐}) :} \zh{中国文化不但有功夫、书法、熊猫，而且中餐也很好吃}. Schéma B : le sujet change pour élargir le champ de la culture (arts $\rightarrow$ gastronomie). \cle{而且} + Sujet B (\zh{中餐}) + \cle{也}. Exemples culturels concrets : \zh{北京烤鸭、四川火锅}.
    \item \textbf{Phrase 3 (suite) — Deuxième \cle{不但……而且……} (Sujet unique \zh{我们喀麦隆的学生}) :} \zh{我们喀麦隆的学生不但爱吃，而且也想学做菜}. Verbes d'état/désir \zh{爱} / \zh{想}. \cle{也} devant \zh{想} pour marquer l'addition du désir à l'appréciation. \zh{学做菜} (apprendre à cuisiner) plus complet que \zh{学做}.
    \item \textbf{Conclusion :} \zh{所以……不但对未来有用，而且让生活更精彩}. Reprise synthétique des deux structures. \cle{让} (ràng) « rendre / faire en sorte que » + comparatif \cle{更} (gèng).
    \item \textbf{Niveau de langue :} HSK 3 (vocabulaire : \zh{投资、合作、沟通、奖学金、越来越多、精彩}), phrases complexes mais claires, tons standards.
\end{enumerate}
\end{exoresolu}

\courssub{Bilan de l'activité}
\noindent
Cette activité d'intégration a permis de vérifier l'acquisition opérationnelle des deux structures cibles dans une situation de communication \textbf{authentique, contextualisée (Cameroun/Chine) et évaluée par les pairs}.

\begin{aretenir}[Bilan synthétique pour le cahier de l'élève]
\begin{itemize}
    \item \cle{第一……第二……第三……} : Indispensable pour \textbf{structurer} un discours argumentatif, une narration chronologique ou une énumération logique. Place : \textbf{Début de proposition}.
    \item \cle{不但……而且……} : Indispensable pour \textbf{renforcer} un argument, montrer l'ampleur, l'accumulation d'avantages ou de faits. Place : \textbf{Entre deux prédicats (même sujet) ou entre deux propositions (sujets différents)}.
    \item \textbf{Association gagnante :} Utiliser \cle{第一/第二} pour poser les grands axes, puis \cle{不但……而且……} à l'intérieur de chaque axe pour développer la densité argumentative. C'est la marque d'un niveau HSK 3/4 (B1/B2).
    \item \textbf{Contexte Cameroun-Chine :} Les arguments « bourses CSC », « entreprises Huawei/China Road/China Railway », « marchés Douala/Yaoundé », « Institut Confucius Yaoundé/Buea », « gastronomie » sont des \emph{faits sociaux réels} qui ancrent la langue dans la vie de l'élève.
\end{itemize}

\end{aretenir}

\begin{methode}[Prolongement vers l'Examen (Épreuve écrite / Expression écrite)]
\noindent
\textbf{Sujet type BAC / Probatoire :} « \emph{Votre lycée jumelé avec un lycée de Pékin organise un échange. Rédigez une lettre de motivation (100-120 caractères) pour être sélectionné. Utilisez \cle{第一……第二……} et \cle{不但……而且……}.} »
\begin{enumerate}
    \item \textbf{Analyser} : Cible = jury de sélection. Ton = formel, motivé. Contraintes = 2 structures obligatoires + longueur.
    \item \textbf{Planifier} : 
    \begin{itemize}
        \item \zh{第一} : Niveau langue / résultats scolaires.
        \item \zh{第二} : Intérêt culture / adaptation.
        \item \zh{第三} : Rôle d'ambassadeur / avenir.
    \end{itemize}
    \item \textbf{Rédiger} (brouillon) en insérant 2 \cle{不但……而且……} (ex: \zh{我不但成绩好，而且很认真} ; \zh{老师不但推荐我，而且父母也支持}).
    \item \textbf{Relire} : Ordre des mots ? \cle{而且} présent ? Sujets clairs ? Ponctuation chinoise ? Caractères lisibles ?
    \item \textbf{Recopier} proprement sur la copie d'examen.
\end{enumerate}
\end{methode}

\begin{experience}[Prolongement numérique — OnBuch+ / Classe inversée]
\noindent
Sur l'application \textbf{OnBuch+}, module « ZHA02 - Intégration » :
\begin{itemize}
    \item \textbf{Exercice interactif 1 :} Glisser-déposer pour reconstruire 5 phrases \cle{不但……而且……} (dont 2 à 2 sujets) à partir de blocs de mots.
    \item \textbf{Exercice interactif 2 :} Dictée à trous (pinyin $\rightarrow$ caractères) sur le témoignage de Paul (Document déclencheur).
    \item \textbf{Atelier Oral (IA) :} L'élève s'enregistre en présentant son paragraphe individuel. L'IA (reconnaissance vocale + LLM) donne un retour immédiat sur : \textit{justesse des tons sur \cle{不但/búdàn} et \cle{而且/érqiě}, fluidité de l'énumération \cle{第一/第二}, respect de l'ordre des mots}.
    \item \textbf{Badge « Argumentateur Sinophone » :} Décerné si score $\ge$ 80\% sur les 3 activités.
\end{itemize}
\end{experience}

\begin{savaistu}[Regard croisé Cameroun — Chine : L'apprentissage des langues]
\noindent
Au Cameroun, pays officiellement bilingue (Français/Anglais), l'apprentissage d'une \textbf{troisième langue} (LV2/LV3) comme le chinois, l'allemand, l'espagnol ou l'arabe est un atout majeur.
\begin{itemize}
    \item \textbf{Chine :} Depuis 2018, le chinois est intégré au programme officiel MINESEC. Les \emph{Instituts Confucius} de Yaoundé (Université de Yaoundé II) et Buea (Université de Buea) sont des hubs de certification HSK, de bourses (CSC, Provincial) et d'échanges universitaires.
    \item \textbf{Économie :} La Chine est le 1\textsuperscript{er} partenaire commercial du Cameroun (importations machines, textiles, télécoms ; exportations bois, coton, pétrole). Parler chinois est un \textbf{levier d'employabilité} direct (interprète, commerce, logistique, BTP, télécoms Huawei/ZTE/Transsion).
    \item \textbf{Culture :} La \emph{Fête du Printemps} (Nouvel An chinois) est célébrée officiellement à l'Ambassade de Chine et dans les lycées partenaires. Le \emph{thé}, la \emph{calligraphie}, le \emph{tai-chi} et la \emph{cuisine} (raviolis \zh{饺子}, nouilles \zh{拉面}) séduisent la jeunesse camerounaise.
    \item \textbf{Parallèle :} Tout comme le Camerounais apprend l'anglais/français pour l'unité nationale et l'ouverture internationale, le Chinois apprend l'anglais (obligatoire dès le primaire) pour l'ouverture mondiale. Apprendre le chinois, c'est \textbf{devenir un acteur du lien Cameroun-Chine}.
\end{itemize}
\end{savaistu}

\newpage

\coursec{Méthodes et exercices résolus}

\courssub{Observation guidée : découvrir les deux structures}

\begin{textebox}[Dialogue : pourquoi apprendre le chinois ?]
\zh{A：你为什么学习汉语？}\\
\emph{A : Nǐ wèishénme xuéxí Hànyǔ ?}\\
« A : Pourquoi apprends-tu le chinois ? »\\
\zh{B：第一，汉语很有用；第二，我想以后去中国。}\\
\emph{B : Dì-yī, Hànyǔ hěn yǒuyòng; dì-èr, wǒ xiǎng yǐhòu qù Zhōngguó.}\\
« B : Premièrement, le chinois est très utile ; deuxièmement, je veux aller en Chine plus tard. »\\
\zh{A：汉语难吗？}\\
\emph{A : Hànyǔ nán ma ?}\\
« A : Le chinois, c'est difficile ? »\\
\zh{B：汉语不但有意思，而且很有用，我越学越喜欢。}\\
\emph{B : Hànyǔ bùdàn yǒu yìsi, érqiě hěn yǒuyòng, wǒ yuè xué yuè xǐhuan.}\\
« B : Le chinois est non seulement intéressant, mais aussi très utile ; plus je l'étudie, plus je l'aime. »
\end{textebox}

\begin{definition}[Deux mots-outils pour organiser le discours]
\cle{第一……第二……} (dì-yī…… dì-èr……) : corrélatifs d'\textbf{énumération}, « premièrement… deuxièmement… ». Ils ordonnent des arguments, des raisons, des étapes.\\
\cle{不但……而且……} (bùdàn…… érqiě……) : corrélatifs de \textbf{gradation}, « non seulement… mais aussi… ». Ils relient deux informations dont la seconde ajoute quelque chose à la première.
\end{definition}

\begin{propriete}[Schémas des deux structures]
Énumération : \cle{第一，proposition complète；第二，proposition complète。} (on peut continuer avec \zh{第三} « troisièmement » et \zh{最后} « enfin »).\\
Gradation, même sujet : \cle{Sujet + 不但 + A，而且 + B。}\\
Gradation, sujets différents : \cle{不但 + Sujet 1 + A，而且 + Sujet 2 + 也 + B。}
\end{propriete}

\begin{center}
\begin{tabularx}{\linewidth}{|X|X|X|}
\hline
\rowcolor{popblueL} Français & Chinois (caractères + pinyin) & Fonction \\
\hline
premièrement… deuxièmement… & \zh{第一……第二……} dì-yī…… dì-èr…… & énumérer des arguments \\
\hline
non seulement… mais aussi… & \zh{不但……而且……} bùdàn…… érqiě…… & gradation (ajout renforcé) \\
\hline
donc, c'est pourquoi & \zh{所以} suǒyǐ & conclure après une énumération \\
\hline
aussi (dans la 2\up{e} proposition) & \zh{也} yě & renforce la gradation avec deux sujets \\
\hline
\end{tabularx}
\end{center}

\courssub{Méthode 1 : Construire une énumération avec 第一……第二……}

\begin{methode}[Énumérer des arguments avec 第一……第二……]
Pour organiser un raisonnement (exposé, rédaction, réponse à une question « pourquoi ? »), suis ces étapes :
\begin{enumerate}
\item \textbf{Annonce le sujet} avec une phrase simple : Sujet + Verbe + Complément.\\ \zh{我喜欢汉语。} \emph{Wǒ xǐhuan Hànyǔ.} « J'aime le chinois. »
\item \textbf{Introduis les raisons} par \zh{因为} (parce que) ou directement par \zh{第一} (premièrement).
\item \textbf{Place 第一 / 第二 / 第三 en tête de proposition}, suivi d'une virgule chinoise \zh{，}, puis d'une phrase complète (sujet + verbe).\\ Schéma : \cle{第一，phrase complète；第二，phrase complète。}
\item \textbf{Conclus} éventuellement avec \zh{所以} (donc) : \zh{所以……} « donc… ».
\item \textbf{Vérifie} : chaque proposition après 第一 / 第二 doit pouvoir exister seule comme phrase correcte.
\end{enumerate}
Réflexe : \zh{第一} = « premièrement », \zh{第二} = « deuxièmement ». On peut continuer : \zh{第三} (troisièmement), \zh{最后} (enfin).
\end{methode}

\begin{exoresolu}[Rédaction guidée : pourquoi j'apprends le chinois]
\textbf{Énoncé (type examen).} Réponds à la question \zh{你为什么学习汉语？} (Pourquoi apprends-tu le chinois ?) en rédigeant un court paragraphe de 3 à 4 phrases utilisant obligatoirement \zh{第一……第二……}. Donne le pinyin et la traduction.
\tcblower
\textbf{Étape 1 — Analyse de la consigne.} La question demande des \emph{raisons} : il faut donc une énumération ordonnée. La structure 第一……第二…… est parfaite pour cela.

\textbf{Étape 2 — Choix des arguments.} Argument 1 : le chinois est utile (\zh{汉语很有用}). Argument 2 : je veux travailler avec la Chine (\zh{我想跟中国人一起工作}).

\textbf{Étape 3 — Construction.} On annonce d'abord le sujet, puis on énumère avec 第一 / 第二 placés en tête de proposition, chacun suivi de la virgule chinoise \zh{，}. Chaque proposition est complète (sujet + verbe).

\textbf{Réponse rédigée :}\\
\zh{我很喜欢学习汉语。第一，汉语很有用，很多中国人在喀麦隆工作。第二，我想以后跟中国人一起做生意。所以，我每天都努力学习汉语。}\\
\emph{Wǒ hěn xǐhuan xuéxí Hànyǔ. Dì-yī, Hànyǔ hěn yǒuyòng, hěn duō Zhōngguó rén zài Kāmàilóng gōngzuò. Dì-èr, wǒ xiǎng yǐhòu gēn Zhōngguó rén yìqǐ zuò shēngyi. Suǒyǐ, wǒ měi tiān dōu nǔlì xuéxí Hànyǔ.}\\
« J'aime beaucoup apprendre le chinois. Premièrement, le chinois est très utile, beaucoup de Chinois travaillent au Cameroun. Deuxièmement, je veux plus tard faire du commerce avec des Chinois. C'est pourquoi j'étudie le chinois avec effort chaque jour. »

\textbf{Étape 4 — Vérification.} 第一 et 第二 sont bien en tête de proposition ; chaque proposition a un sujet et un verbe ; la conclusion avec \zh{所以} est correcte. Note : on écrit \zh{第一} avec le chiffre \zh{一} et non \zh{弟} (petit frère) — homophone piégeux.
\end{exoresolu}

\courssub{Méthode 2 : Construire une gradation avec 不但……而且……}

\begin{methode}[Exprimer « non seulement… mais aussi… » avec 不但……而且……]
\begin{enumerate}
\item \textbf{Identifie les deux informations} à relier : une qualité A, une qualité B qui s'y ajoute (gradation).
\item \textbf{Cas 1 — même sujet pour les deux propositions} : le sujet se place \textbf{avant} \zh{不但}.\\ Schéma : \cle{Sujet + 不但 + A，而且 + B。}\\ \zh{他不但会说汉语，而且会说英语。} \emph{Tā bùdàn huì shuō Hànyǔ, érqiě huì shuō Yīngyǔ.} « Non seulement il parle chinois, mais aussi il parle anglais. »
\item \textbf{Cas 2 — deux sujets différents} : chaque sujet se place \textbf{après} 不但 / 而且.\\ Schéma : \cle{不但 + Sujet 1 + A，而且 + Sujet 2 + 也 + B。}\\ \zh{不但我喜欢足球，而且我弟弟也喜欢。} \emph{Bùdàn wǒ xǐhuan zúqiú, érqiě wǒ dìdi yě xǐhuan.} « Non seulement moi j'aime le football, mais mon petit frère aussi. »
\item \textbf{Ne jamais oublier le deuxième élément} : \zh{不但} seul est incorrect ; il faut toujours \zh{而且} (avec éventuellement \zh{也} ou \zh{还} dans la seconde proposition).
\item \textbf{Traduis mentalement} : 不但 = « non seulement », 而且 = « mais aussi / et de plus ».
\end{enumerate}
\end{methode}

\begin{exoresolu}[Transformation : de deux phrases simples à une phrase à gradation]
\textbf{Énoncé (type examen).} Transforme les paires de phrases suivantes en une seule phrase avec \zh{不但……而且……}. Donne le pinyin et la traduction.\\
1. \zh{她会唱歌。她会跳舞。} (Elle sait chanter. Elle sait danser.)\\
2. \zh{我妈妈喜欢这个城市。我爸爸喜欢这个城市。} (Ma mère aime cette ville. Mon père aime cette ville.)
\tcblower
\textbf{Phrase 1 — raisonnement.} Les deux propositions ont le \textbf{même sujet} (\zh{她}). D'après la règle (cas 1), le sujet se place \emph{avant} \zh{不但} et n'est pas répété dans la seconde proposition.

\textbf{Réponse 1 :} \zh{她不但会唱歌，而且会跳舞。}\\
\emph{Tā bùdàn huì chàng gē, érqiě huì tiào wǔ.}\\
« Non seulement elle sait chanter, mais elle sait aussi danser. »

\textbf{Phrase 2 — raisonnement.} Les deux propositions ont des \textbf{sujets différents} (\zh{妈妈} et \zh{爸爸}). D'après la règle (cas 2), chaque sujet se place \emph{après} 不但 / 而且. On ajoute \zh{也} (aussi) dans la seconde proposition pour renforcer la gradation.

\textbf{Réponse 2 :} \zh{不但我妈妈喜欢这个城市，而且我爸爸也喜欢。}\\
\emph{Bùdàn wǒ māma xǐhuan zhège chéngshì, érqiě wǒ bàba yě xǐhuan.}\\
« Non seulement ma mère aime cette ville, mais mon père l'aime aussi. »

\textbf{Conclusion.} Le choix de la place du sujet dépend d'une seule question : « les deux propositions ont-elles le même sujet ? » Même sujet → sujet avant 不但. Sujets différents → sujets après 不但 et 而且.
\end{exoresolu}

\courssub{Méthode 3 : Traduire du français vers le chinois sans calque}

\begin{methode}[Traduire « premièrement… deuxièmement… » et « non seulement… mais aussi… »]
\begin{enumerate}
\item \textbf{Repère la structure française} dans la phrase à traduire (énumération ou gradation).
\item \textbf{Choisis la structure chinoise correspondante} : 第一……第二…… pour l'énumération ; 不但……而且…… pour la gradation.
\item \textbf{Respecte l'ordre chinois} : Sujet + Temps/Lieu + Verbe + Complément. Ne calque jamais l'ordre français : « il parle bien le chinois » → \zh{他汉语说得很好} (complément de manière avec \zh{得}) ; jamais de traduction mot à mot.
\item \textbf{Place les corrélatifs} : 第一 / 第二 en tête de proposition ; 不但 après le sujet unique, avant les sujets différents.
\item \textbf{Relis en chinois} (pas en français) : la phrase doit être naturelle pour un Sinophone.
\end{enumerate}
\end{methode}

\begin{exoresolu}[Thème : traduction guidée]
\textbf{Énoncé (type examen).} Traduis en chinois : « Premièrement, Douala est une grande ville ; deuxièmement, il y a beaucoup de travail. Non seulement mon frère aîné veut y habiter, mais ma sœur aussi. » Donne le pinyin.
\tcblower
\textbf{Étape 1 — Repérage.} « Premièrement… deuxièmement… » → \zh{第一……第二……}. « Non seulement… mais aussi » avec deux sujets différents (frère / sœur) → \zh{不但 + sujet 1……而且 + sujet 2 + 也……}.

\textbf{Étape 2 — Traduction des éléments.} Douala = \zh{杜阿拉} (Dù'ālā) ; « grande ville » = \zh{大城市} ; « il y a beaucoup de travail » = \zh{有很多工作} ; « y habiter » = \zh{住在那儿} ; « mon frère aîné » = \zh{我哥哥} ; « ma sœur (aînée) » = \zh{我姐姐}.

\textbf{Étape 3 — Assemblage.}\\
\zh{第一，杜阿拉是个大城市；第二，那儿有很多工作。不但我哥哥想住在那儿，而且我姐姐也想住在那儿。}\\
\emph{Dì-yī, Dù'ālā shì ge dà chéngshì; dì-èr, nàr yǒu hěn duō gōngzuò. Bùdàn wǒ gēge xiǎng zhù zài nàr, érqiě wǒ jiějie yě xiǎng zhù zài nàr.}

\textbf{Étape 4 — Vérification.} Le classificateur \zh{个} accompagne bien \zh{城市} ; \zh{住在} + lieu est la structure correcte pour « habiter à » ; les sujets différents sont bien placés après 不但 et 而且, avec \zh{也} dans la seconde proposition. Traduction fidèle, sans calque de l'ordre français.
\end{exoresolu}

\courssub{Méthode 4 : Rédiger un court texte argumenté (exposé ou expression écrite)}

\begin{methode}[Combiner les deux structures dans un paragraphe cohérent]
Pour un sujet du type « Parle de ta famille / de ton école / de ta ville » :
\begin{enumerate}
\item \textbf{Phrase d'introduction} : présente le thème en une phrase simple.
\item \textbf{Développement avec 第一……第二……} : donne deux ou trois idées ordonnées.
\item \textbf{Enrichis une idée avec 不但……而且……} : ajoute une gradation pour montrer ta maîtrise des structures.\\ Exemple : \zh{我们学校不但很大，而且很干净。} \emph{Wǒmen xuéxiào bùdàn hěn dà, érqiě hěn gānjìng.} « Notre école est non seulement très grande, mais aussi très propre. »
\item \textbf{Conclusion} avec \zh{所以} ou une phrase d'opinion : \zh{我觉得……} « je pense que… ».
\item \textbf{Relecture ciblée} : vérifie les tons dans le pinyin, la place des sujets autour de 不但/而且, et la virgule chinoise après 第一 / 第二.
\end{enumerate}
\end{methode}

\begin{exoresolu}[Expression écrite : ma famille]
\textbf{Énoncé (type examen).} Rédige un court texte (5 à 6 phrases) intitulé \zh{我的家} (Ma famille), en utilisant au moins une fois \zh{第一……第二……} et une fois \zh{不但……而且……}. Donne le pinyin et la traduction.
\tcblower
\textbf{Étape 1 — Plan.} Introduction (ma famille) → énumération de deux raisons d'aimer ma famille (第一/第二) → gradation sur mon père (不但/而且) → conclusion (所以).

\textbf{Étape 2 — Rédaction phrase par phrase.}\\
Introduction : \zh{我有一个很幸福的家。} « J'ai une famille très heureuse. »\\
Énumération : \zh{第一，我的爸爸妈妈都很爱我；第二，我们每天晚上一起吃饭、说话。}\\
Gradation (même sujet \zh{爸爸} → sujet avant 不但) : \zh{我爸爸不但工作很努力，而且常常帮助我学习。}\\
Conclusion : \zh{所以，我很爱我的家。}

\textbf{Texte complet :}\\
\zh{我的家}\\\zh{我有一个很幸福的家。第一，我的爸爸妈妈都很爱我；第二，我们每天晚上一起吃饭、说话。我爸爸不但工作很努力，而且常常帮助我学习。所以，我很爱我的家。}\\
\emph{Wǒ de jiā}\\\emph{Wǒ yǒu yí ge hěn xìngfú de jiā. Dì-yī, wǒ de bàba māma dōu hěn ài wǒ; dì-èr, wǒmen měi tiān wǎnshang yìqǐ chī fàn, shuō huà. Wǒ bàba bùdàn gōngzuò hěn nǔlì, érqiě chángcháng bāngzhù wǒ xuéxí. Suǒyǐ, wǒ hěn ài wǒ de jiā.}\\
« Ma famille — J'ai une famille très heureuse. Premièrement, mon père et ma mère m'aiment beaucoup ; deuxièmement, nous dînons et parlons ensemble chaque soir. Mon père non seulement travaille avec beaucoup d'effort, mais il m'aide aussi souvent à étudier. C'est pourquoi j'aime beaucoup ma famille. »

\textbf{Étape 3 — Vérification.} Les deux structures imposées sont présentes et correctes ; le sujet unique \zh{爸爸} est bien placé avant \zh{不但} ; dans \zh{工作很努力}, le verbe \zh{工作} est suivi du complément de manière introduit par l'adverbe \zh{很}, structure correcte sans répéter le verbe ; les virgules chinoises sont bien placées après 第一 et 第二.
\end{exoresolu}

\courssub{Méthode 5 : Identifier et corriger une phrase fautive}

\begin{methode}[Déboguer une phrase avec 第一/第二 ou 不但/而且]
\begin{enumerate}
\item \textbf{Repère la structure} utilisée dans la phrase.
\item \textbf{Teste la complétude} : 不但 est-il bien suivi de 而且 (ou 也/还) ? Sinon, la phrase est incomplète.
\item \textbf{Teste la place du sujet} : un seul sujet → avant 不但 ; deux sujets → après 不但 et 而且.
\item \textbf{Teste l'ordre interne} : chaque proposition doit suivre Sujet + (Temps/Lieu) + Verbe + Complément.
\item \textbf{Corrige puis traduis} la phrase corrigée pour vérifier le sens.
\end{enumerate}
\end{methode}

\begin{exoresolu}[Correction de phrases fautives]
\textbf{Énoncé (type examen).} Les phrases suivantes contiennent chacune une erreur. Corrige-les et justifie. Donne le pinyin et la traduction des phrases corrigées.\\
1. \zh{不但她会说汉语，会说法语。}\\
2. \zh{他不但喜欢足球，而且他。}\\
3. \zh{第一我喜欢音乐，第二我喜欢运动，所以我每天都很高兴。} (phrase à améliorer)
\tcblower
\textbf{Phrase 1 — diagnostic.} \zh{不但} est présent mais \zh{而且} manque : la gradation est incomplète. Les deux propositions ont le même sujet (\zh{她}), donc le sujet doit passer \emph{avant} \zh{不但}.\\
\textbf{Correction :} \zh{她不但会说汉语，而且会说法语。}\\
\emph{Tā bùdàn huì shuō Hànyǔ, érqiě huì shuō Fǎyǔ.} « Non seulement elle parle chinois, mais elle parle aussi français. »

\textbf{Phrase 2 — diagnostic.} La seconde proposition \zh{而且他} est vide : il manque le verbe. Chaque proposition après 不但/而且 doit être complète.\\
\textbf{Correction :} \zh{他不但喜欢足球，而且喜欢篮球。}\\
\emph{Tā bùdàn xǐhuan zúqiú, érqiě xǐhuan lánqiú.} « Non seulement il aime le football, mais il aime aussi le basket-ball. »

\textbf{Phrase 3 — diagnostic.} La phrase est grammaticalement possible mais lourde : le sujet \zh{我} est répété inutilement et il manque les virgules chinoises après 第一 et 第二. En chinois, on évite les répétitions quand le sujet est identique.\\
\textbf{Version améliorée :} \zh{第一，我喜欢音乐；第二，我喜欢运动，所以我每天都很高兴。}\\
\emph{Dì-yī, wǒ xǐhuan yīnyuè; dì-èr, wǒ xǐhuan yùndòng, suǒyǐ wǒ měi tiān dōu hěn gāoxìng.} « Premièrement, j'aime la musique ; deuxièmement, j'aime le sport, donc je suis content chaque jour. » (Les virgules chinoises après 第一 et 第二 sont obligatoires à l'écrit.)

\textbf{Conclusion.} Trois réflexes de correction : compléter la paire 不但/而且, vérifier que chaque proposition a un verbe, alléger les répétitions de sujet.
\end{exoresolu}

\begin{attention}[Les 5 erreurs qui coûtent des points à l'examen]
\begin{enumerate}
\item \textbf{Oublier la deuxième moitié de la gradation.} Écrire \zh{他不但会说汉语。} sans \zh{而且} est une phrase incomplète : la paire \cle{不但……而且……} est inséparable. Calque dangereux du français « non seulement » qui peut seul paraître suffisant.
\item \textbf{Mal placer le sujet autour de 不但.} Même sujet → \zh{她不但……而且……} (sujet avant). Sujets différents → \zh{不但我哥哥……而且我姐姐也……} (sujets après). Inverser les deux cas est la faute la plus fréquente des francophones.
\item \textbf{Confondre les caractères 第 et 弟.} \zh{第一} (dì-yī, « premièrement ») s'écrit avec \zh{第} (clé du bambou \zh{⺮} en haut) ; \zh{弟} (dì) signifie « petit frère ». Homophones, sens totalement différents : une faute de caractère coûte des points en écriture.
\item \textbf{Oublier la virgule chinoise après 第一 / 第二.} À l'écrit, on écrit \zh{第一，……；第二，……} avec la ponctuation pleine chasse \zh{，} et le point-virgule chinois \zh{；} entre les grandes étapes de l'énumération.
\item \textbf{Calquer l'ordre des mots du français dans les propositions.} « Il parle bien le chinois » ne se traduit pas mot à mot : on écrit \zh{他汉语说得很好} (avec le mot-outil \zh{得} pour le complément de manière). De même, le lieu et le temps se placent avant le verbe : \zh{我们每天晚上一起吃饭} et non l'ordre français « nous mangeons ensemble chaque soir » calqué.
\end{enumerate}
\end{attention}

\begin{aretenir}[Synthèse de la leçon]
\begin{itemize}
\item \cle{第一……第二……} : énumération ordonnée ; corrélatifs en tête de proposition, suivis de la virgule chinoise ; chaque proposition est complète ; conclusion possible avec \zh{所以}.
\item \cle{不但……而且……} : gradation « non seulement… mais aussi… » ; la paire est inséparable.
\item Place du sujet avec 不但/而且 : \textbf{même sujet} → sujet avant 不但 ; \textbf{sujets différents} → sujets après 不但 et 而且, avec \zh{也} dans la seconde proposition.
\item Attention aux homophones \zh{第} / \zh{弟} et à la ponctuation chinoise pleine chasse.
\end{itemize}
\end{aretenir}

\begin{exercice}[Pour t'entraîner]
1. Transforme en une phrase avec \zh{不但……而且……} : \zh{他会说汉语。他会写汉字。} (même sujet)\\
2. Transforme en une phrase avec \zh{不但……而且……} : \zh{我喜欢喀麦隆。我哥哥喜欢喀麦隆。} (sujets différents)\\
3. Traduis en chinois : « Premièrement, j'aime le chinois ; deuxièmement, je veux travailler en Chine. »\\
4. Corrige : \zh{他不但会唱歌，会跳舞。}
\end{exercice}

\begin{corrige}[Exercice « Pour t'entraîner »]
1. Même sujet → sujet avant 不但 : \zh{他不但会说汉语，而且会写汉字。} \emph{Tā bùdàn huì shuō Hànyǔ, érqiě huì xiě Hànzì.} « Non seulement il parle chinois, mais il sait aussi écrire les caractères chinois. »\\
2. Sujets différents → sujets après 不但 et 而且, avec \zh{也} : \zh{不但我喜欢喀麦隆，而且我哥哥也喜欢。} \emph{Bùdàn wǒ xǐhuan Kāmàilóng, érqiě wǒ gēge yě xǐhuan.} « Non seulement j'aime le Cameroun, mais mon grand frère l'aime aussi. »\\
3. \zh{第一，我喜欢汉语；第二，我想在中国工作。} \emph{Dì-yī, wǒ xǐhuan Hànyǔ; dì-èr, wǒ xiǎng zài Zhōngguó gōngzuò.}\\
4. Il manque \zh{而且} ; même sujet → sujet avant 不但 : \zh{他不但会唱歌，而且会跳舞。} \emph{Tā bùdàn huì chàng gē, érqiě huì tiào wǔ.} « Non seulement il sait chanter, mais il sait aussi danser. »
\end{corrige}

\newpage

\coursec{Exercices}

\courssub{Je vérifie mes connaissances}

\begin{exercice}[QCM : la bonne structure]
Entoure la bonne réponse.
\begin{enumerate}
\item \zh{他（ ）会说英语，（ ）会说汉语。} \\
A. 不但……而且…… \quad B. 第一……第二…… \quad C. 因为……所以……
\item Pour énumérer des arguments dans un exposé, on utilise : \\
A. \zh{不但……而且……} \quad B. \zh{第一……第二……} \quad C. \zh{虽然……但是……}
\item \zh{第一} signifie : \\
A. premièrement \quad B. le premier jour \quad C. une fois
\item Dans \zh{不但……而且……}, quand les deux parties ont le même sujet, le sujet se place : \\
A. après \zh{不但} \quad B. avant \zh{不但} \quad C. après \zh{而且}
\end{enumerate}
\end{exercice}

\begin{exercice}[Vrai ou faux ? Justifie ta réponse]
Dis si la phrase est correcte (V) ou fautive (F) et justifie en une phrase.
\begin{enumerate}
\item \zh{不但我喜欢足球，而且我喜欢篮球。} (même sujet)
\item \zh{第二，我们学习汉语。}
\item \zh{他不但会说汉语，而且说得很好。}
\item \zh{第两，喀麦隆有很多学校。}
\end{enumerate}
\end{exercice}

\begin{exercice}[Vocabulaire : complète le tableau]
Complète le tableau avec le pinyin et la traduction française.

\begin{center}
\begin{tabularx}{\linewidth}{|X|X|X|X|}
\hline
\rowcolor{popblueL} Caractères & Pinyin & Nature & Traduction \\
\hline
\zh{第一} & & & \\
\hline
\zh{不但} & & & \\
\hline
\zh{而且} & & & \\
\hline
\zh{也} & & & \\
\hline
\zh{优点} & & & \\
\hline
\zh{方面} & & & \\
\hline
\end{tabularx}
\end{center}
\end{exercice}

\begin{exercice}[Associe le pinyin aux caractères]
Relie chaque caractère ou mot à son pinyin.
\begin{enumerate}
\item \zh{第二} \hspace{1cm} a. érqiě
\item \zh{而且} \hspace{1cm} b. dìsān
\item \zh{不但} \hspace{1cm} c. dì'èr
\item \zh{第三} \hspace{1cm} d. bùdàn
\end{enumerate}
\end{exercice}

\courssub{Je m'entraîne}

\begin{exercice}[Traduis en français]
Traduis les phrases suivantes en français.
\begin{enumerate}
\item \zh{第一，汉语很有意思；第二，汉语很有用。} \\
\emph{Dìyī, Hànyǔ hěn yǒu yìsi; dì'èr, Hànyǔ hěn yǒuyòng.}
\item \zh{他不但学习好，而且体育也很好。} \\
\emph{Tā bùdàn xuéxí hǎo, érqiě tǐyù yě hěn hǎo.}
\item \zh{不但我喜欢喀麦隆，我的朋友们也喜欢。} \\
\emph{Bùdàn wǒ xǐhuan Kāmàilóng, wǒ de péngyoumen yě xǐhuan.}
\end{enumerate}
\end{exercice}

\begin{exercice}[Transformation : fusionne les phrases]
Fusionne chaque paire de phrases en une seule phrase avec \zh{不但……而且……}. Attention à la place du sujet !
\begin{enumerate}
\item \zh{他会说法语。他会说英语。}
\item \zh{雅温得很大。雅温得很美。}
\item \zh{我喜欢足球。我弟弟也喜欢足球。}
\item \zh{她学习汉语。她学习中国文化。}
\item \zh{这个菜好吃。这个菜不贵。}
\end{enumerate}
\end{exercice}

\begin{exercice}[Texte à trous]
Complète le texte avec les mots suivants : \zh{第一、第二、不但、而且、也}.

\begin{textebox}[Mon école]
我的学校很好。（ ），老师很好；（ ），同学很友好。\\
\emph{Wǒ de xuéxiào hěn hǎo. ( ), lǎoshī hěn hǎo; ( ), tóngxué hěn yǒuhǎo.}\\
我们学校（ ）有很大的图书馆，（ ）有漂亮的操场。\\
\emph{Wǒmen xuéxiào ( ) yǒu hěn dà de túshūguǎn, ( ) yǒu piàoliang de cāochǎng.}\\
我很爱学习，我的朋友们（ ）很爱学习。\\
\emph{Wǒ hěn ài xuéxí, wǒ de péngyoumen ( ) hěn ài xuéxí.}
\end{textebox}
\end{exercice}

\begin{exercice}[Remets les mots dans l'ordre]
Remets les mots dans l'ordre pour former une phrase correcte.
\begin{enumerate}
\item 不但 / 汉语 / 我 / 学习 / 而且 / 学习 / 英语 / 也
\item 第二 / 杜阿拉 / 是 / 大城市 / 一个
\item 不但 / 他 / 而且 / 足球 / 喜欢 / 篮球 / 喜欢
\item 第一 / 第二 / 喀麦隆 / 很美 / 人很友好 / ； / ，
\end{enumerate}
\end{exercice}

\courssub{Je me prépare au Bac}

\begin{exercice}[Type Bac — Compréhension de texte et questions de langue]
Lis le texte suivant, puis réponds aux questions.

\begin{textebox}[Ma ville : Douala]
我住在杜阿拉。杜阿拉是喀麦隆的大城市。第一，杜阿拉很大，人很多；第二，杜阿拉有港口，经济很发达；第三，这里的人很友好。\\
\emph{Wǒ zhù zài Dù'ālā. Dù'ālā shì Kāmàilóng de dà chéngshì. Dìyī, Dù'ālā hěn dà, rén hěn duō; dì'èr, Dù'ālā yǒu gǎngkǒu, jīngjì hěn fādá; dìsān, zhèlǐ de rén hěn yǒuhǎo.}\\
杜阿拉不但经济发达，而且文化也很丰富。不但年轻人喜欢这个城市，老年人也喜欢。\\
\emph{Dù'ālā bùdàn jīngjì fādá, érqiě wénhuà yě hěn fēngfù. Bùdàn niánqīngrén xǐhuan zhège chéngshì, lǎoniánrén yě xǐhuan.}
\end{textebox}

\textbf{Questions de compréhension} (réponds en français) :
\begin{enumerate}
\item Où habite l'auteur du texte ?
\item Cite les trois avantages de Douala mentionnés dans le texte.
\item Qui aime cette ville, selon la dernière phrase ?
\end{enumerate}

\textbf{Questions de langue} :
\begin{enumerate}
\item Releve dans le texte une phrase avec \zh{不但……而且……} où les deux parties ont le même sujet, et une phrase où les sujets sont différents.
\item Pourquoi le mot \zh{也} apparaît-il dans la dernière phrase ? Explique la règle.
\item Transforme : \zh{杜阿拉很大。杜阿拉人很多。} en une phrase avec \zh{不但……而且……}.
\end{enumerate}
\end{exercice}

\begin{exercice}[Type Bac — Thème et version]
\textbf{A. Thème.} Traduis les phrases suivantes en chinois (caractères obligatoires) :
\begin{enumerate}
\item Premièrement, le chinois est intéressant ; deuxièmement, il est très utile.
\item Mon frère aime non seulement le football, mais aussi le basketball.
\item Non seulement j'aime Yaoundé, mais mes parents l'aiment aussi. (Attention : sujets différents !)
\item Cette école est non seulement grande, mais aussi très belle.
\end{enumerate}

\textbf{B. Version.} Traduis le texte suivant en français :

\begin{textebox}[Version]
学习汉语有很多好处。第一，汉语很有意思；第二，中国很大，文化很丰富。我不但学习汉语，而且学习中国历史。\\
\emph{Xuéxí Hànyǔ yǒu hěn duō hǎochu. Dìyī, Hànyǔ hěn yǒu yìsi; dì'èr, Zhōngguó hěn dà, wénhuà hěn fēngfù. Wǒ bùdàn xuéxí Hànyǔ, érqiě xuéxí Zhōngguó lìshǐ.}
\end{textebox}
\end{exercice}

\begin{exercice}[Type Bac — Correction et production écrite]
\textbf{A. Correction.} Les phrases suivantes contiennent chacune une erreur. Identifie l'erreur, corrige la phrase et justifie la règle en français.
\begin{enumerate}
\item \zh{不但他学习好，而且体育好。} (même sujet)
\item \zh{他不但会说汉语，说得很好。}
\item \zh{不但我喜欢中国，我的朋友们喜欢中国。}
\item \zh{第两，我们学校有图书馆。}
\end{enumerate}

\textbf{B. Production écrite.} Rédige un paragraphe d'au moins 6 phrases (environ 80 caractères chinois) sur le sujet : « Les avantages de ma ville » (\zh{我的城市的优点}). Consignes obligatoires :
\begin{itemize}
\item utilise au moins une fois la structure \zh{第一……第二……} (tu peux ajouter \zh{第三}) ;
\item utilise au moins une fois la structure \zh{不但……而且……} ;
\item emploie le vocabulaire du module (ville, école, marché, habitants, culture…) ;
\item soigne la ponctuation chinoise (， 。 ；).
\end{itemize}
\end{exercice}

\newpage

\coursec{Corrigés des exercices}

\begin{corrige}[Exercice 1 — QCM : la bonne structure]
\begin{enumerate}
\item \textbf{Réponse A} : \zh{他不但会说英语，而且会说汉语。} \emph{Tā bùdàn huì shuō Yīngyǔ, érqiě huì shuō Hànyǔ.} « Il sait parler non seulement anglais, mais aussi chinois. » Les deux compétences s'ajoutent : c'est une gradation, donc \zh{不但……而且……}. \zh{第一……第二……} énumère des arguments séparés par des virgules ou points-virgules ; \zh{因为……所以……} exprime la cause, hors sujet ici.
\item \textbf{Réponse B} : pour énumérer des arguments (premièrement, deuxièmement…), on utilise \zh{第一……第二……}.
\item \textbf{Réponse A} : \zh{第一} \emph{dìyī} signifie « premièrement » (ou « premier » selon le contexte). Attention : « le premier jour » se dirait \zh{第一天}.
\item \textbf{Réponse B} : quand les deux parties ont le même sujet, le sujet se place \textbf{avant} \zh{不但} : Sujet + \zh{不但} + A + \zh{而且} + B. Exemple : \zh{他不但聪明，而且努力。}
\end{enumerate}
\end{corrige}

\begin{corrige}[Exercice 2 — Vrai ou faux ? Justifie ta réponse]
\begin{enumerate}
\item \textbf{F} (fautive). Quand les deux parties ont le même sujet (\zh{我}), le sujet doit se placer \textbf{avant} \zh{不但}. Correction : \zh{我不但喜欢足球，而且喜欢篮球。} \emph{Wǒ bùdàn xǐhuan zúqiú, érqiě xǐhuan lánqiú.} « J'aime non seulement le football, mais aussi le basketball. »
\item \textbf{V} (correcte). \zh{第二} introduit bien le deuxième point d'une énumération : \zh{第二，我们学习汉语。} « Deuxièmement, nous apprenons le chinois. »
\item \textbf{V} (correcte). Même sujet \zh{他} placé avant \zh{不但} ; la seconde partie reprend le verbe : \zh{他不但会说汉语，而且说得很好。} « Non seulement il sait parler chinois, mais en plus il le parle très bien. »
\item \textbf{F} (fautive). L'énumération exige le préfixe ordinal \zh{第} + \zh{一/二/三} ; \zh{两} ne s'emploie jamais après \zh{第}. Correction : \zh{第二，喀麦隆有很多学校。} « Deuxièmement, le Cameroun a beaucoup d'écoles. »
\end{enumerate}
\end{corrige}

\begin{corrige}[Exercice 3 — Vocabulaire : complète le tableau]
\begin{center}
\begin{tabularx}{\linewidth}{|X|X|X|X|}
\hline
\rowcolor{popblueL} Caractères & Pinyin & Nature & Traduction \\
\hline
\zh{第一} & dìyī & mot ordinal & premièrement ; premier \\
\hline
\zh{不但} & bùdàn & conjonction & non seulement \\
\hline
\zh{而且} & érqiě & conjonction & mais aussi ; et de plus \\
\hline
\zh{也} & yě & adverbe & aussi \\
\hline
\zh{优点} & yōudiǎn & nom & avantage, point fort \\
\hline
\zh{方面} & fāngmiàn & nom & aspect, domaine \\
\hline
\end{tabularx}
\end{center}
\textbf{Points de vigilance} : le ton de \zh{不} devient 2\up{e} ton devant un 4\up{e} ton : \zh{不但} se prononce \emph{búdàn}. \zh{第} porte la clé du bambou \zh{⺮} ; ne pas le confondre avec \zh{弟} \emph{dì} « petit frère ».
\end{corrige}

\begin{corrige}[Exercice 4 — Associe le pinyin aux caractères]
\begin{enumerate}
\item \zh{第二} → \textbf{c. dì'èr} (apostrophe obligatoire pour séparer les syllabes)
\item \zh{而且} → \textbf{a. érqiě}
\item \zh{不但} → \textbf{d. bùdàn}
\item \zh{第三} → \textbf{b. dìsān}
\end{enumerate}
\end{corrige}

\begin{corrige}[Exercice 5 — Traduis en français]
\begin{enumerate}
\item \zh{第一，汉语很有意思；第二，汉语很有用。} → « Premièrement, le chinois est très intéressant ; deuxièmement, le chinois est très utile. »
\item \zh{他不但学习好，而且体育也很好。} → « Non seulement il est bon en études, mais il est aussi très bon en sport. »
\item \zh{不但我喜欢喀麦隆，我的朋友们也喜欢。} → « Non seulement j'aime le Cameroun, mais mes amis l'aiment aussi. » (Sujets différents : \zh{不但} en tête de phrase, \zh{也} dans la seconde partie.)
\end{enumerate}
\end{corrige}

\begin{corrige}[Exercice 6 — Transformation : fusionne les phrases]
\begin{enumerate}
\item \zh{他不但会说法语，而且会说英语。} \emph{Tā bùdàn huì shuō Fǎyǔ, érqiě huì shuō Yīngyǔ.} « Il sait parler non seulement français, mais aussi anglais. » (Même sujet : \zh{他} avant \zh{不但}.)
\item \zh{雅温得不但很大，而且很美。} \emph{Yǎwēndé bùdàn hěn dà, érqiě hěn měi.} « Yaoundé est non seulement très grande, mais aussi très belle. »
\item \zh{不但我喜欢足球，我弟弟也喜欢足球。} \emph{Bùdàn wǒ xǐhuan zúqiú, wǒ dìdi yě xǐhuan zúqiú.} « Non seulement j'aime le football, mais mon petit frère l'aime aussi. » (Sujets différents : \zh{不但} avant le premier sujet, \zh{也} obligatoire dans la seconde partie.)
\item \zh{她不但学习汉语，而且学习中国文化。} \emph{Tā bùdàn xuéxí Hànyǔ, érqiě xuéxí Zhōngguó wénhuà.} « Elle étudie non seulement le chinois, mais aussi la culture chinoise. »
\item \zh{这个菜不但好吃，而且不贵。} \emph{Zhège cài bùdàn hǎochī, érqiě bù guì.} « Ce plat est non seulement bon, mais aussi pas cher. »
\end{enumerate}
\textbf{Méthode} : 1) comparer les sujets des deux phrases ; 2) si le sujet est identique, le placer une seule fois avant \zh{不但} ; 3) si les sujets sont différents, mettre \zh{不但} en tête et ajouter \zh{也} dans la seconde partie.
\end{corrige}

\begin{corrige}[Exercice 7 — Texte à trous]
Texte complété :
我的学校很好。\zh{第一}，老师很好；\zh{第二}，同学很友好。我们学校\zh{不但}有很大的图书馆，\zh{而且}有漂亮的操场。我很爱学习，我的朋友们\zh{也}很爱学习。

\emph{Wǒ de xuéxiào hěn hǎo. Dìyī, lǎoshī hěn hǎo; dì'èr, tóngxué hěn yǒuhǎo. Wǒmen xuéxiào bùdàn yǒu hěn dà de túshūguǎn, érqiě yǒu piàoliang de cāochǎng. Wǒ hěn ài xuéxí, wǒ de péngyoumen yě hěn ài xuéxí.}

Traduction : « Mon école est très bien. Premièrement, les professeurs sont très bien ; deuxièmement, les camarades sont très amicaux. Notre école a non seulement une très grande bibliothèque, mais aussi un beau terrain de sport. J'aime beaucoup étudier, et mes amis aiment aussi beaucoup étudier. »

\textbf{Justifications} : les deux premiers trous introduisent une énumération d'arguments → \zh{第一} / \zh{第二} ; la phrase suivante ajoute une qualité à une autre (même sujet \zh{我们学校}) → \zh{不但……而且……} ; le dernier trou exprime « aussi » avec un nouveau sujet → \zh{也}.
\end{corrige}

\begin{corrige}[Exercice 8 — Remets les mots dans l'ordre]
\begin{enumerate}
\item \zh{我不但学习汉语，而且也学习英语。} \emph{Wǒ bùdàn xuéxí Hànyǔ, érqiě yě xuéxí Yīngyǔ.} « J'étudie non seulement le chinois, mais aussi l'anglais. » (Même sujet \zh{我} avant \zh{不但}.)
\item \zh{第二，杜阿拉是一个大城市。} \emph{Dì'èr, Dù'ālā shì yí ge dà chéngshì.} « Deuxièmement, Douala est une grande ville. »
\item \zh{他不但喜欢足球，而且喜欢篮球。} \emph{Tā bùdàn xǐhuan zúqiú, érqiě xǐhuan lánqiú.} « Il aime non seulement le football, mais aussi le basketball. »
\item \zh{第一，喀麦隆很美；第二，人很友好。} \emph{Dìyī, Kāmàilóng hěn měi; dì'èr, rén hěn yǒuhǎo.} « Premièrement, le Cameroun est très beau ; deuxièmement, les gens sont très amicaux. »
\end{enumerate}
\end{corrige}

\begin{corrige}[Exercice 9 — Type Bac — Compréhension de texte et questions de langue]
\textbf{Barème indicatif (sur 10)} : compréhension 4 pts (1 + 2 + 1) ; langue 6 pts (2 + 2 + 2).

\textbf{Questions de compréhension}
\begin{enumerate}
\item L'auteur habite à Douala (\zh{杜阿拉}), une grande ville du Cameroun.
\item Les trois avantages : premièrement, Douala est très grande et très peuplée ; deuxièmement, elle possède un port et son économie est très développée ; troisièmement, ses habitants sont très amicaux.
\item Selon la dernière phrase, les jeunes (\zh{年轻人}) aiment cette ville, et les personnes âgées (\zh{老年人}) l'aiment aussi.
\end{enumerate}

\textbf{Questions de langue}
\begin{enumerate}
\item Même sujet : \zh{杜阿拉不但经济发达，而且文化也很丰富。} (le sujet \zh{杜阿拉} est placé avant \zh{不但}). Sujets différents : \zh{不但年轻人喜欢这个城市，老年人也喜欢。} (\zh{不但} en tête de phrase, avant le premier sujet \zh{年轻人}).
\item \zh{也} « aussi » apparaît parce que les deux sujets sont différents : la règle impose, dans la seconde partie, un marqueur d'addition (\zh{也}) après le nouveau sujet : \zh{不但 + S1 + …, S2 + 也 + …}. Sans \zh{也}, la phrase serait fautive.
\item \zh{杜阿拉不但很大，而且人很多。} \emph{Dù'ālā bùdàn hěn dà, érqiě rén hěn duō.} « Douala est non seulement très grande, mais aussi très peuplée. » (Même sujet → \zh{杜阿拉} avant \zh{不但}.)
\end{enumerate}
\textbf{Points de vigilance} : ne pas traduire \zh{发达} par « développer » (verbe) mais « développé » ; bien repérer la place du sujet avant de transformer.
\end{corrige}

\begin{corrige}[Exercice 10 — Type Bac — Thème et version]
\textbf{Barème indicatif (sur 10)} : thème 6 pts (1,5 pt par phrase) ; version 4 pts.

\textbf{A. Thème — corrigé}
\begin{enumerate}
\item \zh{第一，汉语很有意思；第二，汉语很有用。} \emph{Dìyī, Hànyǔ hěn yǒu yìsi; dì'èr, Hànyǔ hěn yǒuyòng.}
\item \zh{我哥哥不但喜欢足球，而且喜欢篮球。} \emph{Wǒ gēge bùdàn xǐhuan zúqiú, érqiě xǐhuan lánqiú.} (Même sujet : \zh{我哥哥} avant \zh{不但}.)
\item \zh{不但我喜欢雅温得，我父母也喜欢。} \emph{Bùdàn wǒ xǐhuan Yǎwēndé, wǒ fùmǔ yě xǐhuan.} (Sujets différents : \zh{不但} en tête, \zh{也} obligatoire.)
\item \zh{这个学校不但很大，而且很漂亮。} \emph{Zhège xuéxiào bùdàn hěn dà, érqiě hěn piàoliang.}
\end{enumerate}

\textbf{B. Version — traduction}
« Apprendre le chinois a beaucoup d'avantages. Premièrement, le chinois est très intéressant ; deuxièmement, la Chine est très grande et sa culture est très riche. Je n'étudie pas seulement le chinois, j'étudie aussi l'histoire de la Chine. »

\textbf{Points de vigilance} : \zh{好处} \emph{hǎochu} = « avantage, bienfait » ; \zh{历史} \emph{lìshǐ} = « histoire » ; ne pas oublier le point-virgule chinois ；dans la phrase 1 du thème ; dans la phrase 3, l'erreur classique serait de placer \zh{我} avant \zh{不但}.
\end{corrige}

\begin{corrige}[Exercice 11 — Type Bac — Correction et production écrite]
\textbf{Barème indicatif (sur 10)} : correction 4 pts (1 pt par phrase) ; production écrite 6 pts (structures 2 pts, vocabulaire 1 pt, exactitude grammaticale 2 pts, ponctuation et présentation 1 pt).

\textbf{A. Correction}
\begin{enumerate}
\item Erreur : place du sujet. Correction : \zh{他不但学习好，而且体育好。} « Non seulement il est bon en études, mais aussi en sport. » Règle : même sujet → le sujet se place avant \zh{不但}.
\item Erreur : il manque \zh{而且}. Correction : \zh{他不但会说汉语，而且说得很好。} « Non seulement il sait parler chinois, mais il le parle très bien. » Règle : \zh{不但} appelle toujours son corrélat \zh{而且}.
\item Erreur : il manque \zh{也} dans la seconde partie. Correction : \zh{不但我喜欢中国，我的朋友们也喜欢中国。} « Non seulement j'aime la Chine, mais mes amis l'aiment aussi. » Règle : sujets différents → \zh{也} après le second sujet.
\item Erreur : \zh{两} impossible après \zh{第}. Correction : \zh{第二，我们学校有图书馆。} « Deuxièmement, notre école a une bibliothèque. » Règle : les ordinaux se forment avec \zh{第} + \zh{一/二/三} ; \zh{两} sert uniquement devant un classificateur (\zh{两个人}).
\end{enumerate}

\textbf{B. Production écrite — proposition de rédaction modèle}

我的城市是雅温得。雅温得有很多优点。第一，雅温得的学校很多，老师很好；第二，这里的人很友好；第三，雅温得的市场很大，东西不贵。我的城市不但很美，而且文化很丰富。不但我喜欢雅温得，我的家人也喜欢。我很爱我的城市。

\emph{Wǒ de chéngshì shì Yǎwēndé. Yǎwēndé yǒu hěn duō yōudiǎn. Dìyī, Yǎwēndé de xuéxiào hěn duō, lǎoshī hěn hǎo; dì'èr, zhèlǐ de rén hěn yǒuhǎo; dìsān, Yǎwēndé de shìchǎng hěn dà, dōngxi bù guì. Wǒ de chéngshì bùdàn hěn měi, érqiě wénhuà hěn fēngfù. Bùdàn wǒ xǐhuan Yǎwēndé, wǒ de jiārén yě xǐhuan. Wǒ hěn ài wǒ de chéngshì.}

Traduction : « Ma ville est Yaoundé. Yaoundé a beaucoup d'avantages. Premièrement, il y a beaucoup d'écoles à Yaoundé et les professeurs sont très bien ; deuxièmement, les gens d'ici sont très amicaux ; troisièmement, les marchés de Yaoundé sont très grands et les choses ne sont pas chères. Ma ville est non seulement très belle, mais sa culture est aussi très riche. Non seulement j'aime Yaoundé, mais ma famille l'aime aussi. J'aime beaucoup ma ville. »

\textbf{Points de vigilance} : vérifier la place du sujet dans chaque \zh{不但……而且……} ; employer la ponctuation chinoise pleine chasse (， 。 ；) ; compter environ 80 caractères ; éviter la répétition excessive de \zh{很} en variant les adjectifs (\zh{友好、丰富、漂亮}).
\end{corrige}

\newpage

\coursec{Fiche bilan}

\begin{popfigure}
\begin{tikzpicture}[mindmap, grow cyclic, every node/.style={concept, font=\small},
  root concept/.append style={concept color=popblue, font=\large\bfseries},
  level 1/.append style={level distance=3.4cm, sibling angle=60},
  level 1 concept/.append style={font=\small}]
\node[root concept, align=center]{Leçon 2\\ 第一……第二……\\ 不但……而且……}
  child[concept color=poporange]{ node[align=center]{Énumération\\ 第一、第二、第三……}
    child[concept color=poporangeL]{ node[align=center]{dì-yī, dì-èr\\ premier, deuxième} }
  }
  child[concept color=popgreen]{ node[align=center]{Gradation\\ 不但……而且……}
    child[concept color=popgreenL]{ node[align=center]{bùdàn... érqiě...\\ non seulement... mais aussi...} }
  }
  child[concept color=poppurple]{ node[align=center]{Ordre des mots\\ 不但 + SV，而且 + SV}
  }
  child[concept color=poppink]{ node[align=center]{Sujets identiques ou différents\\ 他不但……，而且……\\ 不但他……，而且我……}
  }
  child[concept color=popgold]{ node[align=center]{Pièges des francophones\\ pas de 很 avant l'adjectif après 而且\\ ne pas traduire mot à mot}
  }
  child[concept color=popblueL]{ node[align=center]{Emplois\\ exposé, lettre,\\ description de famille}
  };
\end{tikzpicture}
\legende{Carte mentale de la leçon : énumération avec 第一……第二…… et gradation avec 不但……而且……}
\end{popfigure}

\begin{bilan}
\textbf{1. Lexique clé à connaître par cœur}

\begin{center}
\begin{tabularx}{\linewidth}{|X|X|X|X|}
\hline
\rowcolor{popblueL} Caractères & Pinyin & Nature & Traduction \\
\hline
第一 & dì-yī & nombre ordinal & premier, premièrement \\
\hline
第二 & dì-èr & nombre ordinal & deuxième, deuxièmement \\
\hline
第三 & dì-sān & nombre ordinal & troisième, troisièmement \\
\hline
不但 & bùdàn & conjonction & non seulement \\
\hline
而且 & érqiě & conjonction & mais aussi, et même \\
\hline
最后 & zuìhòu & adverbe & enfin, pour finir \\
\hline
原因 & yuányīn & nom & raison, cause \\
\hline
优点 & yōudiǎn & nom & avantage, qualité \\
\hline
\end{tabularx}
\end{center}

\textbf{2. Règles de grammaire à connaître par cœur}

\begin{center}
\begin{tabularx}{\linewidth}{|X|X|X|}
\hline
\rowcolor{popblueL} Structure & Exemple (caractères + pinyin) & Traduction \\
\hline
第一……，第二……，（第三……） & 我喜欢汉语，第一，汉语很有意思；第二，中国很大。Wǒ xǐhuan Hànyǔ, dì-yī, Hànyǔ hěn yǒu yìsi; dì-èr, Zhōngguó hěn dà. & J'aime le chinois : premièrement, le chinois est intéressant ; deuxièmement, la Chine est très grande. \\
\hline
Sujet + 不但 + SV1，而且 + SV2 (même sujet) & 他不但会说汉语，而且会写汉字。Tā bùdàn huì shuō Hànyǔ, érqiě huì xiě Hànzì. & Non seulement il sait parler chinois, mais il sait aussi écrire les caractères. \\
\hline
不但 + S1 + SV，而且 + S2 + 也 + SV (sujets différents) & 不但我喜欢足球，而且我弟弟也喜欢。Bùdàn wǒ xǐhuan zúqiú, érqiě wǒ dìdi yě xǐhuan. & Non seulement moi j'aime le football, mais mon petit frère l'aime aussi. \\
\hline
\end{tabularx}
\end{center}

\textbf{3. Méthodes express}
\begin{itemize}
\item \cle{Énumération} : pour présenter des raisons ou des arguments, annoncer l'idée générale, puis développer avec 第一……，第二……，（第三……），最后……
\item \cle{Gradation} : vérifier d'abord si les deux propositions ont le même sujet ; si oui, 不但 se place après le sujet ; si non, 不但 se place avant le premier sujet et on ajoute souvent 也 dans la seconde proposition.
\item \cle{Traduction} : ne jamais traduire « non seulement... mais aussi » mot à mot avec 不 + 是 ; utiliser toujours le couple 不但……而且……
\item \cle{Rédaction} : combiner les deux structures dans un court texte (lettre, exposé) : énumérer avec 第一/第二, renforcer avec 不但/而且.
\end{itemize}
\end{bilan}

\begin{aretenir}[Checklist avant le Bac]
\begin{itemize}
\item[$\square$] Je sais former les nombres ordinaux avec 第 + nombre : 第一、第二、第三……
\item[$\square$] Je sais énumérer des arguments avec 第一……，第二……，最后……
\item[$\square$] Je connais le schéma 不但 + SV1，而且 + SV2 quand le sujet est identique.
\item[$\square$] Je sais placer 不但 avant le sujet quand les deux sujets sont différents.
\item[$\square$] Je pense à ajouter 也 dans la seconde proposition quand les sujets sont différents.
\item[$\square$] Je n'oublie pas la virgule chinoise ，entre les deux propositions.
\item[$\square$] Je n'ajoute pas 很 devant un adjectif quand il n'est pas nécessaire au sens.
\item[$\square$] Je ne traduis jamais « non seulement » par 不是.
\item[$\square$] Je sais écrire les caractères 第、但、而、且 dans le bon ordre des traits.
\item[$\square$] Je sais utiliser les deux structures dans un court texte sur ma famille ou mon école.
\item[$\square$] Je vérifie mes tons : dì-yī (4\textsuperscript{e} + 1\textsuperscript{er} ton), érqiě (2\textsuperscript{e} + 3\textsuperscript{e} ton).
\item[$\square$] Je relis ma production pour vérifier l'ordre des mots : sujet + 不但 / 不但 + sujet.
\end{itemize}
\end{aretenir}

\findecours

\end{document}
